%%C080001 p

本章包含以下案例：


%%C080002 li

\begin{itemize}
\item 使用线程


%%C080003 li

\item 处理线程函数抛出的异常


%%C080004 li

\item 用互斥量与锁同步对共享数据的访问


%%C080005 li

\item 避免使用递归互斥量


%%C080006 li

\item 在线程之间发送通知


%%C080007 li

\item 用promise与future从线程返回值


%%C080008 li

\item 异步执行函数


%%C080009 li

\item 使用原子类型


%%C080010 li

\item 用线程实现并行map与fold


%%C080011 li

\item 用任务实现并行map与fold


\end{itemize}


%%C080012 h2

\section{引言}


%%C080013 p

所有计算机都包含多个处理器，或者至少包含多个核心，而利用这种计算能力是许多类别应用的关键。遗憾的是，许多开发者仍然停留在顺序执行代码的思维定式中，尽管那些互不依赖的操作原本可以并发执行。本章介绍标准库对线程、异步任务及相关组件的支持，并在最后给出一些实用示例。


%%C080014 h2

\section{使用线程}


%%C080015 p

大多数现代处理器（除了那些专用于不需要强大计算能力的应用类型，例如物联网应用的处理器）都有两个、四个或更多核心，使你能够并发执行多个执行线程。应用必须以显式方式编写，才能利用已经存在的多个处理单元；你可以通过同时在多个线程上执行函数来编写这样的应用。C++标准库提供了对线程使用、共享数据同步、线程间通信以及异步任务的支持。本章将探讨与线程和任务相关的最重要主题。


%%C080016 p

线程是一段指令序列，可以由调度器（例如操作系统）独立管理。线程可以是软件线程，它们可以在单个处理单元上运行，通常通过时间片轮转实现；也可以是硬件线程，它们可以在多处理器或多核系统上同时运行，也就是并行运行。多个软件线程也可以并发运行在一个硬件线程上。C++库提供了对软件线程的支持。在本章的第一部分，我们将考察库中内建支持的各种线程对象与机制，例如线程、锁对象、条件变量、异常处理等。在本案例中，你将学习如何创建和管理线程。


%%C080017 h3

\subsection{准备工作}


%%C080018 p

执行线程由\texttt{thread}类表示，该类位于\texttt{<\allowbreak{}thread>}头文件的\texttt{std}命名空间中。其他线程工具位于同一个头文件中，但位于\texttt{std::\allowbreak{}this\_\allowbreak{}thread}命名空间中。


%%C080019 p

在下面的示例中使用了\texttt{print\_\allowbreak{}time()}函数：


%%C080020 code

\begin{codeblockbg}
\begin{Verbatim}[commandchars=\\\{\},breaklines=true,breakanywhere=true,fontsize=\codefont,xleftmargin=17bp]
\PYG{k+kr}{inline}\PYG{+w}{ }\PYG{k+kt}{void}\PYG{+w}{ }\PYG{n+nf}{print\PYGZus{}time}\PYG{p}{(}\PYG{p}{)}
\PYG{p}{\PYGZob{}}
\PYG{+w}{  }\PYG{k}{auto}\PYG{+w}{ }\PYG{n}{now}\PYG{+w}{ }\PYG{o}{=}\PYG{+w}{ }\PYG{n}{std}\PYG{o}{:}\PYG{o}{:}\PYG{n}{chrono}\PYG{o}{:}\PYG{o}{:}\PYG{n}{system\PYGZus{}clock}\PYG{o}{:}\PYG{o}{:}\PYG{n}{now}\PYG{p}{(}\PYG{p}{)}\PYG{p}{;}
\PYG{+w}{  }\PYG{k}{auto}\PYG{+w}{ }\PYG{n}{stime}\PYG{+w}{ }\PYG{o}{=}\PYG{+w}{ }\PYG{n}{std}\PYG{o}{:}\PYG{o}{:}\PYG{n}{chrono}\PYG{o}{:}\PYG{o}{:}\PYG{n}{system\PYGZus{}clock}\PYG{o}{:}\PYG{o}{:}\PYG{n}{to\PYGZus{}time\PYGZus{}t}\PYG{p}{(}\PYG{n}{now}\PYG{p}{)}\PYG{p}{;}
\PYG{+w}{  }\PYG{k}{auto}\PYG{+w}{ }\PYG{n}{ltime}\PYG{+w}{ }\PYG{o}{=}\PYG{+w}{ }\PYG{n}{std}\PYG{o}{:}\PYG{o}{:}\PYG{n}{localtime}\PYG{p}{(}\PYG{o}{\PYGZam{}}\PYG{n}{stime}\PYG{p}{)}\PYG{p}{;}
\PYG{+w}{  }\PYG{n}{std}\PYG{o}{:}\PYG{o}{:}\PYG{n}{cout}\PYG{+w}{ }\PYG{o}{\PYGZlt{}}\PYG{o}{\PYGZlt{}}\PYG{+w}{ }\PYG{n}{std}\PYG{o}{:}\PYG{o}{:}\PYG{n}{put\PYGZus{}time}\PYG{p}{(}\PYG{n}{ltime}\PYG{p}{,}\PYG{+w}{ }\PYG{l+s}{\PYGZdq{}}\PYG{l+s}{\PYGZpc{}c}\PYG{l+s}{\PYGZdq{}}\PYG{p}{)}\PYG{+w}{ }\PYG{o}{\PYGZlt{}}\PYG{o}{\PYGZlt{}}\PYG{+w}{ }\PYG{n}{std}\PYG{o}{:}\PYG{o}{:}\PYG{n}{endl}\PYG{p}{;}
\PYG{p}{\PYGZcb{}}
\end{Verbatim}
\end{codeblockbg}


%%C080021 h3

\subsection{具体做法}


%%C080022 p

使用以下方案来管理线程：


%%C080023 li

\begin{itemize}
\item 要创建一个\texttt{std::\allowbreak{}thread}对象而不启动新线程的执行，可以使用它的默认构造函数：


%%C080024 code

\begin{codeblockbg}
\begin{Verbatim}[commandchars=\\\{\},breaklines=true,breakanywhere=true,fontsize=\codefont,xleftmargin=17bp]
\PYG{n}{std}\PYG{o}{:}\PYG{o}{:}\PYG{k+kr}{thread}\PYG{+w}{ }\PYG{n}{t}\PYG{p}{;}
\end{Verbatim}
\end{codeblockbg}


%%C080025 li

\item 通过构造一个\texttt{std::\allowbreak{}thread}对象并把函数作为参数传入，即可在另一个线程上开始执行该函数：


%%C080026 code

\begin{codeblockbg}
\begin{Verbatim}[commandchars=\\\{\},breaklines=true,breakanywhere=true,fontsize=\codefont,xleftmargin=17bp]
\PYG{k+kt}{void}\PYG{+w}{ }\PYG{n+nf}{func1}\PYG{p}{(}\PYG{p}{)}
\PYG{p}{\PYGZob{}}
\PYG{+w}{  }\PYG{n}{std}\PYG{o}{:}\PYG{o}{:}\PYG{n}{cout}\PYG{+w}{ }\PYG{o}{\PYGZlt{}}\PYG{o}{\PYGZlt{}}\PYG{+w}{ }\PYG{l+s}{\PYGZdq{}}\PYG{l+s}{thread func without params}\PYG{l+s}{\PYGZdq{}}\PYG{+w}{ }\PYG{o}{\PYGZlt{}}\PYG{o}{\PYGZlt{}}\PYG{+w}{ }\PYG{n}{std}\PYG{o}{:}\PYG{o}{:}\PYG{n}{endl}\PYG{p}{;}
\PYG{p}{\PYGZcb{}}
\PYG{n}{std}\PYG{o}{:}\PYG{o}{:}\PYG{k+kr}{thread}\PYG{+w}{ }\PYG{n+nf}{t}\PYG{p}{(}\PYG{n}{func1}\PYG{p}{)}\PYG{p}{;}
\PYG{n}{std}\PYG{o}{:}\PYG{o}{:}\PYG{k+kr}{thread}\PYG{+w}{ }\PYG{n+nf}{t}\PYG{p}{(}\PYG{p}{[}\PYG{p}{]}\PYG{p}{(}\PYG{p}{)}\PYG{+w}{ }\PYG{p}{\PYGZob{}}
\PYG{+w}{  }\PYG{n}{std}\PYG{o}{:}\PYG{o}{:}\PYG{n}{cout}\PYG{+w}{ }\PYG{o}{\PYGZlt{}}\PYG{o}{\PYGZlt{}}\PYG{+w}{ }\PYG{l+s}{\PYGZdq{}}\PYG{l+s}{thread func without params}\PYG{l+s}{\PYGZdq{}}
\PYG{+w}{            }\PYG{o}{\PYGZlt{}}\PYG{o}{\PYGZlt{}}\PYG{+w}{ }\PYG{n}{std}\PYG{o}{:}\PYG{o}{:}\PYG{n}{endl}\PYG{p}{;}\PYG{+w}{ }\PYG{p}{\PYGZcb{}}\PYG{p}{)}\PYG{p}{;}
\end{Verbatim}
\end{codeblockbg}


%%C080027 li

\item 通过构造一个\texttt{std::\allowbreak{}thread}对象，把函数作为参数传给构造函数、随后传入它的参数，即可在另一个线程上开始执行带参数的函数：


%%C080028 code

\begin{codeblockbg}
\begin{Verbatim}[commandchars=\\\{\},breaklines=true,breakanywhere=true,fontsize=\codefont,xleftmargin=17bp]
\PYG{k+kt}{void}\PYG{+w}{ }\PYG{n+nf}{func2}\PYG{p}{(}\PYG{k+kt}{int}\PYG{+w}{ }\PYG{k}{const}\PYG{+w}{ }\PYG{n}{i}\PYG{p}{,}\PYG{+w}{ }\PYG{k+kt}{double}\PYG{+w}{ }\PYG{k}{const}\PYG{+w}{ }\PYG{n}{d}\PYG{p}{,}\PYG{+w}{ }\PYG{n}{std}\PYG{o}{:}\PYG{o}{:}\PYG{n}{string}\PYG{+w}{ }\PYG{k}{const}\PYG{+w}{ }\PYG{n}{s}\PYG{p}{)}
\PYG{p}{\PYGZob{}}
\PYG{+w}{  }\PYG{n}{std}\PYG{o}{:}\PYG{o}{:}\PYG{n}{cout}\PYG{+w}{ }\PYG{o}{\PYGZlt{}}\PYG{o}{\PYGZlt{}}\PYG{+w}{ }\PYG{n}{i}\PYG{+w}{ }\PYG{o}{\PYGZlt{}}\PYG{o}{\PYGZlt{}}\PYG{+w}{ }\PYG{l+s}{\PYGZdq{}}\PYG{l+s}{, }\PYG{l+s}{\PYGZdq{}}\PYG{+w}{ }\PYG{o}{\PYGZlt{}}\PYG{o}{\PYGZlt{}}\PYG{+w}{ }\PYG{n}{d}\PYG{+w}{ }\PYG{o}{\PYGZlt{}}\PYG{o}{\PYGZlt{}}\PYG{+w}{ }\PYG{l+s}{\PYGZdq{}}\PYG{l+s}{, }\PYG{l+s}{\PYGZdq{}}\PYG{+w}{ }\PYG{o}{\PYGZlt{}}\PYG{o}{\PYGZlt{}}\PYG{+w}{ }\PYG{n}{s}\PYG{+w}{ }\PYG{o}{\PYGZlt{}}\PYG{o}{\PYGZlt{}}\PYG{+w}{ }\PYG{n}{std}\PYG{o}{:}\PYG{o}{:}\PYG{n}{endl}\PYG{p}{;}
\PYG{p}{\PYGZcb{}}
\PYG{n}{std}\PYG{o}{:}\PYG{o}{:}\PYG{k+kr}{thread}\PYG{+w}{ }\PYG{n}{t}\PYG{p}{(}\PYG{n}{func2}\PYG{p}{,}\PYG{+w}{ }\PYG{l+m+mi}{42}\PYG{p}{,}\PYG{+w}{ }\PYG{l+m+mf}{42.0}\PYG{p}{,}\PYG{+w}{ }\PYG{l+s}{\PYGZdq{}}\PYG{l+s}{42}\PYG{l+s}{\PYGZdq{}}\PYG{p}{)}\PYG{p}{;}
\end{Verbatim}
\end{codeblockbg}


%%C080029 li

\item 要等待线程结束执行，可以对\texttt{thread}对象调用\texttt{join()}方法：


%%C080030 code

\begin{codeblockbg}
\begin{Verbatim}[commandchars=\\\{\},breaklines=true,breakanywhere=true,fontsize=\codefont,xleftmargin=17bp]
\PYG{n}{t}\PYG{p}{.}\PYG{n}{join}\PYG{p}{(}\PYG{p}{)}\PYG{p}{;}
\end{Verbatim}
\end{codeblockbg}


%%C080031 li

\item 要让线程独立于当前线程对象继续执行，可以使用\texttt{detach()}方法：


%%C080032 code

\begin{codeblockbg}
\begin{Verbatim}[commandchars=\\\{\},breaklines=true,breakanywhere=true,fontsize=\codefont,xleftmargin=17bp]
\PYG{n}{t}\PYG{p}{.}\PYG{n}{detach}\PYG{p}{(}\PYG{p}{)}\PYG{p}{;}
\end{Verbatim}
\end{codeblockbg}


%%C080033 li

\item 要以引用方式把参数传给函数，需要把它们包装在\texttt{std::\allowbreak{}ref}或\texttt{std::\allowbreak{}cref}中（如果引用是常量引用）：


%%C080034 code

\begin{codeblockbg}
\begin{Verbatim}[commandchars=\\\{\},breaklines=true,breakanywhere=true,fontsize=\codefont,xleftmargin=17bp]
\PYG{k+kt}{void}\PYG{+w}{ }\PYG{n+nf}{func3}\PYG{p}{(}\PYG{k+kt}{int}\PYG{+w}{ }\PYG{o}{\PYGZam{}}\PYG{+w}{ }\PYG{n}{i}\PYG{p}{)}
\PYG{p}{\PYGZob{}}
\PYG{+w}{  }\PYG{n}{i}\PYG{+w}{ }\PYG{o}{*}\PYG{o}{=}\PYG{+w}{ }\PYG{l+m+mi}{2}\PYG{p}{;}
\PYG{p}{\PYGZcb{}}
\PYG{k+kt}{int}\PYG{+w}{ }\PYG{n}{n}\PYG{+w}{ }\PYG{o}{=}\PYG{+w}{ }\PYG{l+m+mi}{42}\PYG{p}{;}
\PYG{n}{std}\PYG{o}{:}\PYG{o}{:}\PYG{k+kr}{thread}\PYG{+w}{ }\PYG{n+nf}{t}\PYG{p}{(}\PYG{n}{func3}\PYG{p}{,}\PYG{+w}{ }\PYG{n}{std}\PYG{o}{:}\PYG{o}{:}\PYG{n}{ref}\PYG{p}{(}\PYG{n}{n}\PYG{p}{)}\PYG{p}{)}\PYG{p}{;}
\PYG{n}{t}\PYG{p}{.}\PYG{n}{join}\PYG{p}{(}\PYG{p}{)}\PYG{p}{;}
\PYG{n}{std}\PYG{o}{:}\PYG{o}{:}\PYG{n}{cout}\PYG{+w}{ }\PYG{o}{\PYGZlt{}}\PYG{o}{\PYGZlt{}}\PYG{+w}{ }\PYG{n}{n}\PYG{+w}{ }\PYG{o}{\PYGZlt{}}\PYG{o}{\PYGZlt{}}\PYG{+w}{ }\PYG{n}{std}\PYG{o}{:}\PYG{o}{:}\PYG{n}{endl}\PYG{p}{;}\PYG{+w}{ }\PYG{c+c1}{// 84}
\end{Verbatim}
\end{codeblockbg}


%%C080035 li

\item 要让线程暂停执行一段指定的时长，可以使用\myttbr{std::\allowbreak{}this\_\allowbreak{}thread::\allowbreak{}sleep\_\allowbreak{}for()}函数：


%%C080036 code

\begin{codeblockbg}
\begin{Verbatim}[commandchars=\\\{\},breaklines=true,breakanywhere=true,fontsize=\codefont,xleftmargin=17bp]
\PYG{k+kt}{void}\PYG{+w}{ }\PYG{n+nf}{func4}\PYG{p}{(}\PYG{p}{)}
\PYG{p}{\PYGZob{}}
\PYG{+w}{  }\PYG{k}{using}\PYG{+w}{ }\PYG{k}{namespace}\PYG{+w}{ }\PYG{n+nn}{std}\PYG{o}{:}\PYG{o}{:}\PYG{n+nn}{chrono}\PYG{p}{;}
\PYG{+w}{  }\PYG{n}{print\PYGZus{}time}\PYG{p}{(}\PYG{p}{)}\PYG{p}{;}
\PYG{+w}{  }\PYG{n}{std}\PYG{o}{:}\PYG{o}{:}\PYG{n}{this\PYGZus{}thread}\PYG{o}{:}\PYG{o}{:}\PYG{n}{sleep\PYGZus{}for}\PYG{p}{(}\PYG{l+m+mi}{2}\PYG{n}{s}\PYG{p}{)}\PYG{p}{;}
\PYG{+w}{  }\PYG{n}{print\PYGZus{}time}\PYG{p}{(}\PYG{p}{)}\PYG{p}{;}
\PYG{p}{\PYGZcb{}}
\PYG{n}{std}\PYG{o}{:}\PYG{o}{:}\PYG{k+kr}{thread}\PYG{+w}{ }\PYG{n}{t}\PYG{p}{(}\PYG{n}{func4}\PYG{p}{)}\PYG{p}{;}
\PYG{n}{t}\PYG{p}{.}\PYG{n}{join}\PYG{p}{(}\PYG{p}{)}\PYG{p}{;}
\end{Verbatim}
\end{codeblockbg}


%%C080037 li

\item 要让线程暂停执行直到某个指定的时刻，可以使用\myttbr{std::\allowbreak{}this\_\allowbreak{}thread::\allowbreak{}sleep\_\allowbreak{}until()}函数：


%%C080038 code

\begin{codeblockbg}
\begin{Verbatim}[commandchars=\\\{\},breaklines=true,breakanywhere=true,fontsize=\codefont,xleftmargin=17bp]
\PYG{k+kt}{void}\PYG{+w}{ }\PYG{n+nf}{func5}\PYG{p}{(}\PYG{p}{)}
\PYG{p}{\PYGZob{}}
\PYG{+w}{  }\PYG{k}{using}\PYG{+w}{ }\PYG{k}{namespace}\PYG{+w}{ }\PYG{n+nn}{std}\PYG{o}{:}\PYG{o}{:}\PYG{n+nn}{chrono}\PYG{p}{;}
\PYG{+w}{  }\PYG{n}{print\PYGZus{}time}\PYG{p}{(}\PYG{p}{)}\PYG{p}{;}
\PYG{+w}{  }\PYG{n}{std}\PYG{o}{:}\PYG{o}{:}\PYG{n}{this\PYGZus{}thread}\PYG{o}{:}\PYG{o}{:}\PYG{n}{sleep\PYGZus{}until}\PYG{p}{(}
\PYG{+w}{  }\PYG{n}{std}\PYG{o}{:}\PYG{o}{:}\PYG{n}{chrono}\PYG{o}{:}\PYG{o}{:}\PYG{n}{system\PYGZus{}clock}\PYG{o}{:}\PYG{o}{:}\PYG{n}{now}\PYG{p}{(}\PYG{p}{)}\PYG{+w}{ }\PYG{o}{+}\PYG{+w}{ }\PYG{l+m+mi}{2}\PYG{n}{s}\PYG{p}{)}\PYG{p}{;}
\PYG{+w}{  }\PYG{n}{print\PYGZus{}time}\PYG{p}{(}\PYG{p}{)}\PYG{p}{;}
\PYG{p}{\PYGZcb{}}
\PYG{n}{std}\PYG{o}{:}\PYG{o}{:}\PYG{k+kr}{thread}\PYG{+w}{ }\PYG{n}{t}\PYG{p}{(}\PYG{n}{func5}\PYG{p}{)}\PYG{p}{;}
\PYG{n}{t}\PYG{p}{.}\PYG{n}{join}\PYG{p}{(}\PYG{p}{)}\PYG{p}{;}
\end{Verbatim}
\end{codeblockbg}


%%C080039 li

\item 要挂起当前线程的执行，并给另一个线程提供执行机会，可以使用\myttbr{std::\allowbreak{}this\_\allowbreak{}thread::\allowbreak{}yield()}：


%%C080040 code

\begin{codeblockbg}
\begin{Verbatim}[commandchars=\\\{\},breaklines=true,breakanywhere=true,fontsize=\codefont,xleftmargin=17bp]
\PYG{k+kt}{void}\PYG{+w}{ }\PYG{n+nf}{func6}\PYG{p}{(}\PYG{n}{std}\PYG{o}{:}\PYG{o}{:}\PYG{n}{chrono}\PYG{o}{:}\PYG{o}{:}\PYG{n}{seconds}\PYG{+w}{ }\PYG{n}{timeout}\PYG{p}{)}
\PYG{p}{\PYGZob{}}
\PYG{+w}{  }\PYG{k}{auto}\PYG{+w}{ }\PYG{n}{now}\PYG{+w}{ }\PYG{o}{=}\PYG{+w}{ }\PYG{n}{std}\PYG{o}{:}\PYG{o}{:}\PYG{n}{chrono}\PYG{o}{:}\PYG{o}{:}\PYG{n}{system\PYGZus{}clock}\PYG{o}{:}\PYG{o}{:}\PYG{n}{now}\PYG{p}{(}\PYG{p}{)}\PYG{p}{;}
\PYG{+w}{  }\PYG{k}{auto}\PYG{+w}{ }\PYG{n}{then}\PYG{+w}{ }\PYG{o}{=}\PYG{+w}{ }\PYG{n}{now}\PYG{+w}{ }\PYG{o}{+}\PYG{+w}{ }\PYG{n}{timeout}\PYG{p}{;}
\PYG{+w}{  }\PYG{k}{do}
\PYG{+w}{  }\PYG{p}{\PYGZob{}}
\PYG{+w}{    }\PYG{n}{std}\PYG{o}{:}\PYG{o}{:}\PYG{n}{this\PYGZus{}thread}\PYG{o}{:}\PYG{o}{:}\PYG{n}{yield}\PYG{p}{(}\PYG{p}{)}\PYG{p}{;}
\PYG{+w}{  }\PYG{p}{\PYGZcb{}}\PYG{+w}{ }\PYG{k}{while}\PYG{+w}{ }\PYG{p}{(}\PYG{n}{std}\PYG{o}{:}\PYG{o}{:}\PYG{n}{chrono}\PYG{o}{:}\PYG{o}{:}\PYG{n}{system\PYGZus{}clock}\PYG{o}{:}\PYG{o}{:}\PYG{n}{now}\PYG{p}{(}\PYG{p}{)}\PYG{+w}{ }\PYG{o}{\PYGZlt{}}\PYG{+w}{ }\PYG{n}{then}\PYG{p}{)}\PYG{p}{;}
\PYG{p}{\PYGZcb{}}
\PYG{n}{std}\PYG{o}{:}\PYG{o}{:}\PYG{k+kr}{thread}\PYG{+w}{ }\PYG{n+nf}{t}\PYG{p}{(}\PYG{n}{func6}\PYG{p}{,}\PYG{+w}{ }\PYG{n}{std}\PYG{o}{:}\PYG{o}{:}\PYG{n}{chrono}\PYG{o}{:}\PYG{o}{:}\PYG{n}{seconds}\PYG{p}{(}\PYG{l+m+mi}{2}\PYG{p}{)}\PYG{p}{)}\PYG{p}{;}
\PYG{n}{t}\PYG{p}{.}\PYG{n}{join}\PYG{p}{(}\PYG{p}{)}\PYG{p}{;}
\PYG{n}{print\PYGZus{}time}\PYG{p}{(}\PYG{p}{)}\PYG{p}{;}
\end{Verbatim}
\end{codeblockbg}


\end{itemize}


%%C080041 h3

\subsection{工作原理}


%%C080042 p

表示单个执行线程的\texttt{std::\allowbreak{}thread}类有几个构造函数：


%%C080043 li

\begin{itemize}
\item 默认构造函数，它只创建线程对象，但不启动新线程的执行。


%%C080044 li

\item 移动构造函数，它创建一个新的线程对象，用来表示原先由被构造对象所代表的那个执行线程。新对象构造完成后，另一个对象就不再与该执行线程关联。


%%C080045 li

\item 参数个数可变的构造函数：第一个参数是表示顶层线程函数的函数，其余参数是要传给线程函数的实参。实参需要按值传给线程函数。如果线程函数的形参是引用或常量引用，则必须把它们包装在\texttt{std::\allowbreak{}ref}或\texttt{std::\allowbreak{}cref}对象中。


\end{itemize}


%%C080046 p

在这种情况下，线程函数不能返回值。函数实际上具有\texttt{void}之外的返回类型并不违法，但它会忽略函数直接返回的任何值。如果确实需要返回值，可以通过共享变量或函数参数来实现。在后面的案例中，我们将看到线程函数如何借助\textit{promise}把值返回给另一个线程。


%%C080047 p

如果函数因异常而终止，那么在启动线程的上下文中无法用\texttt{try...\allowbreak{}catch}语句捕获该异常，程序会因调用\texttt{std::\allowbreak{}terminate()}而异常终止。所有异常都必须在执行线程内部捕获，但可以通过\myttbr{std::\allowbreak{}exception\_\allowbreak{}ptr}对象在线程之间传送。我们将在下一个案例中讨论这个主题。


%%C080048 p

线程开始执行后，它既可以连接（joinable），也可以分离（detachable）。连接一个线程意味着阻塞当前线程的执行，直到被连接的线程结束执行。分离一个线程则意味着把线程对象与它所代表的执行线程解耦，从而让当前线程和已分离的线程可以同时执行。连接线程用\texttt{join()}完成，分离线程用\texttt{detach()}完成。一旦调用这两个方法中的任何一个，该线程就被称为不可连接的（non-joinable），线程对象也就可以安全地销毁了。线程被分离后，它可能需要访问的共享数据必须在其整个执行期间都保持有效。\texttt{joinable()}方法用于指示线程是否可以连接。


%%C080049 p

每个线程都有一个标识符，可以这样获取——对于当前线程，调用\myttbr{std::\allowbreak{}this\_\allowbreak{}thread::\allowbreak{}get\_\allowbreak{}id()}函数；对于由\texttt{thread}对象表示的其他执行线程，调用它的\texttt{get\_\allowbreak{}id()}方法。


%%C080050 p

\texttt{std::\allowbreak{}this\_\allowbreak{}thread}命名空间中还提供了几个额外的工具函数：


%%C080051 li

\begin{itemize}
\item \texttt{yield()}方法提示调度器激活另一个线程。这在实现忙等待（busy-waiting）例程时很有用，正如上一节的最后一个示例所示。


%%C080052 li

\item \texttt{sleep\_\allowbreak{}for()}方法至少阻塞当前线程执行指定的时间段（由于调度的原因，线程实际休眠的时间可能长于请求的时间）。


%%C080053 li

\item \texttt{sleep\_\allowbreak{}until()}方法至少阻塞当前线程执行到指定的时间点（由于调度的原因，实际休眠的时长可能长于请求的时长）。


\end{itemize}


%%C080054 h3

\subsection{延伸阅读}


%%C080055 li

\begin{itemize}
\item \textit{处理线程函数抛出的异常}


%%C080056 li

\item \textit{用互斥量与锁同步对共享数据的访问}


%%C080057 li

\item \textit{避免使用递归互斥量}


%%C080058 li

\item \textit{在线程之间发送通知}


%%C080059 li

\item \textit{用promise与future从线程返回值}


\end{itemize}


%%C080060 h2

\section{处理线程函数抛出的异常}


%%C080061 p

在上一案例中，我们介绍了线程支持库，并了解了如何用线程做一些基本操作。在那个案例中，我们简要讨论了线程函数中的异常处理，并提到异常不能离开顶层线程函数，因为异常会导致程序因调用\texttt{std::\allowbreak{}terminate()}而异常终止。另一方面，异常可以封装在\myttbr{std::\allowbreak{}exception\_\allowbreak{}ptr}中在线程之间传送；本案例将展示具体做法。


%%C080062 h3

\subsection{准备工作}


%%C080063 p

你现在已经熟悉了上一案例\textit{使用线程}中讨论的线程操作。\texttt{exception\_\allowbreak{}ptr}类位于\texttt{<\allowbreak{}exception>}头文件的\texttt{std}命名空间中；\texttt{mutex}（我们将在下一个案例中更详细地讨论）同样位于该命名空间，但位于\texttt{<\allowbreak{}mutex>}头文件中。


%%C080064 h3

\subsection{具体做法}


%%C080065 p

要从主线程或发起连接的线程中正确处理工作线程抛出的异常，请按以下步骤操作（假设多个线程可能抛出多个异常）：


%%C080066 p

使用一个全局容器来存放\myttbr{std::\allowbreak{}exception\_\allowbreak{}ptr}实例：


%%C080067 code

\begin{codeblockbg}
\begin{Verbatim}[commandchars=\\\{\},breaklines=true,breakanywhere=true,fontsize=\codefont,xleftmargin=17bp]
\PYG{n}{std}\PYG{o}{:}\PYG{o}{:}\PYG{n}{vector}\PYG{o}{\PYGZlt{}}\PYG{n}{std}\PYG{o}{:}\PYG{o}{:}\PYG{n}{exception\PYGZus{}ptr}\PYG{o}{\PYGZgt{}}\PYG{+w}{ }\PYG{n}{g\PYGZus{}exceptions}\PYG{p}{;}
\end{Verbatim}
\end{codeblockbg}


%%C080068 p

使用一个全局\texttt{mutex}来同步对共享容器的访问：


%%C080069 code

\begin{codeblockbg}
\begin{Verbatim}[commandchars=\\\{\},breaklines=true,breakanywhere=true,fontsize=\codefont,xleftmargin=17bp]
\PYG{n}{std}\PYG{o}{:}\PYG{o}{:}\PYG{n}{mutex}\PYG{+w}{ }\PYG{n}{g\PYGZus{}mutex}\PYG{p}{;}
\end{Verbatim}
\end{codeblockbg}


%%C080070 p

对在顶层


%%C080071 p

线程函数中执行的代码使用\texttt{try...\allowbreak{}catch}块。用\myttbr{std::\allowbreak{}current\_\allowbreak{}exception()}捕获当前异常，并把它的副本或其引用包装进\myttbr{std::\allowbreak{}exception\_\allowbreak{}ptr}指针，再添加到用于存放异常的共享容器中：


%%C080072 code

\begin{codeblockbg}
\begin{Verbatim}[commandchars=\\\{\},breaklines=true,breakanywhere=true,fontsize=\codefont,xleftmargin=17bp]
\PYG{k+kt}{void}\PYG{+w}{ }\PYG{n+nf}{func1}\PYG{p}{(}\PYG{p}{)}
\PYG{p}{\PYGZob{}}
\PYG{+w}{  }\PYG{k}{throw}\PYG{+w}{ }\PYG{n}{std}\PYG{o}{:}\PYG{o}{:}\PYG{n}{exception}\PYG{p}{(}\PYG{l+s}{\PYGZdq{}}\PYG{l+s}{exception 1}\PYG{l+s}{\PYGZdq{}}\PYG{p}{)}\PYG{p}{;}
\PYG{p}{\PYGZcb{}}
\PYG{k+kt}{void}\PYG{+w}{ }\PYG{n+nf}{func2}\PYG{p}{(}\PYG{p}{)}
\PYG{p}{\PYGZob{}}
\PYG{+w}{  }\PYG{k}{throw}\PYG{+w}{ }\PYG{n}{std}\PYG{o}{:}\PYG{o}{:}\PYG{n}{exception}\PYG{p}{(}\PYG{l+s}{\PYGZdq{}}\PYG{l+s}{exception 2}\PYG{l+s}{\PYGZdq{}}\PYG{p}{)}\PYG{p}{;}
\PYG{p}{\PYGZcb{}}
\PYG{k+kt}{void}\PYG{+w}{ }\PYG{n+nf}{thread\PYGZus{}func1}\PYG{p}{(}\PYG{p}{)}
\PYG{p}{\PYGZob{}}
\PYG{+w}{  }\PYG{k}{try}
\PYG{+w}{  }\PYG{p}{\PYGZob{}}
\PYG{+w}{    }\PYG{n}{func1}\PYG{p}{(}\PYG{p}{)}\PYG{p}{;}
\PYG{+w}{  }\PYG{p}{\PYGZcb{}}
\PYG{+w}{  }\PYG{k}{catch}\PYG{+w}{ }\PYG{p}{(}\PYG{p}{.}\PYG{p}{.}\PYG{p}{.}\PYG{p}{)}
\PYG{+w}{  }\PYG{p}{\PYGZob{}}
\PYG{+w}{    }\PYG{n}{std}\PYG{o}{:}\PYG{o}{:}\PYG{n}{lock\PYGZus{}guard}\PYG{o}{\PYGZlt{}}\PYG{n}{std}\PYG{o}{:}\PYG{o}{:}\PYG{n}{mutex}\PYG{o}{\PYGZgt{}}\PYG{+w}{ }\PYG{n}{lock}\PYG{p}{(}\PYG{n}{g\PYGZus{}mutex}\PYG{p}{)}\PYG{p}{;}
\PYG{+w}{    }\PYG{n}{g\PYGZus{}exceptions}\PYG{p}{.}\PYG{n}{push\PYGZus{}back}\PYG{p}{(}\PYG{n}{std}\PYG{o}{:}\PYG{o}{:}\PYG{n}{current\PYGZus{}exception}\PYG{p}{(}\PYG{p}{)}\PYG{p}{)}\PYG{p}{;}
\PYG{+w}{  }\PYG{p}{\PYGZcb{}}
\PYG{p}{\PYGZcb{}}
\PYG{k+kt}{void}\PYG{+w}{ }\PYG{n+nf}{thread\PYGZus{}func2}\PYG{p}{(}\PYG{p}{)}
\PYG{p}{\PYGZob{}}
\PYG{+w}{  }\PYG{k}{try}
\PYG{+w}{  }\PYG{p}{\PYGZob{}}
\PYG{+w}{    }\PYG{n}{func2}\PYG{p}{(}\PYG{p}{)}\PYG{p}{;}
\PYG{+w}{  }\PYG{p}{\PYGZcb{}}
\PYG{+w}{  }\PYG{k}{catch}\PYG{+w}{ }\PYG{p}{(}\PYG{p}{.}\PYG{p}{.}\PYG{p}{.}\PYG{p}{)}
\PYG{+w}{  }\PYG{p}{\PYGZob{}}
\PYG{+w}{    }\PYG{n}{std}\PYG{o}{:}\PYG{o}{:}\PYG{n}{lock\PYGZus{}guard}\PYG{o}{\PYGZlt{}}\PYG{n}{std}\PYG{o}{:}\PYG{o}{:}\PYG{n}{mutex}\PYG{o}{\PYGZgt{}}\PYG{+w}{ }\PYG{n}{lock}\PYG{p}{(}\PYG{n}{g\PYGZus{}mutex}\PYG{p}{)}\PYG{p}{;}
\PYG{+w}{    }\PYG{n}{g\PYGZus{}exceptions}\PYG{p}{.}\PYG{n}{push\PYGZus{}back}\PYG{p}{(}\PYG{n}{std}\PYG{o}{:}\PYG{o}{:}\PYG{n}{current\PYGZus{}exception}\PYG{p}{(}\PYG{p}{)}\PYG{p}{)}\PYG{p}{;}
\PYG{+w}{  }\PYG{p}{\PYGZcb{}}
\PYG{p}{\PYGZcb{}}
\end{Verbatim}
\end{codeblockbg}


%%C080073 p

在主线程中，启动这些线程之前先清空容器：


%%C080074 code

\begin{codeblockbg}
\begin{Verbatim}[commandchars=\\\{\},breaklines=true,breakanywhere=true,fontsize=\codefont,xleftmargin=17bp]
\PYG{n}{g\PYGZus{}exceptions}\PYG{p}{.}\PYG{n}{clear}\PYG{p}{(}\PYG{p}{)}\PYG{p}{;}
\end{Verbatim}
\end{codeblockbg}


%%C080075 p

在主线程中，等所有线程执行完毕后，检查


%%C080076 p

捕获到的异常，并对每一个妥善处理：


%%C080077 code

\begin{codeblockbg}
\begin{Verbatim}[commandchars=\\\{\},breaklines=true,breakanywhere=true,fontsize=\codefont,xleftmargin=17bp]
\PYG{n}{std}\PYG{o}{:}\PYG{o}{:}\PYG{k+kr}{thread}\PYG{+w}{ }\PYG{n+nf}{t1}\PYG{p}{(}\PYG{n}{thread\PYGZus{}func1}\PYG{p}{)}\PYG{p}{;}
\PYG{n}{std}\PYG{o}{:}\PYG{o}{:}\PYG{k+kr}{thread}\PYG{+w}{ }\PYG{n+nf}{t2}\PYG{p}{(}\PYG{n}{thread\PYGZus{}func2}\PYG{p}{)}\PYG{p}{;}
\PYG{n}{t1}\PYG{p}{.}\PYG{n}{join}\PYG{p}{(}\PYG{p}{)}\PYG{p}{;}
\PYG{n}{t2}\PYG{p}{.}\PYG{n}{join}\PYG{p}{(}\PYG{p}{)}\PYG{p}{;}
\PYG{k}{for}\PYG{+w}{ }\PYG{p}{(}\PYG{k}{auto}\PYG{+w}{ }\PYG{k}{const}\PYG{+w}{ }\PYG{o}{\PYGZam{}}\PYG{+w}{ }\PYG{n}{e}\PYG{+w}{ }\PYG{o}{:}\PYG{+w}{ }\PYG{n}{g\PYGZus{}exceptions}\PYG{p}{)}
\PYG{p}{\PYGZob{}}
\PYG{+w}{  }\PYG{k}{try}
\PYG{+w}{  }\PYG{p}{\PYGZob{}}
\PYG{+w}{    }\PYG{k}{if}\PYG{p}{(}\PYG{n}{e}\PYG{+w}{ }\PYG{o}{!}\PYG{o}{=}\PYG{+w}{ }\PYG{n+nb}{nullptr}\PYG{p}{)}
\PYG{+w}{      }\PYG{n}{std}\PYG{o}{:}\PYG{o}{:}\PYG{n}{rethrow\PYGZus{}exception}\PYG{p}{(}\PYG{n}{e}\PYG{p}{)}\PYG{p}{;}
\PYG{+w}{  }\PYG{p}{\PYGZcb{}}
\PYG{+w}{  }\PYG{k}{catch}\PYG{p}{(}\PYG{n}{std}\PYG{o}{:}\PYG{o}{:}\PYG{n}{exception}\PYG{+w}{ }\PYG{k}{const}\PYG{+w}{ }\PYG{o}{\PYGZam{}}\PYG{+w}{ }\PYG{n}{ex}\PYG{p}{)}
\PYG{+w}{  }\PYG{p}{\PYGZob{}}
\PYG{+w}{    }\PYG{n}{std}\PYG{o}{:}\PYG{o}{:}\PYG{n}{cout}\PYG{+w}{ }\PYG{o}{\PYGZlt{}}\PYG{o}{\PYGZlt{}}\PYG{+w}{ }\PYG{n}{ex}\PYG{p}{.}\PYG{n}{what}\PYG{p}{(}\PYG{p}{)}\PYG{+w}{ }\PYG{o}{\PYGZlt{}}\PYG{o}{\PYGZlt{}}\PYG{+w}{ }\PYG{n}{std}\PYG{o}{:}\PYG{o}{:}\PYG{n}{endl}\PYG{p}{;}
\PYG{+w}{  }\PYG{p}{\PYGZcb{}}
\PYG{p}{\PYGZcb{}}
\end{Verbatim}
\end{codeblockbg}


%%C080078 h3

\subsection{工作原理}


%%C080079 p

对于上一节的示例，我们假设多个线程都可能抛出异常，因此需要一个容器把它们全部保存下来。如果同一时刻只有一个线程抛出一个异常，那就不需要共享容器，也不需要\texttt{mutex}来同步对它的访问。你可以使用一个类型为\myttbr{std::\allowbreak{}exception\_\allowbreak{}ptr}的全局对象来保存线程之间传送的异常。


%%C080080 p

\myttbr{std::\allowbreak{}current\_\allowbreak{}exception()}是一个通常用在\texttt{catch}子句中的函数，用于捕获当前异常并创建\myttbr{std::\allowbreak{}exception\_\allowbreak{}ptr}实例。这样做的目的是保存指向原始异常的一个副本或引用（取决于实现），只要还存在引用它的\myttbr{std::\allowbreak{}exception\_\allowbreak{}ptr}指针，这个副本或引用就保持有效。如果在没有处理任何异常时调用该函数，它会创建一个空的\myttbr{std::\allowbreak{}exception\_\allowbreak{}ptr}。


%%C080081 p

\myttbr{std::\allowbreak{}exception\_\allowbreak{}ptr}指针是用\myttbr{std::\allowbreak{}current\_\allowbreak{}exception()}捕获到的异常的外包装。如果采用默认构造，它不持有任何异常。如果两个这种类型的对象都为空，或者指向同一个异常对象，那么它们相等。\myttbr{std::\allowbreak{}exception\_\allowbreak{}ptr}对象可以传给其他线程，在那里重新抛出并在\texttt{try...\allowbreak{}catch}块中捕获。


%%C080082 p

\myttbr{std::\allowbreak{}rethrow\_\allowbreak{}exception()}是一个函数，它接受\myttbr{std::\allowbreak{}exception\_\allowbreak{}ptr}作为参数，并抛出该参数所引用的异常对象。


%%C080083 note

\begin{myNote}
\myttbr{std::\allowbreak{}current\_\allowbreak{}exception()}、\myttbr{std::\allowbreak{}rethrow\_\allowbreak{}exception()}和\myttbr{std::\allowbreak{}exception\_\allowbreak{}ptr}在C++11中都已可用。
\end{myNote}


%%C080084 p

在上一节的示例中，每个线程函数都对自己执行的全部代码使用了\texttt{try...\allowbreak{}catch}语句，这样任何异常都不会未经捕获就离开函数。处理到异常时，先获取对全局\texttt{mutex}对象的锁，再把持有当前异常的\myttbr{std::\allowbreak{}exception\_\allowbreak{}ptr}对象添加到共享容器中。采用这种做法，线程函数会在第一个异常处停止；但在其他情况下，即使前一个操作抛出异常，你可能仍需要执行多个操作。这时你会有多个\texttt{try...\allowbreak{}catch}语句，并且可能只把其中一部分异常传送到线程之外。在主线程中，等所有线程执行完毕后，遍历容器，把每个非空的异常重新抛出，并用\texttt{try...\allowbreak{}catch}块捕获后妥善处理。


%%C080085 h3

\subsection{延伸阅读}


%%C080086 li

\begin{itemize}
\item \textit{使用线程}


\end{itemize}


%%C080087 h2

\section{用互斥量与锁同步对共享数据的访问}


%%C080088 p

线程使你能够同时执行多个函数，但这些函数常常需要访问共享资源。对共享资源的访问必须同步，以便同一时刻只有一个线程能够读或写该共享资源。上一案例中已经展示了这样的例子：多个线程都能同时向同一个共享容器添加对象。本案例将说明C++标准定义了哪些机制来同步线程对共享数据的访问，以及它们是如何工作的。


%%C080089 h3

\subsection{准备工作}


%%C080090 p

本案例讨论的\texttt{mutex}和\texttt{lock}类位于\texttt{<\allowbreak{}mutex>}头文件的\texttt{std}命名空间中。


%%C080091 h3

\subsection{具体做法}


%%C080092 p

对单个共享资源的同步访问使用以下模式：


%%C080093 p

在适当的上下文（类作用域或全局作用域）中定义一个\texttt{mutex}：


%%C080094 code

\begin{codeblockbg}
\begin{Verbatim}[commandchars=\\\{\},breaklines=true,breakanywhere=true,fontsize=\codefont,xleftmargin=17bp]
\PYG{n}{std}\PYG{o}{:}\PYG{o}{:}\PYG{n}{mutex}\PYG{+w}{ }\PYG{n}{g\PYGZus{}mutex}\PYG{p}{;}
\end{Verbatim}
\end{codeblockbg}


%%C080095 p

在各自线程中访问共享资源之前，先获取\texttt{mutex}上的\texttt{lock}：


%%C080096 code

\begin{codeblockbg}
\begin{Verbatim}[commandchars=\\\{\},breaklines=true,breakanywhere=true,fontsize=\codefont,xleftmargin=17bp]
\PYG{k+kt}{void}\PYG{+w}{ }\PYG{n+nf}{thread\PYGZus{}func}\PYG{p}{(}\PYG{p}{)}
\PYG{p}{\PYGZob{}}
\PYG{+w}{  }\PYG{k}{using}\PYG{+w}{ }\PYG{k}{namespace}\PYG{+w}{ }\PYG{n+nn}{std}\PYG{o}{:}\PYG{o}{:}\PYG{n+nn}{chrono\PYGZus{}literals}\PYG{p}{;}
\PYG{+w}{  }\PYG{p}{\PYGZob{}}
\PYG{+w}{    }\PYG{n}{std}\PYG{o}{:}\PYG{o}{:}\PYG{n}{lock\PYGZus{}guard}\PYG{o}{\PYGZlt{}}\PYG{n}{std}\PYG{o}{:}\PYG{o}{:}\PYG{n}{mutex}\PYG{o}{\PYGZgt{}}\PYG{+w}{ }\PYG{n}{lock}\PYG{p}{(}\PYG{n}{g\PYGZus{}mutex}\PYG{p}{)}\PYG{p}{;}
\PYG{+w}{    }\PYG{n}{std}\PYG{o}{:}\PYG{o}{:}\PYG{n}{cout}\PYG{+w}{ }\PYG{o}{\PYGZlt{}}\PYG{o}{\PYGZlt{}}\PYG{+w}{ }\PYG{l+s}{\PYGZdq{}}\PYG{l+s}{running thread }\PYG{l+s}{\PYGZdq{}}
\PYG{+w}{              }\PYG{o}{\PYGZlt{}}\PYG{o}{\PYGZlt{}}\PYG{+w}{ }\PYG{n}{std}\PYG{o}{:}\PYG{o}{:}\PYG{n}{this\PYGZus{}thread}\PYG{o}{:}\PYG{o}{:}\PYG{n}{get\PYGZus{}id}\PYG{p}{(}\PYG{p}{)}\PYG{+w}{ }\PYG{o}{\PYGZlt{}}\PYG{o}{\PYGZlt{}}\PYG{+w}{ }\PYG{n}{std}\PYG{o}{:}\PYG{o}{:}\PYG{n}{endl}\PYG{p}{;}
\PYG{+w}{  }\PYG{p}{\PYGZcb{}}
\PYG{n}{std}\PYG{o}{:}\PYG{o}{:}\PYG{n}{this\PYGZus{}thread}\PYG{o}{:}\PYG{o}{:}\PYG{n}{yield}\PYG{p}{(}\PYG{p}{)}\PYG{p}{;}
\PYG{n}{std}\PYG{o}{:}\PYG{o}{:}\PYG{n}{this\PYGZus{}thread}\PYG{o}{:}\PYG{o}{:}\PYG{n}{sleep\PYGZus{}for}\PYG{p}{(}\PYG{l+m+mi}{2}\PYG{n}{s}\PYG{p}{)}\PYG{p}{;}
\PYG{+w}{  }\PYG{p}{\PYGZob{}}
\PYG{+w}{    }\PYG{n}{std}\PYG{o}{:}\PYG{o}{:}\PYG{n}{lock\PYGZus{}guard}\PYG{o}{\PYGZlt{}}\PYG{n}{std}\PYG{o}{:}\PYG{o}{:}\PYG{n}{mutex}\PYG{o}{\PYGZgt{}}\PYG{+w}{ }\PYG{n}{lock}\PYG{p}{(}\PYG{n}{g\PYGZus{}mutex}\PYG{p}{)}\PYG{p}{;}
\PYG{+w}{    }\PYG{n}{std}\PYG{o}{:}\PYG{o}{:}\PYG{n}{cout}\PYG{+w}{ }\PYG{o}{\PYGZlt{}}\PYG{o}{\PYGZlt{}}\PYG{+w}{ }\PYG{l+s}{\PYGZdq{}}\PYG{l+s}{done in thread }\PYG{l+s}{\PYGZdq{}}
\PYG{+w}{              }\PYG{o}{\PYGZlt{}}\PYG{o}{\PYGZlt{}}\PYG{+w}{ }\PYG{n}{std}\PYG{o}{:}\PYG{o}{:}\PYG{n}{this\PYGZus{}thread}\PYG{o}{:}\PYG{o}{:}\PYG{n}{get\PYGZus{}id}\PYG{p}{(}\PYG{p}{)}\PYG{+w}{ }\PYG{o}{\PYGZlt{}}\PYG{o}{\PYGZlt{}}\PYG{+w}{ }\PYG{n}{std}\PYG{o}{:}\PYG{o}{:}\PYG{n}{endl}\PYG{p}{;}
\PYG{+w}{  }\PYG{p}{\PYGZcb{}}
\PYG{p}{\PYGZcb{}}
\end{Verbatim}
\end{codeblockbg}


%%C080097 p

要同时同步对多个共享资源的访问以避免死锁，请使用以下模式：


%%C080098 p

在适当的上下文（全局或


%%C080099 p

类作用域）中为每个共享资源定义一个互斥量：


%%C080100 code

\begin{codeblockbg}
\begin{Verbatim}[commandchars=\\\{\},breaklines=true,breakanywhere=true,fontsize=\codefont,xleftmargin=17bp]
\PYG{k}{template}\PYG{+w}{ }\PYG{o}{\PYGZlt{}}\PYG{k}{typename}\PYG{+w}{ }\PYG{n+nc}{T}\PYG{o}{\PYGZgt{}}
\PYG{k}{struct}\PYG{+w}{ }\PYG{n+nc}{container}
\PYG{p}{\PYGZob{}}
\PYG{+w}{  }\PYG{n}{std}\PYG{o}{:}\PYG{o}{:}\PYG{n}{mutex}\PYG{+w}{     }\PYG{n}{mutex}\PYG{p}{;}
\PYG{+w}{  }\PYG{n}{std}\PYG{o}{:}\PYG{o}{:}\PYG{n}{vector}\PYG{o}{\PYGZlt{}}\PYG{n}{T}\PYG{o}{\PYGZgt{}}\PYG{+w}{ }\PYG{n}{data}\PYG{p}{;}
\PYG{p}{\PYGZcb{}}\PYG{p}{;}
\end{Verbatim}
\end{codeblockbg}


%%C080101 p

使用避免死锁的算法，借助


%%C080102 p

\texttt{std::\allowbreak{}lock()}同时锁定这些互斥量：


%%C080103 code

\begin{codeblockbg}
\begin{Verbatim}[commandchars=\\\{\},breaklines=true,breakanywhere=true,fontsize=\codefont,xleftmargin=17bp]
\PYG{k}{template}\PYG{+w}{ }\PYG{o}{\PYGZlt{}}\PYG{k}{typename}\PYG{+w}{ }\PYG{n+nc}{T}\PYG{o}{\PYGZgt{}}
\PYG{k+kt}{void}\PYG{+w}{ }\PYG{n}{move\PYGZus{}between}\PYG{p}{(}\PYG{n}{container}\PYG{o}{\PYGZlt{}}\PYG{n}{T}\PYG{o}{\PYGZgt{}}\PYG{+w}{ }\PYG{o}{\PYGZam{}}\PYG{+w}{ }\PYG{n}{c1}\PYG{p}{,}\PYG{+w}{ }\PYG{n}{container}\PYG{o}{\PYGZlt{}}\PYG{n}{T}\PYG{o}{\PYGZgt{}}\PYG{+w}{ }\PYG{o}{\PYGZam{}}\PYG{+w}{ }\PYG{n}{c2}\PYG{p}{,}
\PYG{+w}{                  }\PYG{n}{T}\PYG{+w}{ }\PYG{k}{const}\PYG{+w}{ }\PYG{n}{value}\PYG{p}{)}
\PYG{p}{\PYGZob{}}
\PYG{+w}{  }\PYG{n}{std}\PYG{o}{:}\PYG{o}{:}\PYG{n}{lock}\PYG{p}{(}\PYG{n}{c1}\PYG{p}{.}\PYG{n}{mutex}\PYG{p}{,}\PYG{+w}{ }\PYG{n}{c2}\PYG{p}{.}\PYG{n}{mutex}\PYG{p}{)}\PYG{p}{;}
\PYG{+w}{  }\PYG{c+c1}{// 接第3步。}
\PYG{p}{\PYGZcb{}}
\end{Verbatim}
\end{codeblockbg}


%%C080104 p

锁定之后，把每个互斥量的所有权移交给


%%C080105 p

\texttt{std::\allowbreak{}lock\_\allowbreak{}guard}类，以确保它们在函数（或作用域）结束时被安全释放：


%%C080106 code

\begin{codeblockbg}
\begin{Verbatim}[commandchars=\\\{\},breaklines=true,breakanywhere=true,fontsize=\codefont,xleftmargin=17bp]
\PYG{c+c1}{// 接第2步。}
\PYG{n}{std}\PYG{o}{:}\PYG{o}{:}\PYG{n}{lock\PYGZus{}guard}\PYG{o}{\PYGZlt{}}\PYG{n}{std}\PYG{o}{:}\PYG{o}{:}\PYG{n}{mutex}\PYG{o}{\PYGZgt{}}\PYG{+w}{ }\PYG{n}{l1}\PYG{p}{(}\PYG{n}{c1}\PYG{p}{.}\PYG{n}{mutex}\PYG{p}{,}\PYG{+w}{ }\PYG{n}{std}\PYG{o}{:}\PYG{o}{:}\PYG{n}{adopt\PYGZus{}lock}\PYG{p}{)}\PYG{p}{;}
\PYG{n}{std}\PYG{o}{:}\PYG{o}{:}\PYG{n}{lock\PYGZus{}guard}\PYG{o}{\PYGZlt{}}\PYG{n}{std}\PYG{o}{:}\PYG{o}{:}\PYG{n}{mutex}\PYG{o}{\PYGZgt{}}\PYG{+w}{ }\PYG{n}{l2}\PYG{p}{(}\PYG{n}{c2}\PYG{p}{.}\PYG{n}{mutex}\PYG{p}{,}\PYG{+w}{ }\PYG{n}{std}\PYG{o}{:}\PYG{o}{:}\PYG{n}{adopt\PYGZus{}lock}\PYG{p}{)}\PYG{p}{;}
\PYG{n}{c1}\PYG{p}{.}\PYG{n}{data}\PYG{p}{.}\PYG{n}{erase}\PYG{p}{(}
\PYG{+w}{  }\PYG{n}{std}\PYG{o}{:}\PYG{o}{:}\PYG{n}{remove}\PYG{p}{(}\PYG{n}{c1}\PYG{p}{.}\PYG{n}{data}\PYG{p}{.}\PYG{n}{begin}\PYG{p}{(}\PYG{p}{)}\PYG{p}{,}\PYG{+w}{ }\PYG{n}{c1}\PYG{p}{.}\PYG{n}{data}\PYG{p}{.}\PYG{n}{end}\PYG{p}{(}\PYG{p}{)}\PYG{p}{,}\PYG{+w}{ }\PYG{n}{value}\PYG{p}{)}\PYG{p}{,}
\PYG{+w}{  }\PYG{n}{c1}\PYG{p}{.}\PYG{n}{data}\PYG{p}{.}\PYG{n}{end}\PYG{p}{(}\PYG{p}{)}\PYG{p}{)}\PYG{p}{;}
\PYG{n}{c2}\PYG{p}{.}\PYG{n}{data}\PYG{p}{.}\PYG{n}{push\PYGZus{}back}\PYG{p}{(}\PYG{n}{value}\PYG{p}{)}\PYG{p}{;}
\end{Verbatim}
\end{codeblockbg}


%%C080107 h3

\subsection{工作原理}


%%C080108 p

互斥量是一种同步原语，使我们可以保护共享资源不被多个线程同时访问。C++标准库提供了几种实现：

%%C080109 li

\begin{itemize}
\item \texttt{std::\allowbreak{}mutex}是最常用的互斥量类型，前面的代码片段演示的正是它。它提供了获取与释放互斥量的方法：\texttt{lock()}尝试获取互斥量，若不可用则阻塞，\texttt{try\_\allowbreak{}lock()}尝试获取互斥量，若不可用则不加阻塞地立即返回，\texttt{unlock()}释放互斥量。


%%C080110 li

\item \texttt{std::\allowbreak{}timed\_\allowbreak{}mutex}与\texttt{std::\allowbreak{}mutex}类似，但多了两个带超时的获取方法：\texttt{try\_\allowbreak{}lock\_\allowbreak{}for()}尝试获取互斥量，若在指定的时长内互斥量一直不可用就返回，\texttt{try\_\allowbreak{}lock\_\allowbreak{}until()}尝试获取互斥量，若直到指定时间点互斥量仍不可用就返回。


%%C080111 li

\item \myttbr{std::\allowbreak{}recursive\_\allowbreak{}mutex}与\texttt{std::\allowbreak{}mutex}类似，但同一个线程可以多次获取该互斥量而不会被阻塞。


%%C080112 li

\item \myttbr{std::\allowbreak{}recursive\_\allowbreak{}timed\_\allowbreak{}mutex}是递归互斥量与定时互斥量的结合。

\end{itemize}


%%C080113 p

第一个锁定某个可用互斥量的线程获得它的所有权并继续执行。此后任何线程对它的锁定尝试都会失败，包括已经拥有该互斥量的那个线程，\texttt{lock()}方法会阻塞线程，直到有人调用\texttt{unlock()}释放互斥量。如果某个线程需要能够多次锁定同一个互斥量而不被阻塞、从而避免把自己锁死，就应当使用\texttt{recursive\_\allowbreak{}mutex}类模板。


%%C080114 p

用互斥量保护对共享资源的访问，典型做法是锁定互斥量、使用共享资源，然后解锁互斥量：


%%C080115 code

\begin{codeblockbg}
\begin{Verbatim}[commandchars=\\\{\},breaklines=true,breakanywhere=true,fontsize=\codefont,xleftmargin=17bp]
\PYG{n}{g\PYGZus{}mutex}\PYG{p}{.}\PYG{n}{lock}\PYG{p}{(}\PYG{p}{)}\PYG{p}{;}
\PYG{c+c1}{// 使用共享资源，例如 std::cout}
\PYG{n}{std}\PYG{o}{:}\PYG{o}{:}\PYG{n}{cout}\PYG{+w}{ }\PYG{o}{\PYGZlt{}}\PYG{o}{\PYGZlt{}}\PYG{+w}{ }\PYG{l+s}{\PYGZdq{}}\PYG{l+s}{accessing shared resource}\PYG{l+s}{\PYGZdq{}}\PYG{+w}{ }\PYG{o}{\PYGZlt{}}\PYG{o}{\PYGZlt{}}\PYG{+w}{ }\PYG{n}{std}\PYG{o}{:}\PYG{o}{:}\PYG{n}{endl}\PYG{p}{;}
\PYG{n}{g\PYGZus{}mutex}\PYG{p}{.}\PYG{n}{unlock}\PYG{p}{(}\PYG{p}{)}\PYG{p}{;}
\end{Verbatim}
\end{codeblockbg}


%%C080116 p

然而这种使用互斥量的方式很容易出错。因为每一次\texttt{lock()}调用都必须在所有执行路径上与一次\texttt{unlock()}调用配对，既包括正常返回路径，也包括异常返回路径。为了无论函数以何种方式结束执行都能安全地获取与释放互斥量，C++标准定义了若干加锁类：


%%C080117 li

\begin{itemize}
\item \texttt{std::\allowbreak{}lock\_\allowbreak{}guard}就是前面见到的加锁机制，它是以RAII方式实现的互斥量包装器。它在构造时尝试获取互斥量，在析构时释放互斥量。C++11起可用。下面是\texttt{lock\_\allowbreak{}guard}的一个典型实现：


%%C080118 code

\begin{codeblockbg}
\begin{Verbatim}[commandchars=\\\{\},breaklines=true,breakanywhere=true,fontsize=\codefont,xleftmargin=17bp]
\PYG{k}{template}\PYG{+w}{ }\PYG{o}{\PYGZlt{}}\PYG{k}{class}\PYG{+w}{ }\PYG{n+nc}{M}\PYG{o}{\PYGZgt{}}
\PYG{k}{class}\PYG{+w}{ }\PYG{n+nc}{lock\PYGZus{}guard}
\PYG{p}{\PYGZob{}}
\PYG{k}{public}\PYG{o}{:}
\PYG{+w}{  }\PYG{k}{typedef}\PYG{+w}{ }\PYG{n}{M}\PYG{+w}{ }\PYG{n}{mutex\PYGZus{}type}\PYG{p}{;}
\PYG{k}{explicit}\PYG{+w}{ }\PYG{n}{lock\PYGZus{}guard}\PYG{p}{(}\PYG{n}{M}\PYG{o}{\PYGZam{}}\PYG{+w}{ }\PYG{n}{Mtx}\PYG{p}{)}\PYG{+w}{ }\PYG{o}{:}\PYG{+w}{ }\PYG{n}{mtx}\PYG{p}{(}\PYG{n}{Mtx}\PYG{p}{)}
\PYG{p}{\PYGZob{}}
\PYG{+w}{  }\PYG{n}{mtx}\PYG{p}{.}\PYG{n}{lock}\PYG{p}{(}\PYG{p}{)}\PYG{p}{;}
\PYG{p}{\PYGZcb{}}
\PYG{n}{lock\PYGZus{}guard}\PYG{p}{(}\PYG{n}{M}\PYG{o}{\PYGZam{}}\PYG{+w}{ }\PYG{n}{Mtx}\PYG{p}{,}\PYG{+w}{ }\PYG{n}{std}\PYG{o}{:}\PYG{o}{:}\PYG{n}{adopt\PYGZus{}lock\PYGZus{}t}\PYG{p}{)}\PYG{+w}{ }\PYG{o}{:}\PYG{+w}{ }\PYG{n}{mtx}\PYG{p}{(}\PYG{n}{Mtx}\PYG{p}{)}
\PYG{p}{\PYGZob{}}\PYG{+w}{ }\PYG{p}{\PYGZcb{}}
\PYG{o}{\PYGZti{}}\PYG{n}{lock\PYGZus{}guard}\PYG{p}{(}\PYG{p}{)}\PYG{+w}{ }\PYG{k}{noexcept}
\PYG{p}{\PYGZob{}}
\PYG{+w}{  }\PYG{n}{mtx}\PYG{p}{.}\PYG{n}{unlock}\PYG{p}{(}\PYG{p}{)}\PYG{p}{;}
\PYG{p}{\PYGZcb{}}
\PYG{+w}{  }\PYG{n}{lock\PYGZus{}guard}\PYG{p}{(}\PYG{k}{const}\PYG{+w}{ }\PYG{n}{lock\PYGZus{}guard}\PYG{o}{\PYGZam{}}\PYG{p}{)}\PYG{+w}{ }\PYG{o}{=}\PYG{+w}{ }\PYG{k}{delete}\PYG{p}{;}
\PYG{+w}{  }\PYG{n}{lock\PYGZus{}guard}\PYG{o}{\PYGZam{}}\PYG{+w}{ }\PYG{k}{operator}\PYG{o}{=}\PYG{p}{(}\PYG{k}{const}\PYG{+w}{ }\PYG{n}{lock\PYGZus{}guard}\PYG{o}{\PYGZam{}}\PYG{p}{)}\PYG{+w}{ }\PYG{o}{=}\PYG{+w}{ }\PYG{k}{delete}\PYG{p}{;}
\PYG{k}{private}\PYG{o}{:}
\PYG{+w}{  }\PYG{n}{M}\PYG{o}{\PYGZam{}}\PYG{+w}{ }\PYG{n}{mtx}\PYG{p}{;}
\PYG{p}{\PYGZcb{}}\PYG{p}{;}
\end{Verbatim}
\end{codeblockbg}


%%C080119 li

\item \texttt{std::\allowbreak{}unique\_\allowbreak{}lock}是互斥量所有权的包装器，支持延迟加锁、定时加锁、递归加锁、所有权转移，以及配合条件变量使用。C++11起可用。


%%C080120 li

\item \texttt{std::\allowbreak{}shared\_\allowbreak{}lock}是互斥量共享所有权的包装器，支持延迟加锁、定时加锁与所有权转移。C++14起可用。


%%C080121 li

\item \texttt{std::\allowbreak{}scoped\_\allowbreak{}lock}是以RAII方式实现的多互斥量包装器：构造时它如同使用\texttt{std::\allowbreak{}lock()}那样，以死锁避免的方式尝试获取各个互斥量的所有权，析构时按获取顺序的逆序释放这些互斥量。C++17起可用。

\end{itemize}


%%C080122 p

在\textit{具体做法}一节的第一个示例中，我们用\texttt{std::\allowbreak{}mutex}和\texttt{std::\allowbreak{}lock\_\allowbreak{}guard}保护对\texttt{std::\allowbreak{}cout}流对象的访问，该对象由程序中的所有线程共享。下面的示例展示了如何在多个线程上并发执行\texttt{thread\_\allowbreak{}func()}函数：


%%C080123 code

\begin{codeblockbg}
\begin{Verbatim}[commandchars=\\\{\},breaklines=true,breakanywhere=true,fontsize=\codefont,xleftmargin=17bp]
\PYG{n}{std}\PYG{o}{:}\PYG{o}{:}\PYG{n}{vector}\PYG{o}{\PYGZlt{}}\PYG{n}{std}\PYG{o}{:}\PYG{o}{:}\PYG{k+kr}{thread}\PYG{o}{\PYGZgt{}}\PYG{+w}{ }\PYG{n}{threads}\PYG{p}{;}
\PYG{k}{for}\PYG{+w}{ }\PYG{p}{(}\PYG{k+kt}{int}\PYG{+w}{ }\PYG{n}{i}\PYG{+w}{ }\PYG{o}{=}\PYG{+w}{ }\PYG{l+m+mi}{0}\PYG{p}{;}\PYG{+w}{ }\PYG{n}{i}\PYG{+w}{ }\PYG{o}{\PYGZlt{}}\PYG{+w}{ }\PYG{l+m+mi}{5}\PYG{p}{;}\PYG{+w}{ }\PYG{o}{+}\PYG{o}{+}\PYG{n}{i}\PYG{p}{)}
\PYG{+w}{  }\PYG{n}{threads}\PYG{p}{.}\PYG{n}{emplace\PYGZus{}back}\PYG{p}{(}\PYG{n}{thread\PYGZus{}func}\PYG{p}{)}\PYG{p}{;}
\PYG{k}{for}\PYG{+w}{ }\PYG{p}{(}\PYG{k}{auto}\PYG{+w}{ }\PYG{o}{\PYGZam{}}\PYG{+w}{ }\PYG{n}{t}\PYG{+w}{ }\PYG{o}{:}\PYG{+w}{ }\PYG{n}{threads}\PYG{p}{)}
\PYG{+w}{  }\PYG{n}{t}\PYG{p}{.}\PYG{n}{join}\PYG{p}{(}\PYG{p}{)}\PYG{p}{;}
\end{Verbatim}
\end{codeblockbg}


%%C080124 p

程序的一种可能输出如下：


%%C080125 code

\begin{codeblockbg}
\begin{Verbatim}[commandchars=\\\{\},breaklines=true,breakanywhere=true,fontsize=\codefont,xleftmargin=17bp]
\PYG{n}{running}\PYG{+w}{ }\PYG{k+kr}{thread}\PYG{+w}{ }\PYG{l+m+mi}{140296854550272}
\PYG{n}{running}\PYG{+w}{ }\PYG{k+kr}{thread}\PYG{+w}{ }\PYG{l+m+mi}{140296846157568}
\PYG{n}{running}\PYG{+w}{ }\PYG{k+kr}{thread}\PYG{+w}{ }\PYG{l+m+mi}{140296837764864}
\PYG{n}{running}\PYG{+w}{ }\PYG{k+kr}{thread}\PYG{+w}{ }\PYG{l+m+mi}{140296829372160}
\PYG{n}{running}\PYG{+w}{ }\PYG{k+kr}{thread}\PYG{+w}{ }\PYG{l+m+mi}{140296820979456}
\PYG{n}{done}\PYG{+w}{ }\PYG{n}{in}\PYG{+w}{ }\PYG{k+kr}{thread}\PYG{+w}{ }\PYG{l+m+mi}{140296854550272}
\PYG{n}{done}\PYG{+w}{ }\PYG{n}{in}\PYG{+w}{ }\PYG{k+kr}{thread}\PYG{+w}{ }\PYG{l+m+mi}{140296846157568}
\PYG{n}{done}\PYG{+w}{ }\PYG{n}{in}\PYG{+w}{ }\PYG{k+kr}{thread}\PYG{+w}{ }\PYG{l+m+mi}{140296837764864}
\PYG{n}{done}\PYG{+w}{ }\PYG{n}{in}\PYG{+w}{ }\PYG{k+kr}{thread}\PYG{+w}{ }\PYG{l+m+mi}{140296820979456}
\PYG{n}{done}\PYG{+w}{ }\PYG{n}{in}\PYG{+w}{ }\PYG{k+kr}{thread}\PYG{+w}{ }\PYG{l+m+mi}{140296829372160}
\end{Verbatim}
\end{codeblockbg}


%%C080126 p

当某个线程需要获取多个用于保护不同共享资源的互斥量时，逐个获取可能导致死锁。考虑下面的示例（其中\texttt{container}是\textit{具体做法}一节展示的类）：


%%C080127 code

\begin{codeblockbg}
\begin{Verbatim}[commandchars=\\\{\},breaklines=true,breakanywhere=true,fontsize=\codefont,xleftmargin=17bp]
\PYG{k}{template}\PYG{+w}{ }\PYG{o}{\PYGZlt{}}\PYG{k}{typename}\PYG{+w}{ }\PYG{n+nc}{T}\PYG{o}{\PYGZgt{}}
\PYG{k+kt}{void}\PYG{+w}{ }\PYG{n}{move\PYGZus{}between}\PYG{p}{(}\PYG{n}{container}\PYG{o}{\PYGZlt{}}\PYG{n}{T}\PYG{o}{\PYGZgt{}}\PYG{+w}{ }\PYG{o}{\PYGZam{}}\PYG{+w}{ }\PYG{n}{c1}\PYG{p}{,}\PYG{+w}{ }\PYG{n}{container}\PYG{o}{\PYGZlt{}}\PYG{n}{T}\PYG{o}{\PYGZgt{}}\PYG{+w}{ }\PYG{o}{\PYGZam{}}\PYG{+w}{ }\PYG{n}{c2}\PYG{p}{,}\PYG{+w}{ }\PYG{n}{T}\PYG{+w}{ }\PYG{k}{const}\PYG{+w}{ }\PYG{n}{value}\PYG{p}{)}
\PYG{p}{\PYGZob{}}
\PYG{+w}{  }\PYG{n}{std}\PYG{o}{:}\PYG{o}{:}\PYG{n}{lock\PYGZus{}guard}\PYG{o}{\PYGZlt{}}\PYG{n}{std}\PYG{o}{:}\PYG{o}{:}\PYG{n}{mutex}\PYG{o}{\PYGZgt{}}\PYG{+w}{ }\PYG{n}{l1}\PYG{p}{(}\PYG{n}{c1}\PYG{p}{.}\PYG{n}{mutex}\PYG{p}{)}\PYG{p}{;}
\PYG{+w}{  }\PYG{n}{std}\PYG{o}{:}\PYG{o}{:}\PYG{n}{lock\PYGZus{}guard}\PYG{o}{\PYGZlt{}}\PYG{n}{std}\PYG{o}{:}\PYG{o}{:}\PYG{n}{mutex}\PYG{o}{\PYGZgt{}}\PYG{+w}{ }\PYG{n}{l2}\PYG{p}{(}\PYG{n}{c2}\PYG{p}{.}\PYG{n}{mutex}\PYG{p}{)}\PYG{p}{;}
\PYG{+w}{  }\PYG{n}{c1}\PYG{p}{.}\PYG{n}{data}\PYG{p}{.}\PYG{n}{erase}\PYG{p}{(}
\PYG{+w}{    }\PYG{n}{std}\PYG{o}{:}\PYG{o}{:}\PYG{n}{remove}\PYG{p}{(}\PYG{n}{c1}\PYG{p}{.}\PYG{n}{data}\PYG{p}{.}\PYG{n}{begin}\PYG{p}{(}\PYG{p}{)}\PYG{p}{,}\PYG{+w}{ }\PYG{n}{c1}\PYG{p}{.}\PYG{n}{data}\PYG{p}{.}\PYG{n}{end}\PYG{p}{(}\PYG{p}{)}\PYG{p}{,}\PYG{+w}{ }\PYG{n}{value}\PYG{p}{)}\PYG{p}{,}
\PYG{+w}{    }\PYG{n}{c1}\PYG{p}{.}\PYG{n}{data}\PYG{p}{.}\PYG{n}{end}\PYG{p}{(}\PYG{p}{)}\PYG{p}{)}\PYG{p}{;}
\PYG{+w}{  }\PYG{n}{c2}\PYG{p}{.}\PYG{n}{data}\PYG{p}{.}\PYG{n}{push\PYGZus{}back}\PYG{p}{(}\PYG{n}{value}\PYG{p}{)}\PYG{p}{;}
\PYG{p}{\PYGZcb{}}
\PYG{n}{container}\PYG{o}{\PYGZlt{}}\PYG{k+kt}{int}\PYG{o}{\PYGZgt{}}\PYG{+w}{ }\PYG{n}{c1}\PYG{p}{;}
\PYG{n}{c1}\PYG{p}{.}\PYG{n}{data}\PYG{p}{.}\PYG{n}{push\PYGZus{}back}\PYG{p}{(}\PYG{l+m+mi}{1}\PYG{p}{)}\PYG{p}{;}
\PYG{n}{c1}\PYG{p}{.}\PYG{n}{data}\PYG{p}{.}\PYG{n}{push\PYGZus{}back}\PYG{p}{(}\PYG{l+m+mi}{2}\PYG{p}{)}\PYG{p}{;}
\PYG{n}{c1}\PYG{p}{.}\PYG{n}{data}\PYG{p}{.}\PYG{n}{push\PYGZus{}back}\PYG{p}{(}\PYG{l+m+mi}{3}\PYG{p}{)}\PYG{p}{;}
\PYG{n}{container}\PYG{o}{\PYGZlt{}}\PYG{k+kt}{int}\PYG{o}{\PYGZgt{}}\PYG{+w}{ }\PYG{n}{c2}\PYG{p}{;}
\PYG{n}{c2}\PYG{p}{.}\PYG{n}{data}\PYG{p}{.}\PYG{n}{push\PYGZus{}back}\PYG{p}{(}\PYG{l+m+mi}{4}\PYG{p}{)}\PYG{p}{;}
\PYG{n}{c2}\PYG{p}{.}\PYG{n}{data}\PYG{p}{.}\PYG{n}{push\PYGZus{}back}\PYG{p}{(}\PYG{l+m+mi}{5}\PYG{p}{)}\PYG{p}{;}
\PYG{n}{c2}\PYG{p}{.}\PYG{n}{data}\PYG{p}{.}\PYG{n}{push\PYGZus{}back}\PYG{p}{(}\PYG{l+m+mi}{6}\PYG{p}{)}\PYG{p}{;}
\PYG{n}{std}\PYG{o}{:}\PYG{o}{:}\PYG{k+kr}{thread}\PYG{+w}{ }\PYG{n+nf}{t1}\PYG{p}{(}\PYG{n}{move\PYGZus{}between}\PYG{o}{\PYGZlt{}}\PYG{k+kt}{int}\PYG{o}{\PYGZgt{}}\PYG{p}{,}\PYG{+w}{ }\PYG{n}{std}\PYG{o}{:}\PYG{o}{:}\PYG{n}{ref}\PYG{p}{(}\PYG{n}{c1}\PYG{p}{)}\PYG{p}{,}\PYG{+w}{ }\PYG{n}{std}\PYG{o}{:}\PYG{o}{:}\PYG{n}{ref}\PYG{p}{(}\PYG{n}{c2}\PYG{p}{)}\PYG{p}{,}\PYG{+w}{ }\PYG{l+m+mi}{3}\PYG{p}{)}\PYG{p}{;}
\PYG{n}{std}\PYG{o}{:}\PYG{o}{:}\PYG{k+kr}{thread}\PYG{+w}{ }\PYG{n+nf}{t2}\PYG{p}{(}\PYG{n}{move\PYGZus{}between}\PYG{o}{\PYGZlt{}}\PYG{k+kt}{int}\PYG{o}{\PYGZgt{}}\PYG{p}{,}\PYG{+w}{ }\PYG{n}{std}\PYG{o}{:}\PYG{o}{:}\PYG{n}{ref}\PYG{p}{(}\PYG{n}{c2}\PYG{p}{)}\PYG{p}{,}\PYG{+w}{ }\PYG{n}{std}\PYG{o}{:}\PYG{o}{:}\PYG{n}{ref}\PYG{p}{(}\PYG{n}{c1}\PYG{p}{)}\PYG{p}{,}\PYG{+w}{ }\PYG{l+m+mi}{6}\PYG{p}{)}\PYG{p}{;}
\PYG{n}{t1}\PYG{p}{.}\PYG{n}{join}\PYG{p}{(}\PYG{p}{)}\PYG{p}{;}
\PYG{n}{t2}\PYG{p}{.}\PYG{n}{join}\PYG{p}{(}\PYG{p}{)}\PYG{p}{;}
\end{Verbatim}
\end{codeblockbg}


%%C080128 p

在这个示例中，\texttt{container}类持有的数据可能被不同线程同时访问，因此需要用互斥量加锁来保护。\texttt{move\_\allowbreak{}between()}是一个线程安全的函数，它从某个容器中移除一个元素并把它加到另一个容器中。为此，它依次获取两个容器的互斥量，然后从第一个容器中删除该元素，再把它追加到第二个容器末尾。


%%C080129 p

然而这个函数容易发生死锁，因为加锁过程中可能触发竞态条件。假设有如下场景：两个不同的线程执行该函数，但传入参数的顺序不同：


%%C080130 li

\begin{itemize}
\item 第一个线程以\texttt{c1}、\texttt{c2}的顺序作为参数开始执行。


%%C080131 li

\item 第一个线程在获取\texttt{c1}容器的锁之后被挂起。第二个线程以\texttt{c2}、\texttt{c1}的顺序作为参数开始执行。


%%C080132 li

\item 第二个线程在获取\texttt{c2}容器的锁之后被挂起。


%%C080133 li

\item 第一个线程继续执行，尝试获取\texttt{c2}的互斥量，但该互斥量不可用。于是发生死锁（可以在获取第一个互斥量之后让线程休眠一小会儿来模拟这种情况）。

\end{itemize}


%%C080134 p

为了避免这类可能的死锁，应以死锁避免的方式获取互斥量，标准库为此提供了工具函数\texttt{std::\allowbreak{}lock()}。\texttt{move\_\allowbreak{}between()}函数需要把两个锁换成如下代码（如\textit{具体做法}一节所示）：


%%C080135 code

\begin{codeblockbg}
\begin{Verbatim}[commandchars=\\\{\},breaklines=true,breakanywhere=true,fontsize=\codefont,xleftmargin=17bp]
\PYG{n}{std}\PYG{o}{:}\PYG{o}{:}\PYG{n}{lock}\PYG{p}{(}\PYG{n}{c1}\PYG{p}{.}\PYG{n}{mutex}\PYG{p}{,}\PYG{+w}{ }\PYG{n}{c2}\PYG{p}{.}\PYG{n}{mutex}\PYG{p}{)}\PYG{p}{;}
\PYG{n}{std}\PYG{o}{:}\PYG{o}{:}\PYG{n}{lock\PYGZus{}guard}\PYG{o}{\PYGZlt{}}\PYG{n}{std}\PYG{o}{:}\PYG{o}{:}\PYG{n}{mutex}\PYG{o}{\PYGZgt{}}\PYG{+w}{ }\PYG{n}{l1}\PYG{p}{(}\PYG{n}{c1}\PYG{p}{.}\PYG{n}{mutex}\PYG{p}{,}\PYG{+w}{ }\PYG{n}{std}\PYG{o}{:}\PYG{o}{:}\PYG{n}{adopt\PYGZus{}lock}\PYG{p}{)}\PYG{p}{;}
\PYG{n}{std}\PYG{o}{:}\PYG{o}{:}\PYG{n}{lock\PYGZus{}guard}\PYG{o}{\PYGZlt{}}\PYG{n}{std}\PYG{o}{:}\PYG{o}{:}\PYG{n}{mutex}\PYG{o}{\PYGZgt{}}\PYG{+w}{ }\PYG{n}{l2}\PYG{p}{(}\PYG{n}{c2}\PYG{p}{.}\PYG{n}{mutex}\PYG{p}{,}\PYG{+w}{ }\PYG{n}{std}\PYG{o}{:}\PYG{o}{:}\PYG{n}{adopt\PYGZus{}lock}\PYG{p}{)}\PYG{p}{;}
\end{Verbatim}
\end{codeblockbg}


%%C080136 p

互斥量的所有权仍须转移给锁守卫对象，这样它们才能在函数执行结束（或视情况在某个特定作用域结束）之后被正确释放。


%%C080137 p

C++17引入了一个新的互斥量包装器\texttt{std::\allowbreak{}scoped\_\allowbreak{}lock}，可以用它简化上例那样的代码。这种锁能以无死锁的方式获取多个互斥量的所有权，这些互斥量在作用域锁析构时被释放。前面的代码等价于下面这一行代码：


%%C080138 code

\begin{codeblockbg}
\begin{Verbatim}[commandchars=\\\{\},breaklines=true,breakanywhere=true,fontsize=\codefont,xleftmargin=17bp]
\PYG{n}{std}\PYG{o}{:}\PYG{o}{:}\PYG{n}{scoped\PYGZus{}lock}\PYG{+w}{ }\PYG{n+nf}{lock}\PYG{p}{(}\PYG{n}{c1}\PYG{p}{.}\PYG{n}{mutex}\PYG{p}{,}\PYG{+w}{ }\PYG{n}{c2}\PYG{p}{.}\PYG{n}{mutex}\PYG{p}{)}\PYG{p}{;}
\end{Verbatim}
\end{codeblockbg}


%%C080139 h3

\subsection{延伸阅读}


%%C080140 li

\begin{itemize}
\item \textit{使用线程}


%%C080141 li

\item \textit{避免使用递归互斥量}

\end{itemize}


%%C080142 h2

\section{避免使用递归互斥量}


%%C080143 p

标准库提供了多种互斥量类型来保护对共享资源的访问。\myttbr{std::\allowbreak{}recursive\_\allowbreak{}mutex}和\myttbr{std::\allowbreak{}recursive\_\allowbreak{}timed\_\allowbreak{}mutex}是两种允许在同一线程中多次加锁的实现。递归互斥量的典型用途是在递归函数中保护对共享资源的访问。递归互斥量的开销比非递归互斥量大，因此应尽可能避免使用。本案例给出一个用例，把使用递归互斥量的线程安全类型改造成使用非递归互斥量的线程安全类型。


%%C080144 h3

\subsection{准备工作}


%%C080145 p

你需要熟悉标准库中提供的各种互斥量与锁。建议先阅读上一个案例\textit{用互斥量与锁同步对共享数据的访问}，以便对它们有个总体了解。


%%C080146 p

本案例的目标是改造下面这个类，从而避免使用\myttbr{std::\allowbreak{}recursive\_\allowbreak{}mutex}：


%%C080147 code

\begin{codeblockbg}
\begin{Verbatim}[commandchars=\\\{\},breaklines=true,breakanywhere=true,fontsize=\codefont,xleftmargin=17bp]
\PYG{k}{class}\PYG{+w}{ }\PYG{n+nc}{foo\PYGZus{}rec}
\PYG{p}{\PYGZob{}}
\PYG{+w}{  }\PYG{n}{std}\PYG{o}{:}\PYG{o}{:}\PYG{n}{recursive\PYGZus{}mutex}\PYG{+w}{ }\PYG{n}{m}\PYG{p}{;}
\PYG{+w}{  }\PYG{k+kt}{int}\PYG{+w}{ }\PYG{n}{data}\PYG{p}{;}
\PYG{k}{public}\PYG{o}{:}
\PYG{+w}{  }\PYG{n}{foo\PYGZus{}rec}\PYG{p}{(}\PYG{k+kt}{int}\PYG{+w}{ }\PYG{k}{const}\PYG{+w}{ }\PYG{n}{d}\PYG{+w}{ }\PYG{o}{=}\PYG{+w}{ }\PYG{l+m+mi}{0}\PYG{p}{)}\PYG{+w}{ }\PYG{o}{:}\PYG{+w}{ }\PYG{n}{data}\PYG{p}{(}\PYG{n}{d}\PYG{p}{)}\PYG{+w}{ }\PYG{p}{\PYGZob{}}\PYG{p}{\PYGZcb{}}
\PYG{k+kt}{void}\PYG{+w}{ }\PYG{n}{update}\PYG{p}{(}\PYG{k+kt}{int}\PYG{+w}{ }\PYG{k}{const}\PYG{+w}{ }\PYG{n}{d}\PYG{p}{)}
\PYG{p}{\PYGZob{}}
\PYG{+w}{  }\PYG{n}{std}\PYG{o}{:}\PYG{o}{:}\PYG{n}{lock\PYGZus{}guard}\PYG{o}{\PYGZlt{}}\PYG{n}{std}\PYG{o}{:}\PYG{o}{:}\PYG{n}{recursive\PYGZus{}mutex}\PYG{o}{\PYGZgt{}}\PYG{+w}{ }\PYG{n}{lock}\PYG{p}{(}\PYG{n}{m}\PYG{p}{)}\PYG{p}{;}
\PYG{+w}{  }\PYG{n}{data}\PYG{+w}{ }\PYG{o}{=}\PYG{+w}{ }\PYG{n}{d}\PYG{p}{;}
\PYG{p}{\PYGZcb{}}
\PYG{+w}{  }\PYG{k+kt}{int}\PYG{+w}{ }\PYG{n}{update\PYGZus{}with\PYGZus{}return}\PYG{p}{(}\PYG{k+kt}{int}\PYG{+w}{ }\PYG{k}{const}\PYG{+w}{ }\PYG{n}{d}\PYG{p}{)}
\PYG{+w}{  }\PYG{p}{\PYGZob{}}
\PYG{+w}{    }\PYG{n}{std}\PYG{o}{:}\PYG{o}{:}\PYG{n}{lock\PYGZus{}guard}\PYG{o}{\PYGZlt{}}\PYG{n}{std}\PYG{o}{:}\PYG{o}{:}\PYG{n}{recursive\PYGZus{}mutex}\PYG{o}{\PYGZgt{}}\PYG{+w}{ }\PYG{n}{lock}\PYG{p}{(}\PYG{n}{m}\PYG{p}{)}\PYG{p}{;}
\PYG{+w}{    }\PYG{k}{auto}\PYG{+w}{ }\PYG{n}{temp}\PYG{+w}{ }\PYG{o}{=}\PYG{+w}{ }\PYG{n}{data}\PYG{p}{;}
\PYG{+w}{    }\PYG{n}{update}\PYG{p}{(}\PYG{n}{d}\PYG{p}{)}\PYG{p}{;}
\PYG{+w}{    }\PYG{k}{return}\PYG{+w}{ }\PYG{n}{temp}\PYG{p}{;}
\PYG{+w}{  }\PYG{p}{\PYGZcb{}}
\PYG{p}{\PYGZcb{}}\PYG{p}{;}
\end{Verbatim}
\end{codeblockbg}


%%C080148 h3

\subsection{具体做法}


%%C080149 p

要把前面的实现改造成使用非递归互斥量的线程安全类型，请这样做：


%%C080150 p

把\myttbr{std::\allowbreak{}recursive\_\allowbreak{}mutex}替换为\texttt{std::\allowbreak{}mutex}：


%%C080151 code

\begin{codeblockbg}
\begin{Verbatim}[commandchars=\\\{\},breaklines=true,breakanywhere=true,fontsize=\codefont,xleftmargin=17bp]
\PYG{k}{class}\PYG{+w}{ }\PYG{n+nc}{foo}
\PYG{p}{\PYGZob{}}
\PYG{+w}{  }\PYG{n}{std}\PYG{o}{:}\PYG{o}{:}\PYG{n}{mutex}\PYG{+w}{ }\PYG{n}{m}\PYG{p}{;}
\PYG{+w}{  }\PYG{k+kt}{int}\PYG{+w}{        }\PYG{n}{data}\PYG{p}{;}
\PYG{+w}{  }\PYG{c+c1}{// 续第 2 步}
\PYG{p}{\PYGZcb{}}\PYG{p}{;}
\end{Verbatim}
\end{codeblockbg}


%%C080152 p

为公有方法定义私有的、非线程安全的版本或辅助函数


%%C080153 p

供线程安全的公有方法使用：


%%C080154 code

\begin{codeblockbg}
\begin{Verbatim}[commandchars=\\\{\},breaklines=true,breakanywhere=true,fontsize=\codefont,xleftmargin=17bp]
\PYG{k+kt}{void}\PYG{+w}{ }\PYG{n+nf}{internal\PYGZus{}update}\PYG{p}{(}\PYG{k+kt}{int}\PYG{+w}{ }\PYG{k}{const}\PYG{+w}{ }\PYG{n}{d}\PYG{p}{)}\PYG{+w}{ }\PYG{p}{\PYGZob{}}\PYG{+w}{ }\PYG{n}{data}\PYG{+w}{ }\PYG{o}{=}\PYG{+w}{ }\PYG{n}{d}\PYG{p}{;}\PYG{+w}{ }\PYG{p}{\PYGZcb{}}
\PYG{c+c1}{// 续第 3 步}
\end{Verbatim}
\end{codeblockbg}


%%C080155 p

改写公有方法，让它们使用刚刚定义的非线程安全的私有


%%C080156 p

方法：


%%C080157 code

\begin{codeblockbg}
\begin{Verbatim}[commandchars=\\\{\},breaklines=true,breakanywhere=true,fontsize=\codefont,xleftmargin=17bp]
\PYG{k}{public}\PYG{o}{:}
\PYG{+w}{  }\PYG{n}{foo}\PYG{p}{(}\PYG{k+kt}{int}\PYG{+w}{ }\PYG{k}{const}\PYG{+w}{ }\PYG{n}{d}\PYG{+w}{ }\PYG{o}{=}\PYG{+w}{ }\PYG{l+m+mi}{0}\PYG{p}{)}\PYG{+w}{ }\PYG{o}{:}\PYG{+w}{ }\PYG{n}{data}\PYG{p}{(}\PYG{n}{d}\PYG{p}{)}\PYG{+w}{ }\PYG{p}{\PYGZob{}}\PYG{p}{\PYGZcb{}}
\PYG{+w}{  }\PYG{k+kt}{void}\PYG{+w}{ }\PYG{n}{update}\PYG{p}{(}\PYG{k+kt}{int}\PYG{+w}{ }\PYG{k}{const}\PYG{+w}{ }\PYG{n}{d}\PYG{p}{)}
\PYG{+w}{  }\PYG{p}{\PYGZob{}}
\PYG{+w}{    }\PYG{n}{std}\PYG{o}{:}\PYG{o}{:}\PYG{n}{lock\PYGZus{}guard}\PYG{o}{\PYGZlt{}}\PYG{n}{std}\PYG{o}{:}\PYG{o}{:}\PYG{n}{mutex}\PYG{o}{\PYGZgt{}}\PYG{+w}{ }\PYG{n}{lock}\PYG{p}{(}\PYG{n}{m}\PYG{p}{)}\PYG{p}{;}
\PYG{+w}{    }\PYG{n}{internal\PYGZus{}update}\PYG{p}{(}\PYG{n}{d}\PYG{p}{)}\PYG{p}{;}
\PYG{+w}{  }\PYG{p}{\PYGZcb{}}
\PYG{k+kt}{int}\PYG{+w}{ }\PYG{n}{update\PYGZus{}with\PYGZus{}return}\PYG{p}{(}\PYG{k+kt}{int}\PYG{+w}{ }\PYG{k}{const}\PYG{+w}{ }\PYG{n}{d}\PYG{p}{)}
\PYG{p}{\PYGZob{}}
\PYG{+w}{  }\PYG{n}{std}\PYG{o}{:}\PYG{o}{:}\PYG{n}{lock\PYGZus{}guard}\PYG{o}{\PYGZlt{}}\PYG{n}{std}\PYG{o}{:}\PYG{o}{:}\PYG{n}{mutex}\PYG{o}{\PYGZgt{}}\PYG{+w}{ }\PYG{n}{lock}\PYG{p}{(}\PYG{n}{m}\PYG{p}{)}\PYG{p}{;}
\PYG{+w}{  }\PYG{k}{auto}\PYG{+w}{ }\PYG{n}{temp}\PYG{+w}{ }\PYG{o}{=}\PYG{+w}{ }\PYG{n}{data}\PYG{p}{;}
\PYG{+w}{  }\PYG{n}{internal\PYGZus{}update}\PYG{p}{(}\PYG{n}{d}\PYG{p}{)}\PYG{p}{;}
\PYG{+w}{  }\PYG{k}{return}\PYG{+w}{ }\PYG{n}{temp}\PYG{p}{;}
\PYG{p}{\PYGZcb{}}
\end{Verbatim}
\end{codeblockbg}


%%C080158 h3

\subsection{工作原理}


%%C080159 p

\myttbr{std::\allowbreak{}recursive\_\allowbreak{}mutex}类可以被同一线程多次锁定，无论是调用\texttt{lock()}还是\texttt{try\_\allowbreak{}lock()}。当某个线程锁定一个可用的递归互斥量时，它就获得了所有权；正因如此，同一线程后续的锁定尝试不会阻塞该线程的执行、造成死锁。不过，只有调用\texttt{unlock()}的次数与加锁次数相等时，递归互斥量才会被释放。


%%C080160 p

刚才讨论的\texttt{foo\_\allowbreak{}rec}类用递归互斥量保护对共享数据的访问；这里的共享数据是一个整型成员变量，由两个线程安全的公有函数访问：


%%C080161 li

\begin{itemize}
\item \texttt{update()}在私有变量中设置新值。


%%C080162 li

\item \myttbr{update\_\allowbreak{}and\_\allowbreak{}return()}在私有变量中设置新值，并把旧值返回给调用方。该函数调用\texttt{update()}来设置新值。

\end{itemize}


%%C080163 p

\texttt{foo\_\allowbreak{}rec}的实现或许本意是避免代码重复，但这种做法其实是一个设计错误，可以按\textit{具体做法}一节所示加以改进。与其复用公有的线程安全函数，不如提供私有的非线程安全函数，再由公有接口调用它们。


%%C080164 p

同样的方案也可用于其他类似问题：先定义代码的非线程安全版本，再提供（或许是轻量级的）线程安全包装。


%%C080165 h3

\subsection{延伸阅读}


%%C080166 li

\begin{itemize}
\item \textit{使用线程}


%%C080167 li

\item \textit{用互斥量与锁同步对共享数据的访问}

\end{itemize}


%%C080168 h2

\section{在线程之间发送通知}


%%C080169 p

互斥量是可用于保护共享数据访问的同步原语。不过标准库还提供了一种称为条件变量（condition variable）的同步原语，它让一个线程能够向其他线程发出信号，表示某个条件已经发生。等待在该条件变量上的线程会被阻塞，直到条件变量被通知，或者发生超时、虚假唤醒。本案例将介绍如何用条件变量在生产数据的线程与消费数据的线程之间发送通知。


%%C080170 h3

\subsection{准备工作}


%%C080171 p

本案例需要你熟悉线程、互斥量与锁。条件变量位于\myttbr{<\allowbreak{}condition\_\allowbreak{}variable>}头文件的\texttt{std}命名空间中。


%%C080172 h3

\subsection{具体做法}


%%C080173 p

在条件变量上通过通知来同步线程，请使用下面的模式：


%%C080174 p

定义条件变量（放在合适的上下文中）：


%%C080175 code

\begin{codeblockbg}
\begin{Verbatim}[commandchars=\\\{\},breaklines=true,breakanywhere=true,fontsize=\codefont,xleftmargin=17bp]
\PYG{n}{std}\PYG{o}{:}\PYG{o}{:}\PYG{n}{condition\PYGZus{}variable}\PYG{+w}{ }\PYG{n}{cv}\PYG{p}{;}
\end{Verbatim}
\end{codeblockbg}


%%C080176 p

定义一个供线程加锁用的互斥量：


%%C080177 code

\begin{codeblockbg}
\begin{Verbatim}[commandchars=\\\{\},breaklines=true,breakanywhere=true,fontsize=\codefont,xleftmargin=17bp]
\PYG{n}{std}\PYG{o}{:}\PYG{o}{:}\PYG{n}{mutex}\PYG{+w}{ }\PYG{n}{cv\PYGZus{}mutex}\PYG{p}{;}
\end{Verbatim}
\end{codeblockbg}


%%C080178 p

定义线程之间使用的共享数据：


%%C080179 code

\begin{codeblockbg}
\begin{Verbatim}[commandchars=\\\{\},breaklines=true,breakanywhere=true,fontsize=\codefont,xleftmargin=17bp]
\PYG{k+kt}{int}\PYG{+w}{ }\PYG{n}{data}\PYG{+w}{ }\PYG{o}{=}\PYG{+w}{ }\PYG{l+m+mi}{0}\PYG{p}{;}
\end{Verbatim}
\end{codeblockbg}


%%C080180 p

在生产线程中，修改数据之前先锁定互斥量：


%%C080181 code

\begin{codeblockbg}
\begin{Verbatim}[commandchars=\\\{\},breaklines=true,breakanywhere=true,fontsize=\codefont,xleftmargin=17bp]
\PYG{n}{std}\PYG{o}{:}\PYG{o}{:}\PYG{k+kr}{thread}\PYG{+w}{ }\PYG{n+nf}{p}\PYG{p}{(}\PYG{p}{[}\PYG{o}{\PYGZam{}}\PYG{p}{]}\PYG{p}{(}\PYG{p}{)}\PYG{p}{\PYGZob{}}
\PYG{+w}{  }\PYG{c+c1}{// 模拟耗时较长的操作}
\PYG{+w}{  }\PYG{p}{\PYGZob{}}
\PYG{+w}{    }\PYG{k}{using}\PYG{+w}{ }\PYG{k}{namespace}\PYG{+w}{ }\PYG{n+nn}{std}\PYG{o}{:}\PYG{o}{:}\PYG{n+nn}{chrono\PYGZus{}literals}\PYG{p}{;}
\PYG{+w}{    }\PYG{n}{std}\PYG{o}{:}\PYG{o}{:}\PYG{n}{this\PYGZus{}thread}\PYG{o}{:}\PYG{o}{:}\PYG{n}{sleep\PYGZus{}for}\PYG{p}{(}\PYG{l+m+mi}{2}\PYG{n}{s}\PYG{p}{)}\PYG{p}{;}
\PYG{+w}{  }\PYG{p}{\PYGZcb{}}
\PYG{c+c1}{// 生产数据}
\PYG{p}{\PYGZob{}}
\PYG{+w}{  }\PYG{n}{std}\PYG{o}{:}\PYG{o}{:}\PYG{n}{unique\PYGZus{}lock}\PYG{+w}{ }\PYG{n+nf}{lock}\PYG{p}{(}\PYG{n}{cv\PYGZus{}mutex}\PYG{p}{)}\PYG{p}{;}
\PYG{+w}{  }\PYG{n}{data}\PYG{+w}{ }\PYG{o}{=}\PYG{+w}{ }\PYG{l+m+mi}{42}\PYG{p}{;}
\PYG{p}{\PYGZcb{}}
\PYG{c+c1}{// 打印消息}
\PYG{p}{\PYGZob{}}
\PYG{+w}{  }\PYG{n}{std}\PYG{o}{:}\PYG{o}{:}\PYG{n}{lock\PYGZus{}guard}\PYG{+w}{ }\PYG{n+nf}{l}\PYG{p}{(}\PYG{n}{io\PYGZus{}mutex}\PYG{p}{)}\PYG{p}{;}
\PYG{+w}{  }\PYG{n}{std}\PYG{o}{:}\PYG{o}{:}\PYG{n}{cout}\PYG{+w}{ }\PYG{o}{\PYGZlt{}}\PYG{o}{\PYGZlt{}}\PYG{+w}{ }\PYG{l+s}{\PYGZdq{}}\PYG{l+s}{produced }\PYG{l+s}{\PYGZdq{}}\PYG{+w}{ }\PYG{o}{\PYGZlt{}}\PYG{o}{\PYGZlt{}}\PYG{+w}{ }\PYG{n}{data}\PYG{+w}{ }\PYG{o}{\PYGZlt{}}\PYG{o}{\PYGZlt{}}\PYG{+w}{ }\PYG{n}{std}\PYG{o}{:}\PYG{o}{:}\PYG{n}{endl}\PYG{p}{;}
\PYG{p}{\PYGZcb{}}
\PYG{+w}{  }\PYG{c+c1}{// 续第 5 步}
\PYG{p}{\PYGZcb{}}\PYG{p}{)}\PYG{p}{;}
\end{Verbatim}
\end{codeblockbg}


%%C080182 p

在生产线程中，调用


%%C080183 p

\texttt{notify\_\allowbreak{}one()}或\texttt{notify\_\allowbreak{}all()}通知条件变量（应在保护共享数据的互斥量解锁之后进行）：


%%C080184 code

\begin{codeblockbg}
\begin{Verbatim}[commandchars=\\\{\},breaklines=true,breakanywhere=true,fontsize=\codefont,xleftmargin=17bp]
\PYG{c+c1}{// 接第 4 步}
\PYG{n}{cv}\PYG{p}{.}\PYG{n}{notify\PYGZus{}one}\PYG{p}{(}\PYG{p}{)}\PYG{p}{;}
\end{Verbatim}
\end{codeblockbg}


%%C080185 p

在消费线程中，先获取互斥量的唯一锁，再用它等待


%%C080186 p

条件变量：


%%C080187 code

\begin{codeblockbg}
\begin{Verbatim}[commandchars=\\\{\},breaklines=true,breakanywhere=true,fontsize=\codefont,xleftmargin=17bp]
\PYG{n}{std}\PYG{o}{:}\PYG{o}{:}\PYG{k+kr}{thread}\PYG{+w}{ }\PYG{n+nf}{c}\PYG{p}{(}\PYG{p}{[}\PYG{o}{\PYGZam{}}\PYG{p}{]}\PYG{p}{(}\PYG{p}{)}\PYG{p}{\PYGZob{}}
\PYG{+w}{  }\PYG{c+c1}{// 等待通知}
\PYG{+w}{  }\PYG{p}{\PYGZob{}}
\PYG{+w}{    }\PYG{n}{std}\PYG{o}{:}\PYG{o}{:}\PYG{n}{unique\PYGZus{}lock}\PYG{+w}{ }\PYG{n}{lock}\PYG{p}{(}\PYG{n}{cv\PYGZus{}mutex}\PYG{p}{)}\PYG{p}{;}
\PYG{+w}{    }\PYG{n}{cv}\PYG{p}{.}\PYG{n}{wait}\PYG{p}{(}\PYG{n}{lock}\PYG{p}{)}\PYG{p}{;}
\PYG{+w}{  }\PYG{p}{\PYGZcb{}}
\PYG{+w}{  }\PYG{c+c1}{// 续第 7 步}
\PYG{p}{\PYGZcb{}}\PYG{p}{)}\PYG{p}{;}
\end{Verbatim}
\end{codeblockbg}


%%C080188 p

在消费线程中，条件被通知之后使用共享数据：


%%C080189 code

\begin{codeblockbg}
\begin{Verbatim}[commandchars=\\\{\},breaklines=true,breakanywhere=true,fontsize=\codefont,xleftmargin=17bp]
\PYG{c+c1}{// 接第 6 步}
\PYG{p}{\PYGZob{}}
\PYG{+w}{  }\PYG{n}{std}\PYG{o}{:}\PYG{o}{:}\PYG{n}{lock\PYGZus{}guard}\PYG{+w}{ }\PYG{n+nf}{lock}\PYG{p}{(}\PYG{n}{io\PYGZus{}mutex}\PYG{p}{)}\PYG{p}{;}
\PYG{+w}{  }\PYG{n}{std}\PYG{o}{:}\PYG{o}{:}\PYG{n}{cout}\PYG{+w}{ }\PYG{o}{\PYGZlt{}}\PYG{o}{\PYGZlt{}}\PYG{+w}{ }\PYG{l+s}{\PYGZdq{}}\PYG{l+s}{consumed }\PYG{l+s}{\PYGZdq{}}\PYG{+w}{ }\PYG{o}{\PYGZlt{}}\PYG{o}{\PYGZlt{}}\PYG{+w}{ }\PYG{n}{data}\PYG{+w}{ }\PYG{o}{\PYGZlt{}}\PYG{o}{\PYGZlt{}}\PYG{+w}{ }\PYG{n}{std}\PYG{o}{:}\PYG{o}{:}\PYG{n}{endl}\PYG{p}{;}
\PYG{p}{\PYGZcb{}}
\end{Verbatim}
\end{codeblockbg}


%%C080190 h3

\subsection{工作原理}


%%C080191 p

前面的示例中有两个线程共享同一份数据（这里是一个整型变量）。一个线程在耗时较长的计算之后生产数据（用休眠来模拟），另一个线程只在数据生产出来之后才消费它。为此，它们使用由互斥量和条件变量构成的同步机制，后者会把消费线程阻塞，直到生产线程发出数据已可用的通知。这条通信通道的关键是条件变量：消费线程一直等待它，直到生产线程通知它。两个线程几乎同时启动。生产线程开始一段长计算，这段计算应当为消费线程产出数据。与此同时，在数据可用之前消费线程实际上无法继续，它必须保持阻塞，直到收到数据已生产完毕的通知。一旦收到通知，它就可以继续执行。整个机制的工作方式如下：


%%C080192 li

\begin{itemize}
\item 必须至少有一个线程在等待条件变量被通知。


%%C080193 li

\item 必须至少有一个线程在向条件变量发信号。


%%C080194 li

\item 等待的线程必须先在某个互斥量上加锁（\myttbr{std::\allowbreak{}unique\_\allowbreak{}lock<\allowbreak{}std::\allowbreak{}mutex>}），并把该锁传给条件变量的\texttt{wait()}、\texttt{wait\_\allowbreak{}for()}或\texttt{wait\_\allowbreak{}until()}方法。所有这些等待方法都会原子地释放互斥量并阻塞线程，直到条件变量被通知。此时线程解除阻塞，并原子地重新获取互斥量。


%%C080195 li

\item 发出信号的线程可以用\texttt{notify\_\allowbreak{}one()}（解除一个被阻塞线程的阻塞），也可以用\texttt{notify\_\allowbreak{}all()}（解除所有等待该条件变量的被阻塞线程的阻塞）。

\end{itemize}


%%C080196 note

\begin{myNote}
在多处理器系统上，条件变量无法做到完全可预测。因此可能出现\textit{虚假唤醒}（spurious wakeup），即使没有任何线程通知条件变量，某个线程也会被解除阻塞。所以，线程解除阻塞之后必须检查条件是否真的成立。此外，虚假唤醒可能多次发生，因此需要在一个循环中检查条件变量。
\end{myNote}


%%C080197 p

C++标准提供了两种条件变量的实现：


%%C080198 li

\begin{itemize}
\item \myttbr{std::\allowbreak{}condition\_\allowbreak{}variable}，本案例使用的就是它，它定义的条件变量与\texttt{std::\allowbreak{}unique\_\allowbreak{}lock}配合使用。


%%C080199 li

\item \myttbr{std::\allowbreak{}condition\_\allowbreak{}variable\_\allowbreak{}any}是一个更通用的实现，可与任何满足基本锁要求（实现了\texttt{lock()}与\texttt{unlock()}方法）的锁配合使用。这种实现的一个可能用途是提供可中断的等待，正如Anthony Williams在\textit{C++ concurrency in action}（2012）中所解释的：

\end{itemize}


%%C080200 p

\textit{自定义的锁操作既会按预期锁定关联的互斥量，也会在收到中断信号时完成通知该条件变量这一必要工作。}


%%C080201 p

条件变量的所有等待方法都有两个重载：

%%C080202 li

\begin{itemize}
\item 第一个重载接受\myttbr{std::\allowbreak{}unique\_\allowbreak{}lock<\allowbreak{}std::\allowbreak{}mutex>}（取决于类型，即时长或时间点），并使线程保持阻塞，直到条件变量被通知为止。该重载会原子地释放互斥量、阻塞当前线程，并把它加入等待该条件变量的线程列表。当条件通过\texttt{notify\_\allowbreak{}one()}或\texttt{notify\_\allowbreak{}all()}被通知、发生虚假唤醒或发生超时（取决于具体重载）时，线程解除阻塞。此时互斥量会被再次原子地获取。


%%C080203 li

\item 第二个重载在其他重载的参数之外还接受一个谓词。这个谓词可用于在等待条件变为\texttt{true}时避免虚假唤醒。该重载等价于下面的代码：


%%C080204 code

\begin{codeblockbg}
\begin{Verbatim}[commandchars=\\\{\},breaklines=true,breakanywhere=true,fontsize=\codefont,xleftmargin=17bp]
\PYG{k}{while}\PYG{p}{(}\PYG{o}{!}\PYG{n}{pred}\PYG{p}{(}\PYG{p}{)}\PYG{p}{)}
\PYG{+w}{  }\PYG{n}{wait}\PYG{p}{(}\PYG{n}{lock}\PYG{p}{)}\PYG{p}{;}
\end{Verbatim}
\end{codeblockbg}

\end{itemize}


%%C080205 p

下面的代码给出了一个与上一节所给示例类似但更复杂的例子。生产者线程在循环中生成数据（在本例中是有限循环），消费者线程等待新数据变为可用并消费它（打印到控制台）。生产者线程在完成数据生成后终止，消费者线程在不再有数据可消费时终止。数据被添加到\texttt{queue<\allowbreak{}int>}中，并用一个布尔变量向消费者线程表明数据生成过程已经结束：


%%C080206 code

\begin{codeblockbg}
\begin{Verbatim}[commandchars=\\\{\},breaklines=true,breakanywhere=true,fontsize=\codefont,xleftmargin=17bp]
\PYG{n}{std}\PYG{o}{:}\PYG{o}{:}\PYG{n}{mutex}\PYG{+w}{ }\PYG{n}{g\PYGZus{}lockprint}\PYG{p}{;}
\PYG{n}{std}\PYG{o}{:}\PYG{o}{:}\PYG{n}{mutex}\PYG{+w}{ }\PYG{n}{g\PYGZus{}lockqueue}\PYG{p}{;}
\PYG{n}{std}\PYG{o}{:}\PYG{o}{:}\PYG{n}{condition\PYGZus{}variable}\PYG{+w}{ }\PYG{n}{g\PYGZus{}queuecheck}\PYG{p}{;}
\PYG{n}{std}\PYG{o}{:}\PYG{o}{:}\PYG{n}{queue}\PYG{o}{\PYGZlt{}}\PYG{k+kt}{int}\PYG{o}{\PYGZgt{}}\PYG{+w}{ }\PYG{n}{g\PYGZus{}buffer}\PYG{p}{;}
\PYG{k+kt}{bool}\PYG{+w}{ }\PYG{n}{g\PYGZus{}done}\PYG{p}{;}
\PYG{k+kt}{void}\PYG{+w}{ }\PYG{n+nf}{producer}\PYG{p}{(}
\PYG{+w}{ }\PYG{k+kt}{int}\PYG{+w}{ }\PYG{k}{const}\PYG{+w}{ }\PYG{n}{id}\PYG{p}{,}
\PYG{+w}{ }\PYG{n}{std}\PYG{o}{:}\PYG{o}{:}\PYG{n}{mt19937}\PYG{o}{\PYGZam{}}\PYG{+w}{ }\PYG{n}{generator}\PYG{p}{,}
\PYG{+w}{ }\PYG{n}{std}\PYG{o}{:}\PYG{o}{:}\PYG{n}{uniform\PYGZus{}int\PYGZus{}distribution}\PYG{o}{\PYGZlt{}}\PYG{k+kt}{int}\PYG{o}{\PYGZgt{}}\PYG{o}{\PYGZam{}}\PYG{+w}{ }\PYG{n}{dsleep}\PYG{p}{,}
\PYG{+w}{ }\PYG{n}{std}\PYG{o}{:}\PYG{o}{:}\PYG{n}{uniform\PYGZus{}int\PYGZus{}distribution}\PYG{o}{\PYGZlt{}}\PYG{k+kt}{int}\PYG{o}{\PYGZgt{}}\PYG{o}{\PYGZam{}}\PYG{+w}{ }\PYG{n}{dcode}\PYG{p}{)}
\PYG{p}{\PYGZob{}}
\PYG{+w}{  }\PYG{k}{for}\PYG{+w}{ }\PYG{p}{(}\PYG{k+kt}{int}\PYG{+w}{ }\PYG{n}{i}\PYG{+w}{ }\PYG{o}{=}\PYG{+w}{ }\PYG{l+m+mi}{0}\PYG{p}{;}\PYG{+w}{ }\PYG{n}{i}\PYG{+w}{ }\PYG{o}{\PYGZlt{}}\PYG{+w}{ }\PYG{l+m+mi}{5}\PYG{p}{;}\PYG{+w}{ }\PYG{o}{+}\PYG{o}{+}\PYG{n}{i}\PYG{p}{)}
\PYG{+w}{  }\PYG{p}{\PYGZob{}}
\PYG{+w}{    }\PYG{c+c1}{// 模拟工作}
\PYG{+w}{    }\PYG{n}{std}\PYG{o}{:}\PYG{o}{:}\PYG{n}{this\PYGZus{}thread}\PYG{o}{:}\PYG{o}{:}\PYG{n}{sleep\PYGZus{}for}\PYG{p}{(}
\PYG{+w}{      }\PYG{n}{std}\PYG{o}{:}\PYG{o}{:}\PYG{n}{chrono}\PYG{o}{:}\PYG{o}{:}\PYG{n}{milliseconds}\PYG{p}{(}\PYG{n}{dsleep}\PYG{p}{(}\PYG{n}{generator}\PYG{p}{)}\PYG{p}{)}\PYG{p}{)}\PYG{p}{;}
\PYG{c+c1}{// 生成数据}
\PYG{p}{\PYGZob{}}
\PYG{+w}{  }\PYG{n}{std}\PYG{o}{:}\PYG{o}{:}\PYG{n}{unique\PYGZus{}lock}\PYG{o}{\PYGZlt{}}\PYG{n}{std}\PYG{o}{:}\PYG{o}{:}\PYG{n}{mutex}\PYG{o}{\PYGZgt{}}\PYG{+w}{ }\PYG{n}{locker}\PYG{p}{(}\PYG{n}{g\PYGZus{}lockqueue}\PYG{p}{)}\PYG{p}{;}
\PYG{+w}{  }\PYG{k+kt}{int}\PYG{+w}{ }\PYG{n}{value}\PYG{+w}{ }\PYG{o}{=}\PYG{+w}{ }\PYG{n}{id}\PYG{+w}{ }\PYG{o}{*}\PYG{+w}{ }\PYG{l+m+mi}{100}\PYG{+w}{ }\PYG{o}{+}\PYG{+w}{ }\PYG{n}{dcode}\PYG{p}{(}\PYG{n}{generator}\PYG{p}{)}\PYG{p}{;}
\PYG{+w}{  }\PYG{n}{g\PYGZus{}buffer}\PYG{p}{.}\PYG{n}{push}\PYG{p}{(}\PYG{n}{value}\PYG{p}{)}\PYG{p}{;}
\PYG{+w}{  }\PYG{p}{\PYGZob{}}
\PYG{+w}{    }\PYG{n}{std}\PYG{o}{:}\PYG{o}{:}\PYG{n}{unique\PYGZus{}lock}\PYG{o}{\PYGZlt{}}\PYG{n}{std}\PYG{o}{:}\PYG{o}{:}\PYG{n}{mutex}\PYG{o}{\PYGZgt{}}\PYG{+w}{ }\PYG{n}{locker}\PYG{p}{(}\PYG{n}{g\PYGZus{}lockprint}\PYG{p}{)}\PYG{p}{;}
\PYG{+w}{    }\PYG{n}{std}\PYG{o}{:}\PYG{o}{:}\PYG{n}{cout}\PYG{+w}{ }\PYG{o}{\PYGZlt{}}\PYG{o}{\PYGZlt{}}\PYG{+w}{ }\PYG{l+s}{\PYGZdq{}}\PYG{l+s}{[produced(}\PYG{l+s}{\PYGZdq{}}\PYG{+w}{ }\PYG{o}{\PYGZlt{}}\PYG{o}{\PYGZlt{}}\PYG{+w}{ }\PYG{n}{id}\PYG{+w}{ }\PYG{o}{\PYGZlt{}}\PYG{o}{\PYGZlt{}}\PYG{+w}{ }\PYG{l+s}{\PYGZdq{}}\PYG{l+s}{)]: }\PYG{l+s}{\PYGZdq{}}\PYG{+w}{ }\PYG{o}{\PYGZlt{}}\PYG{o}{\PYGZlt{}}\PYG{+w}{ }\PYG{n}{value}
\PYG{+w}{              }\PYG{o}{\PYGZlt{}}\PYG{o}{\PYGZlt{}}\PYG{+w}{ }\PYG{n}{std}\PYG{o}{:}\PYG{o}{:}\PYG{n}{endl}\PYG{p}{;}
\PYG{+w}{  }\PYG{p}{\PYGZcb{}}
\PYG{p}{\PYGZcb{}}
\PYG{+w}{    }\PYG{c+c1}{// 通知消费者}
\PYG{+w}{    }\PYG{n}{g\PYGZus{}queuecheck}\PYG{p}{.}\PYG{n}{notify\PYGZus{}one}\PYG{p}{(}\PYG{p}{)}\PYG{p}{;}
\PYG{+w}{  }\PYG{p}{\PYGZcb{}}
\PYG{p}{\PYGZcb{}}
\PYG{k+kt}{void}\PYG{+w}{ }\PYG{n+nf}{consumer}\PYG{p}{(}\PYG{p}{)}
\PYG{p}{\PYGZob{}}
\PYG{+w}{  }\PYG{c+c1}{// 循环直到收到结束信号}
\PYG{+w}{  }\PYG{k}{while}\PYG{+w}{ }\PYG{p}{(}\PYG{o}{!}\PYG{n}{g\PYGZus{}done}\PYG{p}{)}
\PYG{+w}{  }\PYG{p}{\PYGZob{}}
\PYG{+w}{    }\PYG{n}{std}\PYG{o}{:}\PYG{o}{:}\PYG{n}{unique\PYGZus{}lock}\PYG{o}{\PYGZlt{}}\PYG{n}{std}\PYG{o}{:}\PYG{o}{:}\PYG{n}{mutex}\PYG{o}{\PYGZgt{}}\PYG{+w}{ }\PYG{n}{locker}\PYG{p}{(}\PYG{n}{g\PYGZus{}lockqueue}\PYG{p}{)}\PYG{p}{;}
\PYG{n}{g\PYGZus{}queuecheck}\PYG{p}{.}\PYG{n}{wait\PYGZus{}for}\PYG{p}{(}
\PYG{+w}{  }\PYG{n}{locker}\PYG{p}{,}
\PYG{+w}{  }\PYG{n}{std}\PYG{o}{:}\PYG{o}{:}\PYG{n}{chrono}\PYG{o}{:}\PYG{o}{:}\PYG{n}{seconds}\PYG{p}{(}\PYG{l+m+mi}{1}\PYG{p}{)}\PYG{p}{,}
\PYG{+w}{  }\PYG{p}{[}\PYG{o}{\PYGZam{}}\PYG{p}{]}\PYG{p}{(}\PYG{p}{)}\PYG{+w}{ }\PYG{p}{\PYGZob{}}\PYG{k}{return}\PYG{+w}{ }\PYG{o}{!}\PYG{n}{g\PYGZus{}buffer}\PYG{p}{.}\PYG{n}{empty}\PYG{p}{(}\PYG{p}{)}\PYG{p}{;}\PYG{+w}{ }\PYG{p}{\PYGZcb{}}\PYG{p}{)}\PYG{p}{;}
\PYG{c+c1}{// 如果队列中有值就处理它们}
\PYG{k}{while}\PYG{+w}{ }\PYG{p}{(}\PYG{o}{!}\PYG{n}{g\PYGZus{}done}\PYG{+w}{ }\PYG{o}{\PYGZam{}}\PYG{o}{\PYGZam{}}\PYG{+w}{ }\PYG{o}{!}\PYG{n}{g\PYGZus{}buffer}\PYG{p}{.}\PYG{n}{empty}\PYG{p}{(}\PYG{p}{)}\PYG{p}{)}
\PYG{p}{\PYGZob{}}
\PYG{+w}{  }\PYG{n}{std}\PYG{o}{:}\PYG{o}{:}\PYG{n}{unique\PYGZus{}lock}\PYG{o}{\PYGZlt{}}\PYG{n}{std}\PYG{o}{:}\PYG{o}{:}\PYG{n}{mutex}\PYG{o}{\PYGZgt{}}\PYG{+w}{ }\PYG{n}{locker}\PYG{p}{(}\PYG{n}{g\PYGZus{}lockprint}\PYG{p}{)}\PYG{p}{;}
\PYG{+w}{  }\PYG{n}{std}\PYG{o}{:}\PYG{o}{:}\PYG{n}{cout}
\PYG{+w}{    }\PYG{o}{\PYGZlt{}}\PYG{o}{\PYGZlt{}}\PYG{+w}{ }\PYG{l+s}{\PYGZdq{}}\PYG{l+s}{[consumed]: }\PYG{l+s}{\PYGZdq{}}\PYG{+w}{ }\PYG{o}{\PYGZlt{}}\PYG{o}{\PYGZlt{}}\PYG{+w}{ }\PYG{n}{g\PYGZus{}buffer}\PYG{p}{.}\PYG{n}{front}\PYG{p}{(}\PYG{p}{)}
\PYG{+w}{    }\PYG{o}{\PYGZlt{}}\PYG{o}{\PYGZlt{}}\PYG{+w}{ }\PYG{n}{std}\PYG{o}{:}\PYG{o}{:}\PYG{n}{endl}\PYG{p}{;}
\PYG{+w}{  }\PYG{n}{g\PYGZus{}buffer}\PYG{p}{.}\PYG{n}{pop}\PYG{p}{(}\PYG{p}{)}\PYG{p}{;}
\PYG{p}{\PYGZcb{}}
\PYG{+w}{  }\PYG{p}{\PYGZcb{}}
\PYG{p}{\PYGZcb{}}
\end{Verbatim}
\end{codeblockbg}


%%C080207 p

消费者线程执行以下操作：


%%C080208 li

\begin{itemize}
\item 循环，直到收到数据生成过程已经结束的信号。


%%C080209 li

\item 获取与条件变量关联的\texttt{mutex}上的唯一锁。


%%C080210 li

\item 使用接受谓词的\texttt{wait\_\allowbreak{}for()}重载，在发生唤醒时检查缓冲区非空（以避免虚假唤醒）。该方法使用1秒的超时，即使条件被通知，也会在超时后返回。


%%C080211 li

\item 通过条件变量收到通知后，消费队列中的全部数据。

\end{itemize}


%%C080212 p

为了测试这一点，我们启动若干生产者线程和一个消费者线程。生产者线程生成随机数据，因此共享伪随机数生成器引擎和分布。所有这些都在下面的代码示例中展示：


%%C080213 code

\begin{codeblockbg}
\begin{Verbatim}[commandchars=\\\{\},breaklines=true,breakanywhere=true,fontsize=\codefont,xleftmargin=17bp]
\PYG{k}{auto}\PYG{+w}{ }\PYG{n}{seed\PYGZus{}data}\PYG{+w}{ }\PYG{o}{=}\PYG{+w}{ }\PYG{n}{std}\PYG{o}{:}\PYG{o}{:}\PYG{n}{array}\PYG{o}{\PYGZlt{}}\PYG{k+kt}{int}\PYG{p}{,}\PYG{+w}{ }\PYG{n}{std}\PYG{o}{:}\PYG{o}{:}\PYG{n}{mt19937}\PYG{o}{:}\PYG{o}{:}\PYG{n}{state\PYGZus{}size}\PYG{o}{\PYGZgt{}}\PYG{+w}{ }\PYG{p}{\PYGZob{}}\PYG{p}{\PYGZcb{}}\PYG{p}{;}
\PYG{n}{std}\PYG{o}{:}\PYG{o}{:}\PYG{n}{random\PYGZus{}device}\PYG{+w}{ }\PYG{n}{rd}\PYG{+w}{ }\PYG{p}{\PYGZob{}}\PYG{p}{\PYGZcb{}}\PYG{p}{;}
\PYG{n}{std}\PYG{o}{:}\PYG{o}{:}\PYG{n}{generate}\PYG{p}{(}\PYG{n}{std}\PYG{o}{:}\PYG{o}{:}\PYG{n}{begin}\PYG{p}{(}\PYG{n}{seed\PYGZus{}data}\PYG{p}{)}\PYG{p}{,}\PYG{+w}{ }\PYG{n}{std}\PYG{o}{:}\PYG{o}{:}\PYG{n}{end}\PYG{p}{(}\PYG{n}{seed\PYGZus{}data}\PYG{p}{)}\PYG{p}{,}
\PYG{+w}{              }\PYG{n}{std}\PYG{o}{:}\PYG{o}{:}\PYG{n}{ref}\PYG{p}{(}\PYG{n}{rd}\PYG{p}{)}\PYG{p}{)}\PYG{p}{;}
\PYG{n}{std}\PYG{o}{:}\PYG{o}{:}\PYG{n}{seed\PYGZus{}seq}\PYG{+w}{ }\PYG{n+nf}{seq}\PYG{p}{(}\PYG{n}{std}\PYG{o}{:}\PYG{o}{:}\PYG{n}{begin}\PYG{p}{(}\PYG{n}{seed\PYGZus{}data}\PYG{p}{)}\PYG{p}{,}\PYG{+w}{ }\PYG{n}{std}\PYG{o}{:}\PYG{o}{:}\PYG{n}{end}\PYG{p}{(}\PYG{n}{seed\PYGZus{}data}\PYG{p}{)}\PYG{p}{)}\PYG{p}{;}
\PYG{k}{auto}\PYG{+w}{ }\PYG{n}{generator}\PYG{+w}{ }\PYG{o}{=}\PYG{+w}{ }\PYG{n}{std}\PYG{o}{:}\PYG{o}{:}\PYG{n}{mt19937}\PYG{p}{\PYGZob{}}\PYG{+w}{ }\PYG{n}{seq}\PYG{+w}{ }\PYG{p}{\PYGZcb{}}\PYG{p}{;}
\PYG{k}{auto}\PYG{+w}{ }\PYG{n}{dsleep}\PYG{+w}{ }\PYG{o}{=}\PYG{+w}{ }\PYG{n}{std}\PYG{o}{:}\PYG{o}{:}\PYG{n}{uniform\PYGZus{}int\PYGZus{}distribution}\PYG{o}{\PYGZlt{}}\PYG{o}{\PYGZgt{}}\PYG{p}{\PYGZob{}}\PYG{+w}{ }\PYG{l+m+mi}{100}\PYG{p}{,}\PYG{+w}{ }\PYG{l+m+mi}{500}\PYG{+w}{ }\PYG{p}{\PYGZcb{}}\PYG{p}{;}
\PYG{k}{auto}\PYG{+w}{ }\PYG{n}{dcode}\PYG{+w}{ }\PYG{o}{=}\PYG{+w}{ }\PYG{n}{std}\PYG{o}{:}\PYG{o}{:}\PYG{n}{uniform\PYGZus{}int\PYGZus{}distribution}\PYG{o}{\PYGZlt{}}\PYG{o}{\PYGZgt{}}\PYG{p}{\PYGZob{}}\PYG{+w}{ }\PYG{l+m+mi}{1}\PYG{p}{,}\PYG{+w}{ }\PYG{l+m+mi}{99}\PYG{+w}{ }\PYG{p}{\PYGZcb{}}\PYG{p}{;}
\PYG{n}{std}\PYG{o}{:}\PYG{o}{:}\PYG{n}{cout}\PYG{+w}{ }\PYG{o}{\PYGZlt{}}\PYG{o}{\PYGZlt{}}\PYG{+w}{ }\PYG{l+s}{\PYGZdq{}}\PYG{l+s}{start producing and consuming...}\PYG{l+s}{\PYGZdq{}}\PYG{+w}{ }\PYG{o}{\PYGZlt{}}\PYG{o}{\PYGZlt{}}\PYG{+w}{ }\PYG{n}{std}\PYG{o}{:}\PYG{o}{:}\PYG{n}{endl}\PYG{p}{;}
\PYG{n}{std}\PYG{o}{:}\PYG{o}{:}\PYG{k+kr}{thread}\PYG{+w}{ }\PYG{n+nf}{consumerthread}\PYG{p}{(}\PYG{n}{consumer}\PYG{p}{)}\PYG{p}{;}
\PYG{n}{std}\PYG{o}{:}\PYG{o}{:}\PYG{n}{vector}\PYG{o}{\PYGZlt{}}\PYG{n}{std}\PYG{o}{:}\PYG{o}{:}\PYG{k+kr}{thread}\PYG{o}{\PYGZgt{}}\PYG{+w}{ }\PYG{n}{threads}\PYG{p}{;}
\PYG{k}{for}\PYG{+w}{ }\PYG{p}{(}\PYG{k+kt}{int}\PYG{+w}{ }\PYG{n}{i}\PYG{+w}{ }\PYG{o}{=}\PYG{+w}{ }\PYG{l+m+mi}{0}\PYG{p}{;}\PYG{+w}{ }\PYG{n}{i}\PYG{+w}{ }\PYG{o}{\PYGZlt{}}\PYG{+w}{ }\PYG{l+m+mi}{5}\PYG{p}{;}\PYG{+w}{ }\PYG{o}{+}\PYG{o}{+}\PYG{n}{i}\PYG{p}{)}
\PYG{p}{\PYGZob{}}
\PYG{+w}{  }\PYG{n}{threads}\PYG{p}{.}\PYG{n}{emplace\PYGZus{}back}\PYG{p}{(}\PYG{n}{producer}\PYG{p}{,}
\PYG{+w}{                       }\PYG{n}{i}\PYG{+w}{ }\PYG{o}{+}\PYG{+w}{ }\PYG{l+m+mi}{1}\PYG{p}{,}
\PYG{+w}{                       }\PYG{n}{std}\PYG{o}{:}\PYG{o}{:}\PYG{n}{ref}\PYG{p}{(}\PYG{n}{generator}\PYG{p}{)}\PYG{p}{,}
\PYG{+w}{                       }\PYG{n}{std}\PYG{o}{:}\PYG{o}{:}\PYG{n}{ref}\PYG{p}{(}\PYG{n}{dsleep}\PYG{p}{)}\PYG{p}{,}
\PYG{+w}{                       }\PYG{n}{std}\PYG{o}{:}\PYG{o}{:}\PYG{n}{ref}\PYG{p}{(}\PYG{n}{dcode}\PYG{p}{)}\PYG{p}{)}\PYG{p}{;}
\PYG{p}{\PYGZcb{}}
\PYG{c+c1}{// 等待工作线程完成}
\PYG{k}{for}\PYG{+w}{ }\PYG{p}{(}\PYG{k}{auto}\PYG{o}{\PYGZam{}}\PYG{+w}{ }\PYG{n}{t}\PYG{+w}{ }\PYG{o}{:}\PYG{+w}{ }\PYG{n}{threads}\PYG{p}{)}
\PYG{+w}{  }\PYG{n}{t}\PYG{p}{.}\PYG{n}{join}\PYG{p}{(}\PYG{p}{)}\PYG{p}{;}
\PYG{c+c1}{// 通知记录器结束并等待它}
\PYG{n}{g\PYGZus{}done}\PYG{+w}{ }\PYG{o}{=}\PYG{+w}{ }\PYG{n+nb}{true}\PYG{p}{;}
\PYG{n}{consumerthread}\PYG{p}{.}\PYG{n}{join}\PYG{p}{(}\PYG{p}{)}\PYG{p}{;}
\PYG{n}{std}\PYG{o}{:}\PYG{o}{:}\PYG{n}{cout}\PYG{+w}{ }\PYG{o}{\PYGZlt{}}\PYG{o}{\PYGZlt{}}\PYG{+w}{ }\PYG{l+s}{\PYGZdq{}}\PYG{l+s}{done producing and consuming}\PYG{l+s}{\PYGZdq{}}\PYG{+w}{ }\PYG{o}{\PYGZlt{}}\PYG{o}{\PYGZlt{}}\PYG{+w}{ }\PYG{n}{std}\PYG{o}{:}\PYG{o}{:}\PYG{n}{endl}\PYG{p}{;}
\end{Verbatim}
\end{codeblockbg}


%%C080214 p

该程序的一种可能输出如下：


%%C080215 code

\begin{codeblockbg}
\begin{Verbatim}[commandchars=\\\{\},breaklines=true,breakanywhere=true,fontsize=\codefont,xleftmargin=17bp]
\PYG{n}{start}\PYG{+w}{ }\PYG{n}{producing}\PYG{+w}{ }\PYG{k}{and}\PYG{+w}{ }\PYG{n}{consuming}\PYG{p}{.}\PYG{p}{.}\PYG{p}{.}
\PYG{p}{[}\PYG{n}{produced}\PYG{p}{(}\PYG{l+m+mi}{5}\PYG{p}{)}\PYG{p}{]}\PYG{o}{:}\PYG{+w}{ }\PYG{l+m+mi}{550}
\PYG{p}{[}\PYG{n}{consumed}\PYG{p}{]}\PYG{o}{:}\PYG{+w}{ }\PYG{l+m+mi}{550}
\PYG{p}{[}\PYG{n}{produced}\PYG{p}{(}\PYG{l+m+mi}{5}\PYG{p}{)}\PYG{p}{]}\PYG{o}{:}\PYG{+w}{ }\PYG{l+m+mi}{529}
\PYG{p}{[}\PYG{n}{consumed}\PYG{p}{]}\PYG{o}{:}\PYG{+w}{ }\PYG{l+m+mi}{529}
\PYG{p}{[}\PYG{n}{produced}\PYG{p}{(}\PYG{l+m+mi}{5}\PYG{p}{)}\PYG{p}{]}\PYG{o}{:}\PYG{+w}{ }\PYG{l+m+mi}{537}
\PYG{p}{[}\PYG{n}{consumed}\PYG{p}{]}\PYG{o}{:}\PYG{+w}{ }\PYG{l+m+mi}{537}
\PYG{p}{[}\PYG{n}{produced}\PYG{p}{(}\PYG{l+m+mi}{1}\PYG{p}{)}\PYG{p}{]}\PYG{o}{:}\PYG{+w}{ }\PYG{l+m+mi}{122}
\PYG{p}{[}\PYG{n}{produced}\PYG{p}{(}\PYG{l+m+mi}{2}\PYG{p}{)}\PYG{p}{]}\PYG{o}{:}\PYG{+w}{ }\PYG{l+m+mi}{224}
\PYG{p}{[}\PYG{n}{produced}\PYG{p}{(}\PYG{l+m+mi}{3}\PYG{p}{)}\PYG{p}{]}\PYG{o}{:}\PYG{+w}{ }\PYG{l+m+mi}{326}
\PYG{p}{[}\PYG{n}{produced}\PYG{p}{(}\PYG{l+m+mi}{4}\PYG{p}{)}\PYG{p}{]}\PYG{o}{:}\PYG{+w}{ }\PYG{l+m+mi}{458}
\PYG{p}{[}\PYG{n}{consumed}\PYG{p}{]}\PYG{o}{:}\PYG{+w}{ }\PYG{l+m+mi}{122}
\PYG{p}{[}\PYG{n}{consumed}\PYG{p}{]}\PYG{o}{:}\PYG{+w}{ }\PYG{l+m+mi}{224}
\PYG{p}{[}\PYG{n}{consumed}\PYG{p}{]}\PYG{o}{:}\PYG{+w}{ }\PYG{l+m+mi}{326}
\PYG{p}{[}\PYG{n}{consumed}\PYG{p}{]}\PYG{o}{:}\PYG{+w}{ }\PYG{l+m+mi}{458}
\PYG{p}{.}\PYG{p}{.}\PYG{p}{.}
\PYG{n}{done}\PYG{+w}{ }\PYG{n}{producing}\PYG{+w}{ }\PYG{k}{and}\PYG{+w}{ }\PYG{n}{consuming}
\end{Verbatim}
\end{codeblockbg}


%%C080216 h3

\subsection{延伸阅读}


%%C080217 li

\begin{itemize}
\item \textit{使用线程}


%%C080218 li

\item \textit{用互斥量与锁同步对共享数据的访问}

\end{itemize}


%%C080219 h2

\section{用promise与future从线程返回值}


%%C080220 p

在本章的第一个案例中，我们讨论了如何使用线程。你还了解到，线程函数不能返回值，线程应当借助共享数据等其他手段来实现这一点；不过这样做需要同步。与主线程或其他线程沟通返回值或异常的另一种方式是使用\texttt{std::\allowbreak{}promise}。本案例将解释这种机制如何工作。


%%C080221 h3

\subsection{准备工作}


%%C080222 p

本案例中用到的\texttt{promise}和\texttt{future}类位于\texttt{<\allowbreak{}future>}头文件的\texttt{std}命名空间中。


%%C080223 h3

\subsection{具体做法}


%%C080224 p

要通过promise与future把一个值从一个线程传递到另一个线程，请这样做：


%%C080225 p

通过参数，让线程函数可以获得一个promise。


%%C080226 p

例如：


%%C080227 code

\begin{codeblockbg}
\begin{Verbatim}[commandchars=\\\{\},breaklines=true,breakanywhere=true,fontsize=\codefont,xleftmargin=17bp]
\PYG{k+kt}{void}\PYG{+w}{ }\PYG{n+nf}{produce\PYGZus{}value}\PYG{p}{(}\PYG{n}{std}\PYG{o}{:}\PYG{o}{:}\PYG{n}{promise}\PYG{o}{\PYGZlt{}}\PYG{k+kt}{int}\PYG{o}{\PYGZgt{}}\PYG{o}{\PYGZam{}}\PYG{+w}{ }\PYG{n}{p}\PYG{p}{)}
\PYG{p}{\PYGZob{}}
\PYG{+w}{  }\PYG{c+c1}{// 模拟长时间运行的操作}
\PYG{+w}{  }\PYG{p}{\PYGZob{}}
\PYG{+w}{    }\PYG{k}{using}\PYG{+w}{ }\PYG{k}{namespace}\PYG{+w}{ }\PYG{n+nn}{std}\PYG{o}{:}\PYG{o}{:}\PYG{n+nn}{chrono\PYGZus{}literals}\PYG{p}{;}
\PYG{+w}{    }\PYG{n}{std}\PYG{o}{:}\PYG{o}{:}\PYG{n}{this\PYGZus{}thread}\PYG{o}{:}\PYG{o}{:}\PYG{n}{sleep\PYGZus{}for}\PYG{p}{(}\PYG{l+m+mi}{2}\PYG{n}{s}\PYG{p}{)}\PYG{p}{;}
\PYG{+w}{  }\PYG{p}{\PYGZcb{}}
\PYG{+w}{  }\PYG{c+c1}{// 接第2步。}
\PYG{p}{\PYGZcb{}}
\end{Verbatim}
\end{codeblockbg}


%%C080228 p

在promise上调用\texttt{set\_\allowbreak{}value()}，把结果设为表示一个值，或者调用


%%C080229 p

\texttt{set\_\allowbreak{}exception()}，把结果设为表示一个异常：


%%C080230 code

\begin{codeblockbg}
\begin{Verbatim}[commandchars=\\\{\},breaklines=true,breakanywhere=true,fontsize=\codefont,xleftmargin=17bp]
\PYG{c+c1}{// 接第1步。}
\PYG{n}{p}\PYG{p}{.}\PYG{n}{set\PYGZus{}value}\PYG{p}{(}\PYG{l+m+mi}{42}\PYG{p}{)}\PYG{p}{;}
\end{Verbatim}
\end{codeblockbg}


%%C080231 p

通过参数，让另一个线程函数可以获得与该promise关联的future。


%%C080232 p

例如：


%%C080233 code

\begin{codeblockbg}
\begin{Verbatim}[commandchars=\\\{\},breaklines=true,breakanywhere=true,fontsize=\codefont,xleftmargin=17bp]
\PYG{k+kt}{void}\PYG{+w}{ }\PYG{n+nf}{consume\PYGZus{}value}\PYG{p}{(}\PYG{n}{std}\PYG{o}{:}\PYG{o}{:}\PYG{n}{future}\PYG{o}{\PYGZlt{}}\PYG{k+kt}{int}\PYG{o}{\PYGZgt{}}\PYG{o}{\PYGZam{}}\PYG{+w}{ }\PYG{n}{f}\PYG{p}{)}
\PYG{p}{\PYGZob{}}
\PYG{+w}{  }\PYG{c+c1}{// 接第4步。}
\PYG{p}{\PYGZcb{}}
\end{Verbatim}
\end{codeblockbg}


%%C080234 p

在\texttt{future}对象上调用\texttt{get()}，以取得设置到promise上的结果：


%%C080235 code

\begin{codeblockbg}
\begin{Verbatim}[commandchars=\\\{\},breaklines=true,breakanywhere=true,fontsize=\codefont,xleftmargin=17bp]
\PYG{c+c1}{// 接第3步。}
\PYG{k}{auto}\PYG{+w}{ }\PYG{n}{value}\PYG{+w}{ }\PYG{o}{=}\PYG{+w}{ }\PYG{n}{f}\PYG{p}{.}\PYG{n}{get}\PYG{p}{(}\PYG{p}{)}\PYG{p}{;}
\end{Verbatim}
\end{codeblockbg}


%%C080236 p

在调用线程中，对promise使用\texttt{get\_\allowbreak{}future()}，以获得与该promise关联的


%%C080237 p

\texttt{future}：


%%C080238 code

\begin{codeblockbg}
\begin{Verbatim}[commandchars=\\\{\},breaklines=true,breakanywhere=true,fontsize=\codefont,xleftmargin=17bp]
\PYG{n}{std}\PYG{o}{:}\PYG{o}{:}\PYG{n}{promise}\PYG{o}{\PYGZlt{}}\PYG{k+kt}{int}\PYG{o}{\PYGZgt{}}\PYG{+w}{ }\PYG{n}{p}\PYG{p}{;}
\PYG{n}{std}\PYG{o}{:}\PYG{o}{:}\PYG{k+kr}{thread}\PYG{+w}{ }\PYG{n+nf}{t1}\PYG{p}{(}\PYG{n}{produce\PYGZus{}value}\PYG{p}{,}\PYG{+w}{ }\PYG{n}{std}\PYG{o}{:}\PYG{o}{:}\PYG{n}{ref}\PYG{p}{(}\PYG{n}{p}\PYG{p}{)}\PYG{p}{)}\PYG{p}{;}
\PYG{n}{std}\PYG{o}{:}\PYG{o}{:}\PYG{n}{future}\PYG{o}{\PYGZlt{}}\PYG{k+kt}{int}\PYG{o}{\PYGZgt{}}\PYG{+w}{ }\PYG{n}{f}\PYG{+w}{ }\PYG{o}{=}\PYG{+w}{ }\PYG{n}{p}\PYG{p}{.}\PYG{n}{get\PYGZus{}future}\PYG{p}{(}\PYG{p}{)}\PYG{p}{;}
\PYG{n}{std}\PYG{o}{:}\PYG{o}{:}\PYG{k+kr}{thread}\PYG{+w}{ }\PYG{n+nf}{t2}\PYG{p}{(}\PYG{n}{consume\PYGZus{}value}\PYG{p}{,}\PYG{+w}{ }\PYG{n}{std}\PYG{o}{:}\PYG{o}{:}\PYG{n}{ref}\PYG{p}{(}\PYG{n}{f}\PYG{p}{)}\PYG{p}{)}\PYG{p}{;}
\PYG{n}{t1}\PYG{p}{.}\PYG{n}{join}\PYG{p}{(}\PYG{p}{)}\PYG{p}{;}
\PYG{n}{t2}\PYG{p}{.}\PYG{n}{join}\PYG{p}{(}\PYG{p}{)}\PYG{p}{;}
\end{Verbatim}
\end{codeblockbg}


%%C080239 h3

\subsection{工作原理}


%%C080240 p

promise-future对本质上是一条通信通道，它使一个线程能够通过共享状态把值或异常传递给另一个线程。\texttt{promise}是结果的异步提供者，它有一个关联的\texttt{future}，表示异步返回对象。要建立这条通道，必须首先创建一个promise。这同时会创建共享状态，之后可以通过与该promise关联的future读取它。


%%C080241 p

要给promise设置结果，可以使用以下任意一种方法：


%%C080242 li

\begin{itemize}
\item 用\texttt{set\_\allowbreak{}value()}或\myttbr{set\_\allowbreak{}value\_\allowbreak{}at\_\allowbreak{}thread\_\allowbreak{}exit()}方法设置返回值；后者把值存入共享状态，但只有在线程退出时才通过关联的future使其可用。


%%C080243 li

\item 用\texttt{set\_\allowbreak{}exception()}或\myttbr{set\_\allowbreak{}exception\_\allowbreak{}at\_\allowbreak{}thread\_\allowbreak{}exit()}方法把异常设为返回值。该异常被包装在\myttbr{std::\allowbreak{}exception\_\allowbreak{}ptr}对象中。后者把异常存入共享状态，但只有在线程退出时才使其可用。

\end{itemize}


%%C080244 p

要获取与\texttt{promise}关联的\texttt{future}对象，请使用\texttt{get\_\allowbreak{}future()}方法。要从future中取值，请使用\texttt{get()}方法。这会阻塞调用线程，直到共享状态中的值变为可用。future类提供了若干方法，用于阻塞线程直到共享状态中的结果变为可用：


%%C080245 li

\begin{itemize}
\item \texttt{wait()}只在结果可用时才返回。


%%C080246 li

\item \texttt{wait\_\allowbreak{}for()}在结果可用或指定的超时到期时返回。


%%C080247 li

\item \texttt{wait\_\allowbreak{}until()}在结果可用或到达指定的时间点时返回。

\end{itemize}


%%C080248 p

如果给\texttt{promise}设置的是异常，那么在future上调用\texttt{get()}方法将抛出该异常。上一节的示例改写为下面的代码，以抛出异常而不是设置结果。对\texttt{get()}的调用被放在\texttt{try...\allowbreak{}catch}块中，如果捕获到异常，就把它的消息打印到控制台：


%%C080249 code

\begin{codeblockbg}
\begin{Verbatim}[commandchars=\\\{\},breaklines=true,breakanywhere=true,fontsize=\codefont,xleftmargin=17bp]
\PYG{k+kt}{void}\PYG{+w}{ }\PYG{n+nf}{produce\PYGZus{}value}\PYG{p}{(}\PYG{n}{std}\PYG{o}{:}\PYG{o}{:}\PYG{n}{promise}\PYG{o}{\PYGZlt{}}\PYG{k+kt}{int}\PYG{o}{\PYGZgt{}}\PYG{o}{\PYGZam{}}\PYG{+w}{ }\PYG{n}{p}\PYG{p}{)}
\PYG{p}{\PYGZob{}}
\PYG{+w}{  }\PYG{c+c1}{// 模拟长时间运行的操作}
\PYG{+w}{  }\PYG{p}{\PYGZob{}}
\PYG{+w}{    }\PYG{k}{using}\PYG{+w}{ }\PYG{k}{namespace}\PYG{+w}{ }\PYG{n+nn}{std}\PYG{o}{:}\PYG{o}{:}\PYG{n+nn}{chrono\PYGZus{}literals}\PYG{p}{;}
\PYG{+w}{    }\PYG{n}{std}\PYG{o}{:}\PYG{o}{:}\PYG{n}{this\PYGZus{}thread}\PYG{o}{:}\PYG{o}{:}\PYG{n}{sleep\PYGZus{}for}\PYG{p}{(}\PYG{l+m+mi}{2}\PYG{n}{s}\PYG{p}{)}\PYG{p}{;}
\PYG{+w}{  }\PYG{p}{\PYGZcb{}}
\PYG{+w}{  }\PYG{k}{try}
\PYG{+w}{  }\PYG{p}{\PYGZob{}}
\PYG{+w}{    }\PYG{k}{throw}\PYG{+w}{ }\PYG{n}{std}\PYG{o}{:}\PYG{o}{:}\PYG{n}{runtime\PYGZus{}error}\PYG{p}{(}\PYG{l+s}{\PYGZdq{}}\PYG{l+s}{an error has occurred!}\PYG{l+s}{\PYGZdq{}}\PYG{p}{)}\PYG{p}{;}
\PYG{+w}{  }\PYG{p}{\PYGZcb{}}
\PYG{+w}{  }\PYG{k}{catch}\PYG{p}{(}\PYG{p}{.}\PYG{p}{.}\PYG{p}{.}\PYG{p}{)}
\PYG{+w}{  }\PYG{p}{\PYGZob{}}
\PYG{+w}{    }\PYG{n}{p}\PYG{p}{.}\PYG{n}{set\PYGZus{}exception}\PYG{p}{(}\PYG{n}{std}\PYG{o}{:}\PYG{o}{:}\PYG{n}{current\PYGZus{}exception}\PYG{p}{(}\PYG{p}{)}\PYG{p}{)}\PYG{p}{;}
\PYG{+w}{  }\PYG{p}{\PYGZcb{}}
\PYG{p}{\PYGZcb{}}
\PYG{k+kt}{void}\PYG{+w}{ }\PYG{n}{consume\PYGZus{}value}\PYG{p}{(}\PYG{n}{std}\PYG{o}{:}\PYG{o}{:}\PYG{n}{future}\PYG{o}{\PYGZlt{}}\PYG{k+kt}{int}\PYG{o}{\PYGZgt{}}\PYG{o}{\PYGZam{}}\PYG{+w}{ }\PYG{n}{f}\PYG{p}{)}
\PYG{p}{\PYGZob{}}
\PYG{+w}{  }\PYG{n}{std}\PYG{o}{:}\PYG{o}{:}\PYG{n}{lock\PYGZus{}guard}\PYG{o}{\PYGZlt{}}\PYG{n}{std}\PYG{o}{:}\PYG{o}{:}\PYG{n}{mutex}\PYG{o}{\PYGZgt{}}\PYG{+w}{ }\PYG{n}{lock}\PYG{p}{(}\PYG{n}{g\PYGZus{}mutex}\PYG{p}{)}\PYG{p}{;}
\PYG{+w}{  }\PYG{k}{try}
\PYG{+w}{  }\PYG{p}{\PYGZob{}}
\PYG{+w}{    }\PYG{n}{std}\PYG{o}{:}\PYG{o}{:}\PYG{n}{cout}\PYG{+w}{ }\PYG{o}{\PYGZlt{}}\PYG{o}{\PYGZlt{}}\PYG{+w}{ }\PYG{n}{f}\PYG{p}{.}\PYG{n}{get}\PYG{p}{(}\PYG{p}{)}\PYG{+w}{ }\PYG{o}{\PYGZlt{}}\PYG{o}{\PYGZlt{}}\PYG{+w}{ }\PYG{n}{std}\PYG{o}{:}\PYG{o}{:}\PYG{n}{endl}\PYG{p}{;}
\PYG{+w}{  }\PYG{p}{\PYGZcb{}}
\PYG{+w}{  }\PYG{k}{catch}\PYG{p}{(}\PYG{n}{std}\PYG{o}{:}\PYG{o}{:}\PYG{n}{exception}\PYG{+w}{ }\PYG{k}{const}\PYG{+w}{ }\PYG{o}{\PYGZam{}}\PYG{+w}{ }\PYG{n}{e}\PYG{p}{)}
\PYG{+w}{  }\PYG{p}{\PYGZob{}}
\PYG{+w}{    }\PYG{n}{std}\PYG{o}{:}\PYG{o}{:}\PYG{n}{cout}\PYG{+w}{ }\PYG{o}{\PYGZlt{}}\PYG{o}{\PYGZlt{}}\PYG{+w}{ }\PYG{n}{e}\PYG{p}{.}\PYG{n}{what}\PYG{p}{(}\PYG{p}{)}\PYG{+w}{ }\PYG{o}{\PYGZlt{}}\PYG{o}{\PYGZlt{}}\PYG{+w}{ }\PYG{n}{std}\PYG{o}{:}\PYG{o}{:}\PYG{n}{endl}\PYG{p}{;}
\PYG{+w}{  }\PYG{p}{\PYGZcb{}}
\PYG{p}{\PYGZcb{}}
\end{Verbatim}
\end{codeblockbg}


%%C080250 h3

\subsection{更多内容}


%%C080251 p

以这种方式建立promise-future通道是相当显式的操作，使用\texttt{std::\allowbreak{}async()}函数可以避免这些操作；它是一个更高层的工具，会异步运行函数，在内部创建一个promise和共享状态，并返回与该共享状态关联的future。我们将在下一个案例\textit{异步执行函数}中看到\texttt{std::\allowbreak{}async()}如何工作。


%%C080252 h3

\subsection{延伸阅读}


%%C080253 li

\begin{itemize}
\item \textit{使用线程}


%%C080254 li

\item \textit{处理线程函数抛出的异常}

\end{itemize}


%%C080255 h2

\section{异步执行函数}


%%C080256 p

线程使我们能够同时运行多个函数，这有助于利用多处理器或多核系统中的硬件设施。然而，线程需要显式的底层操作。线程的替代方案是任务，任务是在某个特定线程中运行的工作单元。C++标准没有提供完整的任务库，但它允许开发者在不同线程上异步执行函数，并通过promise-future通道把结果传回，如上一案例所示。本案例将展示如何使用\texttt{std::\allowbreak{}async()}和\texttt{std::\allowbreak{}future}做到这一点。


%%C080257 h3

\subsection{准备工作}


%%C080258 p

我们会用到future，所以请阅读上一案例，快速了解它们如何工作。\texttt{async()}和\texttt{future}都位于\texttt{<\allowbreak{}future>}头文件的\texttt{std}命名空间中。


%%C080259 p

本案例的示例将使用以下函数：


%%C080260 code

\begin{codeblockbg}
\begin{Verbatim}[commandchars=\\\{\},breaklines=true,breakanywhere=true,fontsize=\codefont,xleftmargin=17bp]
\PYG{k+kt}{void}\PYG{+w}{ }\PYG{n+nf}{do\PYGZus{}something}\PYG{p}{(}\PYG{p}{)}
\PYG{p}{\PYGZob{}}
\PYG{+w}{  }\PYG{c+c1}{// 模拟长时间运行的操作}
\PYG{+w}{  }\PYG{p}{\PYGZob{}}
\PYG{+w}{    }\PYG{k}{using}\PYG{+w}{ }\PYG{k}{namespace}\PYG{+w}{ }\PYG{n+nn}{std}\PYG{o}{:}\PYG{o}{:}\PYG{n+nn}{chrono\PYGZus{}literals}\PYG{p}{;}
\PYG{+w}{    }\PYG{n}{std}\PYG{o}{:}\PYG{o}{:}\PYG{n}{this\PYGZus{}thread}\PYG{o}{:}\PYG{o}{:}\PYG{n}{sleep\PYGZus{}for}\PYG{p}{(}\PYG{l+m+mi}{2}\PYG{n}{s}\PYG{p}{)}\PYG{p}{;}
\PYG{+w}{  }\PYG{p}{\PYGZcb{}}
\PYG{+w}{  }\PYG{n}{std}\PYG{o}{:}\PYG{o}{:}\PYG{n}{lock\PYGZus{}guard}\PYG{o}{\PYGZlt{}}\PYG{n}{std}\PYG{o}{:}\PYG{o}{:}\PYG{n}{mutex}\PYG{o}{\PYGZgt{}}\PYG{+w}{ }\PYG{n}{lock}\PYG{p}{(}\PYG{n}{g\PYGZus{}mutex}\PYG{p}{)}\PYG{p}{;}
\PYG{+w}{  }\PYG{n}{std}\PYG{o}{:}\PYG{o}{:}\PYG{n}{cout}\PYG{+w}{ }\PYG{o}{\PYGZlt{}}\PYG{o}{\PYGZlt{}}\PYG{+w}{ }\PYG{l+s}{\PYGZdq{}}\PYG{l+s}{operation 1 done}\PYG{l+s}{\PYGZdq{}}\PYG{+w}{ }\PYG{o}{\PYGZlt{}}\PYG{o}{\PYGZlt{}}\PYG{+w}{ }\PYG{n}{std}\PYG{o}{:}\PYG{o}{:}\PYG{n}{endl}\PYG{p}{;}
\PYG{p}{\PYGZcb{}}
\PYG{k+kt}{void}\PYG{+w}{ }\PYG{n}{do\PYGZus{}something\PYGZus{}else}\PYG{p}{(}\PYG{p}{)}
\PYG{p}{\PYGZob{}}
\PYG{+w}{  }\PYG{c+c1}{// 模拟长时间运行的操作}
\PYG{+w}{  }\PYG{p}{\PYGZob{}}
\PYG{+w}{    }\PYG{k}{using}\PYG{+w}{ }\PYG{k}{namespace}\PYG{+w}{ }\PYG{n+nn}{std}\PYG{o}{:}\PYG{o}{:}\PYG{n+nn}{chrono\PYGZus{}literals}\PYG{p}{;}
\PYG{+w}{    }\PYG{n}{std}\PYG{o}{:}\PYG{o}{:}\PYG{n}{this\PYGZus{}thread}\PYG{o}{:}\PYG{o}{:}\PYG{n}{sleep\PYGZus{}for}\PYG{p}{(}\PYG{l+m+mi}{1}\PYG{n}{s}\PYG{p}{)}\PYG{p}{;}
\PYG{+w}{  }\PYG{p}{\PYGZcb{}}
\PYG{+w}{  }\PYG{n}{std}\PYG{o}{:}\PYG{o}{:}\PYG{n}{lock\PYGZus{}guard}\PYG{o}{\PYGZlt{}}\PYG{n}{std}\PYG{o}{:}\PYG{o}{:}\PYG{n}{mutex}\PYG{o}{\PYGZgt{}}\PYG{+w}{ }\PYG{n}{lock}\PYG{p}{(}\PYG{n}{g\PYGZus{}mutex}\PYG{p}{)}\PYG{p}{;}
\PYG{+w}{  }\PYG{n}{std}\PYG{o}{:}\PYG{o}{:}\PYG{n}{cout}\PYG{+w}{ }\PYG{o}{\PYGZlt{}}\PYG{o}{\PYGZlt{}}\PYG{+w}{ }\PYG{l+s}{\PYGZdq{}}\PYG{l+s}{operation 2 done}\PYG{l+s}{\PYGZdq{}}\PYG{+w}{ }\PYG{o}{\PYGZlt{}}\PYG{o}{\PYGZlt{}}\PYG{+w}{ }\PYG{n}{std}\PYG{o}{:}\PYG{o}{:}\PYG{n}{endl}\PYG{p}{;}
\PYG{p}{\PYGZcb{}}
\PYG{k+kt}{int}\PYG{+w}{ }\PYG{n}{compute\PYGZus{}something}\PYG{p}{(}\PYG{p}{)}
\PYG{p}{\PYGZob{}}
\PYG{+w}{  }\PYG{c+c1}{// 模拟长时间运行的操作}
\PYG{+w}{  }\PYG{p}{\PYGZob{}}
\PYG{+w}{    }\PYG{k}{using}\PYG{+w}{ }\PYG{k}{namespace}\PYG{+w}{ }\PYG{n+nn}{std}\PYG{o}{:}\PYG{o}{:}\PYG{n+nn}{chrono\PYGZus{}literals}\PYG{p}{;}
\PYG{+w}{    }\PYG{n}{std}\PYG{o}{:}\PYG{o}{:}\PYG{n}{this\PYGZus{}thread}\PYG{o}{:}\PYG{o}{:}\PYG{n}{sleep\PYGZus{}for}\PYG{p}{(}\PYG{l+m+mi}{2}\PYG{n}{s}\PYG{p}{)}\PYG{p}{;}
\PYG{+w}{  }\PYG{p}{\PYGZcb{}}
\PYG{+w}{  }\PYG{k}{return}\PYG{+w}{ }\PYG{l+m+mi}{42}\PYG{p}{;}
\PYG{p}{\PYGZcb{}}
\PYG{k+kt}{int}\PYG{+w}{ }\PYG{n}{compute\PYGZus{}something\PYGZus{}else}\PYG{p}{(}\PYG{p}{)}
\PYG{p}{\PYGZob{}}
\PYG{+w}{  }\PYG{c+c1}{// 模拟长时间运行的操作}
\PYG{+w}{  }\PYG{p}{\PYGZob{}}
\PYG{+w}{    }\PYG{k}{using}\PYG{+w}{ }\PYG{k}{namespace}\PYG{+w}{ }\PYG{n+nn}{std}\PYG{o}{:}\PYG{o}{:}\PYG{n+nn}{chrono\PYGZus{}literals}\PYG{p}{;}
\PYG{+w}{    }\PYG{n}{std}\PYG{o}{:}\PYG{o}{:}\PYG{n}{this\PYGZus{}thread}\PYG{o}{:}\PYG{o}{:}\PYG{n}{sleep\PYGZus{}for}\PYG{p}{(}\PYG{l+m+mi}{1}\PYG{n}{s}\PYG{p}{)}\PYG{p}{;}
\PYG{+w}{  }\PYG{p}{\PYGZcb{}}
\PYG{+w}{  }\PYG{k}{return}\PYG{+w}{ }\PYG{l+m+mi}{24}\PYG{p}{;}
\PYG{p}{\PYGZcb{}}
\end{Verbatim}
\end{codeblockbg}


%%C080261 h3

\subsection{具体做法}


%%C080262 p

当当前线程继续执行且不等待结果时，要在另一个线程上异步执行函数，请这样做：


%%C080263 p

用\texttt{std::\allowbreak{}async()}启动一个新线程来执行指定的函数，这会创建一个异步提供者，并返回与之关联的


%%C080264 p

\texttt{future}。要把\texttt{std::\allowbreak{}launch::\allowbreak{}async}策略作为函数的第一个参数传入，以确保函数会异步运行：


%%C080265 code

\begin{codeblockbg}
\begin{Verbatim}[commandchars=\\\{\},breaklines=true,breakanywhere=true,fontsize=\codefont,xleftmargin=17bp]
\PYG{k}{auto}\PYG{+w}{ }\PYG{n}{f}\PYG{+w}{ }\PYG{o}{=}\PYG{+w}{ }\PYG{n}{std}\PYG{o}{:}\PYG{o}{:}\PYG{n}{async}\PYG{p}{(}\PYG{n}{std}\PYG{o}{:}\PYG{o}{:}\PYG{n}{launch}\PYG{o}{:}\PYG{o}{:}\PYG{n}{async}\PYG{p}{,}\PYG{+w}{ }\PYG{n}{do\PYGZus{}something}\PYG{p}{)}\PYG{p}{;}
\end{Verbatim}
\end{codeblockbg}


%%C080266 p

让当前线程继续执行：


%%C080267 code

\begin{codeblockbg}
\begin{Verbatim}[commandchars=\\\{\},breaklines=true,breakanywhere=true,fontsize=\codefont,xleftmargin=17bp]
\PYG{n}{do\PYGZus{}something\PYGZus{}else}\PYG{p}{(}\PYG{p}{)}\PYG{p}{;}
\end{Verbatim}
\end{codeblockbg}


%%C080268 p

当你需要确保异步操作已经完成时，对\texttt{std::\allowbreak{}async()}返回的\texttt{future}对象调用


%%C080269 p

\texttt{wait()}方法：


%%C080270 code

\begin{codeblockbg}
\begin{Verbatim}[commandchars=\\\{\},breaklines=true,breakanywhere=true,fontsize=\codefont,xleftmargin=17bp]
\PYG{n}{f}\PYG{p}{.}\PYG{n}{wait}\PYG{p}{(}\PYG{p}{)}\PYG{p}{;}
\end{Verbatim}
\end{codeblockbg}


%%C080271 p

要在工作线程上异步执行函数，同时让当前线程继续执行，直到当前线程需要该异步函数的结果为止，请这样做：


%%C080272 p

用\texttt{std::\allowbreak{}async()}启动一个新线程来执行指定的函数，创建一个异步提供者，并返回与之关联的


%%C080273 p

\texttt{future}。要用函数的第一个参数指定\texttt{std::\allowbreak{}launch::\allowbreak{}async}策略，以确保函数确实异步运行：


%%C080274 code

\begin{codeblockbg}
\begin{Verbatim}[commandchars=\\\{\},breaklines=true,breakanywhere=true,fontsize=\codefont,xleftmargin=17bp]
\PYG{k}{auto}\PYG{+w}{ }\PYG{n}{f}\PYG{+w}{ }\PYG{o}{=}\PYG{+w}{ }\PYG{n}{std}\PYG{o}{:}\PYG{o}{:}\PYG{n}{async}\PYG{p}{(}\PYG{n}{std}\PYG{o}{:}\PYG{o}{:}\PYG{n}{launch}\PYG{o}{:}\PYG{o}{:}\PYG{n}{async}\PYG{p}{,}\PYG{+w}{ }\PYG{n}{compute\PYGZus{}something}\PYG{p}{)}\PYG{p}{;}
\end{Verbatim}
\end{codeblockbg}


%%C080275 p

让当前线程继续执行：


%%C080276 code

\begin{codeblockbg}
\begin{Verbatim}[commandchars=\\\{\},breaklines=true,breakanywhere=true,fontsize=\codefont,xleftmargin=17bp]
\PYG{k}{auto}\PYG{+w}{ }\PYG{n}{value}\PYG{+w}{ }\PYG{o}{=}\PYG{+w}{ }\PYG{n}{compute\PYGZus{}something\PYGZus{}else}\PYG{p}{(}\PYG{p}{)}\PYG{p}{;}
\end{Verbatim}
\end{codeblockbg}


%%C080277 p

当你需要取得异步执行的函数的结果时，对\texttt{std::\allowbreak{}async()}返回的\texttt{future}对象调用


%%C080278 p

\texttt{get()}方法：


%%C080279 code

\begin{codeblockbg}
\begin{Verbatim}[commandchars=\\\{\},breaklines=true,breakanywhere=true,fontsize=\codefont,xleftmargin=17bp]
\PYG{n}{value}\PYG{+w}{ }\PYG{o}{+}\PYG{o}{=}\PYG{+w}{ }\PYG{n}{f}\PYG{p}{.}\PYG{n}{get}\PYG{p}{(}\PYG{p}{)}\PYG{p}{;}
\end{Verbatim}
\end{codeblockbg}


%%C080280 h3

\subsection{工作原理}


%%C080281 p

\texttt{std::\allowbreak{}async()}是一个可变参数函数模板，有两个重载：一个把启动策略指定为第一个参数，另一个不指定。\texttt{std::\allowbreak{}async()}的其他参数是要执行的函数及其参数（如果有）。启动策略由名为\texttt{std::\allowbreak{}launch}的限定作用域枚举（scoped enumeration）定义，该枚举位于\texttt{<\allowbreak{}future>}头文件中：


%%C080282 code

\begin{codeblockbg}
\begin{Verbatim}[commandchars=\\\{\},breaklines=true,breakanywhere=true,fontsize=\codefont,xleftmargin=17bp]
\PYG{k}{enum}\PYG{+w}{ }\PYG{k}{class}\PYG{+w}{ }\PYG{n+nc}{launch}\PYG{+w}{ }\PYG{o}{:}\PYG{+w}{ }\PYG{c+cm}{/* 未指定 */}
\PYG{p}{\PYGZob{}}
\PYG{+w}{  }\PYG{n}{async}\PYG{+w}{ }\PYG{o}{=}\PYG{+w}{ }\PYG{c+cm}{/* 未指定 */}\PYG{p}{,}
\PYG{+w}{  }\PYG{n}{deferred}\PYG{+w}{ }\PYG{o}{=}\PYG{+w}{ }\PYG{c+cm}{/* 未指定 */}\PYG{p}{,}
\PYG{+w}{  }\PYG{c+cm}{/* 实现定义 */}
\PYG{p}{\PYGZcb{}}\PYG{p}{;}
\end{Verbatim}
\end{codeblockbg}


%%C080283 p

两个可用的启动策略含义如下：


%%C080284 li

\begin{itemize}
\item 使用\texttt{async}时，会启动一个新线程来异步执行任务。


%%C080285 li

\item 使用\texttt{deferred}时，任务在第一次请求其结果时在调用线程上执行。

\end{itemize}


%%C080286 p

当两个标志同时被指定时（\myttbr{std::\allowbreak{}launch::\allowbreak{}async | std::\allowbreak{}launch::\allowbreak{}deferred}），是在新线程上异步运行任务，还是在当前线程上同步运行任务，由实现自行决定。不指定启动策略的\texttt{std::\allowbreak{}async()}重载就是这种行为。这种行为不是确定性的。


%%C080287 note

\begin{myNote}[tip]
不要使用\texttt{std::\allowbreak{}async()}的非确定性重载来异步运行任务。始终使用要求指定启动策略的重载，并始终使用\texttt{std::\allowbreak{}launch::\allowbreak{}async}。
\end{myNote}


%%C080288 p

\texttt{std::\allowbreak{}async()}的两个重载都返回一个\texttt{future}对象，它引用\texttt{std::\allowbreak{}async()}为其建立的promise-future通道在内部创建的共享状态。当你需要异步操作的结果时，在future上调用\texttt{get()}方法。这会阻塞当前线程，直到结果值或异常变为可用。如果future不传送任何值，或者你并不关心那个值、只想确保异步操作最终会完成，那么就使用\texttt{wait()}方法；它会阻塞当前线程，直到共享状态通过future变为可用。


%%C080289 p

future类还有两个等待方法：\texttt{wait\_\allowbreak{}for()}指定一个时长，超过该时长后调用就结束并返回，即使共享状态尚未通过future变为可用；\texttt{wait\_\allowbreak{}until()}指定一个时间点，超过该时间点后调用就返回，即使共享状态尚不可用。可以用这些方法编写轮询例程，向用户显示状态消息，如下例所示：


%%C080290 code

\begin{codeblockbg}
\begin{Verbatim}[commandchars=\\\{\},breaklines=true,breakanywhere=true,fontsize=\codefont,xleftmargin=17bp]
\PYG{k}{auto}\PYG{+w}{ }\PYG{n}{f}\PYG{+w}{ }\PYG{o}{=}\PYG{+w}{ }\PYG{n}{std}\PYG{o}{:}\PYG{o}{:}\PYG{n}{async}\PYG{p}{(}\PYG{n}{std}\PYG{o}{:}\PYG{o}{:}\PYG{n}{launch}\PYG{o}{:}\PYG{o}{:}\PYG{n}{async}\PYG{p}{,}\PYG{+w}{ }\PYG{n}{do\PYGZus{}something}\PYG{p}{)}\PYG{p}{;}
\PYG{k}{while}\PYG{p}{(}\PYG{n+nb}{true}\PYG{p}{)}
\PYG{p}{\PYGZob{}}
\PYG{+w}{  }\PYG{k}{using}\PYG{+w}{ }\PYG{k}{namespace}\PYG{+w}{ }\PYG{n+nn}{std}\PYG{o}{:}\PYG{o}{:}\PYG{n+nn}{chrono\PYGZus{}literals}\PYG{p}{;}
\PYG{+w}{  }\PYG{k}{auto}\PYG{+w}{ }\PYG{n}{status}\PYG{+w}{ }\PYG{o}{=}\PYG{+w}{ }\PYG{n}{f}\PYG{p}{.}\PYG{n}{wait\PYGZus{}for}\PYG{p}{(}\PYG{l+m+mi}{500}\PYG{n}{ms}\PYG{p}{)}\PYG{p}{;}
\PYG{k}{if}\PYG{p}{(}\PYG{n}{status}\PYG{+w}{ }\PYG{o}{=}\PYG{o}{=}\PYG{+w}{ }\PYG{n}{std}\PYG{o}{:}\PYG{o}{:}\PYG{n}{future\PYGZus{}status}\PYG{o}{:}\PYG{o}{:}\PYG{n}{ready}\PYG{p}{)}
\PYG{+w}{  }\PYG{k}{break}\PYG{p}{;}
\PYG{+w}{  }\PYG{n}{std}\PYG{o}{:}\PYG{o}{:}\PYG{n}{cout}\PYG{+w}{ }\PYG{o}{\PYGZlt{}}\PYG{o}{\PYGZlt{}}\PYG{+w}{ }\PYG{l+s}{\PYGZdq{}}\PYG{l+s}{waiting...}\PYG{l+s}{\PYGZdq{}}\PYG{+w}{ }\PYG{o}{\PYGZlt{}}\PYG{o}{\PYGZlt{}}\PYG{+w}{ }\PYG{n}{std}\PYG{o}{:}\PYG{o}{:}\PYG{n}{endl}\PYG{p}{;}
\PYG{p}{\PYGZcb{}}
\PYG{n}{std}\PYG{o}{:}\PYG{o}{:}\PYG{n}{cout}\PYG{+w}{ }\PYG{o}{\PYGZlt{}}\PYG{o}{\PYGZlt{}}\PYG{+w}{ }\PYG{l+s}{\PYGZdq{}}\PYG{l+s}{done!}\PYG{l+s}{\PYGZdq{}}\PYG{+w}{ }\PYG{o}{\PYGZlt{}}\PYG{o}{\PYGZlt{}}\PYG{+w}{ }\PYG{n}{std}\PYG{o}{:}\PYG{o}{:}\PYG{n}{endl}\PYG{p}{;}
\end{Verbatim}
\end{codeblockbg}


%%C080291 p

运行该程序的结果如下：


%%C080292 code

\begin{codeblockbg}
\begin{Verbatim}[commandchars=\\\{\},breaklines=true,breakanywhere=true,fontsize=\codefont,xleftmargin=17bp]
\PYG{n}{waiting}\PYG{p}{.}\PYG{p}{.}\PYG{p}{.}
\PYG{n}{waiting}\PYG{p}{.}\PYG{p}{.}\PYG{p}{.}
\PYG{n}{waiting}\PYG{p}{.}\PYG{p}{.}\PYG{p}{.}
\PYG{n}{operation}\PYG{+w}{ }\PYG{l+m+mi}{1}\PYG{+w}{ }\PYG{n}{done}
\PYG{n}{done}\PYG{o}{!}
\end{Verbatim}
\end{codeblockbg}


%%C080293 h3

\subsection{延伸阅读}


%%C080294 li

\begin{itemize}
\item \textit{用promise与future从线程返回值}

\end{itemize}


%%C080295 h2

\section{使用原子类型}


%%C080296 p

线程库提供了管理线程以及用互斥量和锁同步对共享数据访问的支持。标准库还提供了与之互补的、更底层的原子操作支持，即不可分割的操作，它们可以在不同线程中对共享数据并发执行，既不会产生竞争条件，也不需要加锁。这些支持包括原子类型、原子操作和内存同步顺序。本案例将展示如何使用其中一些类型和函数。


%%C080297 h3

\subsection{准备工作}


%%C080298 p

所有原子类型和操作都定义在\texttt{<\allowbreak{}atomic>}头文件的\texttt{std}命名空间中。


%%C080299 h3

\subsection{具体做法}


%%C080300 p

下面是一系列使用原子类型的典型操作：

%%C080301 li

\begin{itemize}
\item 使用\texttt{std::\allowbreak{}atomic}类模板创建支持原子操作的对象，这些操作包括加载、存储，以及算术运算或位运算：


%%C080302 code

\begin{codeblockbg}
\begin{Verbatim}[commandchars=\\\{\},breaklines=true,breakanywhere=true,fontsize=\codefont,xleftmargin=17bp]
\PYG{n}{std}\PYG{o}{:}\PYG{o}{:}\PYG{n}{atomic}\PYG{o}{\PYGZlt{}}\PYG{k+kt}{int}\PYG{o}{\PYGZgt{}}\PYG{+w}{ }\PYG{n}{counter}\PYG{+w}{ }\PYG{p}{\PYGZob{}}\PYG{l+m+mi}{0}\PYG{p}{\PYGZcb{}}\PYG{p}{;}
\PYG{n}{std}\PYG{o}{:}\PYG{o}{:}\PYG{n}{vector}\PYG{o}{\PYGZlt{}}\PYG{n}{std}\PYG{o}{:}\PYG{o}{:}\PYG{k+kr}{thread}\PYG{o}{\PYGZgt{}}\PYG{+w}{ }\PYG{n}{threads}\PYG{p}{;}
\PYG{k}{for}\PYG{p}{(}\PYG{k+kt}{int}\PYG{+w}{ }\PYG{n}{i}\PYG{+w}{ }\PYG{o}{=}\PYG{+w}{ }\PYG{l+m+mi}{0}\PYG{p}{;}\PYG{+w}{ }\PYG{n}{i}\PYG{+w}{ }\PYG{o}{\PYGZlt{}}\PYG{+w}{ }\PYG{l+m+mi}{10}\PYG{p}{;}\PYG{+w}{ }\PYG{o}{+}\PYG{o}{+}\PYG{n}{i}\PYG{p}{)}
\PYG{p}{\PYGZob{}}
\PYG{+w}{  }\PYG{n}{threads}\PYG{p}{.}\PYG{n}{emplace\PYGZus{}back}\PYG{p}{(}\PYG{p}{[}\PYG{o}{\PYGZam{}}\PYG{n}{counter}\PYG{p}{]}\PYG{p}{(}\PYG{p}{)}\PYG{p}{\PYGZob{}}
\PYG{+w}{    }\PYG{k}{for}\PYG{p}{(}\PYG{k+kt}{int}\PYG{+w}{ }\PYG{n}{i}\PYG{+w}{ }\PYG{o}{=}\PYG{+w}{ }\PYG{l+m+mi}{0}\PYG{p}{;}\PYG{+w}{ }\PYG{n}{i}\PYG{+w}{ }\PYG{o}{\PYGZlt{}}\PYG{+w}{ }\PYG{l+m+mi}{10}\PYG{p}{;}\PYG{+w}{ }\PYG{o}{+}\PYG{o}{+}\PYG{n}{i}\PYG{p}{)}
\PYG{+w}{      }\PYG{o}{+}\PYG{o}{+}\PYG{n}{counter}\PYG{p}{;}
\PYG{+w}{  }\PYG{p}{\PYGZcb{}}\PYG{p}{)}\PYG{p}{;}
\PYG{p}{\PYGZcb{}}
\PYG{k}{for}\PYG{p}{(}\PYG{k}{auto}\PYG{+w}{ }\PYG{o}{\PYGZam{}}\PYG{+w}{ }\PYG{n}{t}\PYG{+w}{ }\PYG{o}{:}\PYG{+w}{ }\PYG{n}{threads}\PYG{p}{)}\PYG{+w}{ }\PYG{n}{t}\PYG{p}{.}\PYG{n}{join}\PYG{p}{(}\PYG{p}{)}\PYG{p}{;}
\PYG{n}{std}\PYG{o}{:}\PYG{o}{:}\PYG{n}{cout}\PYG{+w}{ }\PYG{o}{\PYGZlt{}}\PYG{o}{\PYGZlt{}}\PYG{+w}{ }\PYG{n}{counter}\PYG{+w}{ }\PYG{o}{\PYGZlt{}}\PYG{o}{\PYGZlt{}}\PYG{+w}{ }\PYG{n}{std}\PYG{o}{:}\PYG{o}{:}\PYG{n}{endl}\PYG{p}{;}\PYG{+w}{ }\PYG{c+c1}{// 100}
\end{Verbatim}
\end{codeblockbg}


%%C080303 li

\item 使用\texttt{std::\allowbreak{}atomic\_\allowbreak{}flag}类表示原子布尔类型：


%%C080304 code

\begin{codeblockbg}
\begin{Verbatim}[commandchars=\\\{\},breaklines=true,breakanywhere=true,fontsize=\codefont,xleftmargin=17bp]
\PYG{n}{std}\PYG{o}{:}\PYG{o}{:}\PYG{n}{atomic\PYGZus{}flag}\PYG{+w}{ }\PYG{n}{lock}\PYG{+w}{ }\PYG{o}{=}\PYG{+w}{ }\PYG{n}{ATOMIC\PYGZus{}FLAG\PYGZus{}INIT}\PYG{p}{;}
\PYG{k+kt}{int}\PYG{+w}{ }\PYG{n}{counter}\PYG{+w}{ }\PYG{o}{=}\PYG{+w}{ }\PYG{l+m+mi}{0}\PYG{p}{;}
\PYG{n}{std}\PYG{o}{:}\PYG{o}{:}\PYG{n}{vector}\PYG{o}{\PYGZlt{}}\PYG{n}{std}\PYG{o}{:}\PYG{o}{:}\PYG{k+kr}{thread}\PYG{o}{\PYGZgt{}}\PYG{+w}{ }\PYG{n}{threads}\PYG{p}{;}
\PYG{k}{for}\PYG{p}{(}\PYG{k+kt}{int}\PYG{+w}{ }\PYG{n}{i}\PYG{+w}{ }\PYG{o}{=}\PYG{+w}{ }\PYG{l+m+mi}{0}\PYG{p}{;}\PYG{+w}{ }\PYG{n}{i}\PYG{+w}{ }\PYG{o}{\PYGZlt{}}\PYG{+w}{ }\PYG{l+m+mi}{10}\PYG{p}{;}\PYG{+w}{ }\PYG{o}{+}\PYG{o}{+}\PYG{n}{i}\PYG{p}{)}
\PYG{p}{\PYGZob{}}
\PYG{+w}{  }\PYG{n}{threads}\PYG{p}{.}\PYG{n}{emplace\PYGZus{}back}\PYG{p}{(}\PYG{p}{[}\PYG{o}{\PYGZam{}}\PYG{p}{]}\PYG{p}{(}\PYG{p}{)}\PYG{p}{\PYGZob{}}
\PYG{+w}{    }\PYG{k}{while}\PYG{p}{(}\PYG{n}{lock}\PYG{p}{.}\PYG{n}{test\PYGZus{}and\PYGZus{}set}\PYG{p}{(}\PYG{n}{std}\PYG{o}{:}\PYG{o}{:}\PYG{n}{memory\PYGZus{}order\PYGZus{}acquire}\PYG{p}{)}\PYG{p}{)}\PYG{p}{;}
\PYG{+w}{    }\PYG{o}{+}\PYG{o}{+}\PYG{n}{counter}\PYG{p}{;}
\PYG{+w}{    }\PYG{n}{lock}\PYG{p}{.}\PYG{n}{clear}\PYG{p}{(}\PYG{n}{std}\PYG{o}{:}\PYG{o}{:}\PYG{n}{memory\PYGZus{}order\PYGZus{}release}\PYG{p}{)}\PYG{p}{;}
\PYG{+w}{  }\PYG{p}{\PYGZcb{}}\PYG{p}{)}\PYG{p}{;}
\PYG{p}{\PYGZcb{}}
\PYG{k}{for}\PYG{p}{(}\PYG{k}{auto}\PYG{+w}{ }\PYG{o}{\PYGZam{}}\PYG{+w}{ }\PYG{n}{t}\PYG{+w}{ }\PYG{o}{:}\PYG{+w}{ }\PYG{n}{threads}\PYG{p}{)}\PYG{+w}{ }\PYG{n}{t}\PYG{p}{.}\PYG{n}{join}\PYG{p}{(}\PYG{p}{)}\PYG{p}{;}
\PYG{n}{std}\PYG{o}{:}\PYG{o}{:}\PYG{n}{cout}\PYG{+w}{ }\PYG{o}{\PYGZlt{}}\PYG{o}{\PYGZlt{}}\PYG{+w}{ }\PYG{n}{counter}\PYG{+w}{ }\PYG{o}{\PYGZlt{}}\PYG{o}{\PYGZlt{}}\PYG{+w}{ }\PYG{n}{std}\PYG{o}{:}\PYG{o}{:}\PYG{n}{endl}\PYG{p}{;}\PYG{+w}{ }\PYG{c+c1}{// 10}
\end{Verbatim}
\end{codeblockbg}


%%C080305 li

\item 使用原子类型的成员函数——\texttt{load()}、\texttt{store()}与\texttt{exchange()}——或者非成员函数——\texttt{atomic\_\allowbreak{}load()}/\myttbr{atomic\_\allowbreak{}load\_\allowbreak{}explicit()}、\texttt{atomic\_\allowbreak{}store()}/\myttbr{atomic\_\allowbreak{}store\_\allowbreak{}explicit()}与\texttt{atomic\_\allowbreak{}exchange()}/\myttbr{atomic\_\allowbreak{}exchange\_\allowbreak{}explicit()}——以原子方式读取、设置或交换原子对象的值。


%%C080306 li

\item 使用成员函数\texttt{fetch\_\allowbreak{}add()}与\texttt{fetch\_\allowbreak{}sub()}，或非成员函数\myttbr{atomic\_\allowbreak{}fetch\_\allowbreak{}add()}/\myttbr{atomic\_\allowbreak{}fetch\_\allowbreak{}add\_\allowbreak{}explicit()}与\myttbr{atomic\_\allowbreak{}fetch\_\allowbreak{}sub()}/\myttbr{atomic\_\allowbreak{}fetch\_\allowbreak{}sub\_\allowbreak{}explicit()}，以原子方式给原子对象加上或减去一个值，并返回运算前的值：


%%C080307 code

\begin{codeblockbg}
\begin{Verbatim}[commandchars=\\\{\},breaklines=true,breakanywhere=true,fontsize=\codefont,xleftmargin=17bp]
\PYG{n}{std}\PYG{o}{:}\PYG{o}{:}\PYG{n}{atomic}\PYG{o}{\PYGZlt{}}\PYG{k+kt}{int}\PYG{o}{\PYGZgt{}}\PYG{+w}{ }\PYG{n}{sum}\PYG{+w}{ }\PYG{p}{\PYGZob{}}\PYG{l+m+mi}{0}\PYG{p}{\PYGZcb{}}\PYG{p}{;}
\PYG{n}{std}\PYG{o}{:}\PYG{o}{:}\PYG{n}{vector}\PYG{o}{\PYGZlt{}}\PYG{k+kt}{int}\PYG{o}{\PYGZgt{}}\PYG{+w}{ }\PYG{n}{numbers}\PYG{+w}{ }\PYG{o}{=}\PYG{+w}{ }\PYG{n}{generate\PYGZus{}random}\PYG{p}{(}\PYG{p}{)}\PYG{p}{;}
\PYG{k+kt}{size\PYGZus{}t}\PYG{+w}{ }\PYG{n}{size}\PYG{+w}{ }\PYG{o}{=}\PYG{+w}{ }\PYG{n}{numbers}\PYG{p}{.}\PYG{n}{size}\PYG{p}{(}\PYG{p}{)}\PYG{p}{;}
\PYG{n}{std}\PYG{o}{:}\PYG{o}{:}\PYG{n}{vector}\PYG{o}{\PYGZlt{}}\PYG{n}{std}\PYG{o}{:}\PYG{o}{:}\PYG{k+kr}{thread}\PYG{o}{\PYGZgt{}}\PYG{+w}{ }\PYG{n}{threads}\PYG{p}{;}
\PYG{k}{for}\PYG{p}{(}\PYG{k+kt}{int}\PYG{+w}{ }\PYG{n}{i}\PYG{+w}{ }\PYG{o}{=}\PYG{+w}{ }\PYG{l+m+mi}{0}\PYG{p}{;}\PYG{+w}{ }\PYG{n}{i}\PYG{+w}{ }\PYG{o}{\PYGZlt{}}\PYG{+w}{ }\PYG{l+m+mi}{10}\PYG{p}{;}\PYG{+w}{ }\PYG{o}{+}\PYG{o}{+}\PYG{n}{i}\PYG{p}{)}
\PYG{p}{\PYGZob{}}
\PYG{+w}{  }\PYG{n}{threads}\PYG{p}{.}\PYG{n}{emplace\PYGZus{}back}\PYG{p}{(}\PYG{p}{[}\PYG{o}{\PYGZam{}}\PYG{n}{sum}\PYG{p}{,}\PYG{+w}{ }\PYG{o}{\PYGZam{}}\PYG{n}{numbers}\PYG{p}{]}\PYG{p}{(}\PYG{k+kt}{size\PYGZus{}t}\PYG{+w}{ }\PYG{k}{const}\PYG{+w}{ }\PYG{n}{start}\PYG{p}{,}
\PYG{+w}{                                        }\PYG{k+kt}{size\PYGZus{}t}\PYG{+w}{ }\PYG{k}{const}\PYG{+w}{ }\PYG{n}{end}\PYG{p}{)}\PYG{+w}{ }\PYG{p}{\PYGZob{}}
\PYG{+w}{    }\PYG{k}{for}\PYG{p}{(}\PYG{k+kt}{size\PYGZus{}t}\PYG{+w}{ }\PYG{n}{i}\PYG{+w}{ }\PYG{o}{=}\PYG{+w}{ }\PYG{n}{start}\PYG{p}{;}\PYG{+w}{ }\PYG{n}{i}\PYG{+w}{ }\PYG{o}{\PYGZlt{}}\PYG{+w}{ }\PYG{n}{end}\PYG{p}{;}\PYG{+w}{ }\PYG{o}{+}\PYG{o}{+}\PYG{n}{i}\PYG{p}{)}
\PYG{+w}{    }\PYG{p}{\PYGZob{}}
\PYG{+w}{      }\PYG{n}{std}\PYG{o}{:}\PYG{o}{:}\PYG{n}{atomic\PYGZus{}fetch\PYGZus{}add\PYGZus{}explicit}\PYG{p}{(}
\PYG{+w}{        }\PYG{o}{\PYGZam{}}\PYG{n}{sum}\PYG{p}{,}\PYG{+w}{ }\PYG{n}{numbers}\PYG{p}{[}\PYG{n}{i}\PYG{p}{]}\PYG{p}{,}
\PYG{+w}{        }\PYG{n}{std}\PYG{o}{:}\PYG{o}{:}\PYG{n}{memory\PYGZus{}order\PYGZus{}acquire}\PYG{p}{)}\PYG{p}{;}
\PYG{+w}{    }\PYG{c+c1}{// 等价于}
\PYG{+w}{    }\PYG{c+c1}{// sum.fetch\PYGZus{}add(numbers[i], std::memory\PYGZus{}order\PYGZus{}acquire);}
\PYG{+w}{  }\PYG{p}{\PYGZcb{}}\PYG{p}{\PYGZcb{}}\PYG{p}{,}
\PYG{+w}{  }\PYG{n}{i}\PYG{o}{*}\PYG{p}{(}\PYG{n}{size}\PYG{o}{/}\PYG{l+m+mi}{10}\PYG{p}{)}\PYG{p}{,}
\PYG{+w}{  }\PYG{p}{(}\PYG{n}{i}\PYG{o}{+}\PYG{l+m+mi}{1}\PYG{p}{)}\PYG{o}{*}\PYG{p}{(}\PYG{n}{size}\PYG{o}{/}\PYG{l+m+mi}{10}\PYG{p}{)}\PYG{p}{)}\PYG{p}{;}
\PYG{p}{\PYGZcb{}}
\PYG{k}{for}\PYG{p}{(}\PYG{k}{auto}\PYG{+w}{ }\PYG{o}{\PYGZam{}}\PYG{+w}{ }\PYG{n}{t}\PYG{+w}{ }\PYG{o}{:}\PYG{+w}{ }\PYG{n}{threads}\PYG{p}{)}\PYG{+w}{ }\PYG{n}{t}\PYG{p}{.}\PYG{n}{join}\PYG{p}{(}\PYG{p}{)}\PYG{p}{;}
\end{Verbatim}
\end{codeblockbg}


%%C080308 li

\item 使用成员函数\texttt{fetch\_\allowbreak{}and()}、\texttt{fetch\_\allowbreak{}or()}与\texttt{fetch\_\allowbreak{}xor()}，或非成员函数\myttbr{atomic\_\allowbreak{}fetch\_\allowbreak{}and()}/\myttbr{atomic\_\allowbreak{}fetch\_\allowbreak{}and\_\allowbreak{}explicit()}、\myttbr{atomic\_\allowbreak{}fetch\_\allowbreak{}or()}/\myttbr{atomic\_\allowbreak{}fetch\_\allowbreak{}or\_\allowbreak{}explicit()}与\myttbr{atomic\_\allowbreak{}fetch\_\allowbreak{}xor()}/\myttbr{atomic\_\allowbreak{}fetch\_\allowbreak{}xor\_\allowbreak{}explicit()}，用指定的实参分别执行按位与、按位或、按位异或原子操作，并返回运算前原子对象的值。


%%C080309 li

\item 使用\texttt{std::\allowbreak{}atomic\_\allowbreak{}flag}的成员函数\texttt{test\_\allowbreak{}and\_\allowbreak{}set()}与\texttt{clear()}，或非成员函数\myttbr{atomic\_\allowbreak{}flag\_\allowbreak{}test\_\allowbreak{}and\_\allowbreak{}set()}/\myttbr{atomic\_\allowbreak{}flag\_\allowbreak{}test\_\allowbreak{}and\_\allowbreak{}set\_\allowbreak{}explicit()}与\myttbr{atomic\_\allowbreak{}flag\_\allowbreak{}clear()}/\myttbr{atomic\_\allowbreak{}flag\_\allowbreak{}clear\_\allowbreak{}explicit()}，来设置或复位原子标志。

\end{itemize}


%%C080310 h3

\subsection{工作原理}


%%C080311 p

\texttt{std::\allowbreak{}atomic}是一个类模板，它（连同其各种特化）定义了原子类型。当一个线程写入对象、另一个线程读取数据，且不使用锁来保护访问时，原子类型对象的行为是有明确定义的。\texttt{std::\allowbreak{}atomic}类提供了若干特化：


%%C080312 li

\begin{itemize}
\item 针对\texttt{bool}的完全特化，并提供了名为\texttt{atomic\_\allowbreak{}bool}的typedef。


%%C080313 li

\item 针对所有整数类型的完全特化，并提供了\texttt{atomic\_\allowbreak{}int}、\texttt{atomic\_\allowbreak{}long}、\texttt{atomic\_\allowbreak{}char}、\texttt{atomic\_\allowbreak{}wchar}等众多typedef。


%%C080314 li

\item 针对指针类型的偏特化。

\end{itemize}


%%C080315 p

\texttt{atomic}类模板提供了多种执行原子操作的成员函数，如下所示：


%%C080316 li

\begin{itemize}
\item \texttt{load()}，以原子方式加载并返回对象的值。


%%C080317 li

\item \texttt{store()}，以原子方式把非原子值存入对象；该函数不返回任何内容。


%%C080318 li

\item \texttt{exchange()}，以原子方式把非原子值存入对象，并返回此前的值。


%%C080319 li

\item \texttt{operator=}，效果与\texttt{store(\allowbreak{}arg)}相同。


%%C080320 li

\item \texttt{fetch\_\allowbreak{}add()}，以原子方式把一个非原子实参加到原子值上，并返回此前存储的值。


%%C080321 li

\item \texttt{fetch\_\allowbreak{}sub()}，以原子方式从原子值中减去一个非原子实参，并返回此前存储的值。


%%C080322 li

\item \texttt{fetch\_\allowbreak{}and()}、\texttt{fetch\_\allowbreak{}or()}与\texttt{fetch\_\allowbreak{}xor()}，以原子方式在实参与原子值之间执行按位与、按位或、按位异或操作，把新值存入原子对象，并返回此前的值。


%%C080323 li

\item 前置与后置的\texttt{operator++}与\texttt{operator--}，以原子方式把原子对象的值加1或减1。这些操作等价于使用\texttt{fetch\_\allowbreak{}add()}或\texttt{fetch\_\allowbreak{}sub()}。


%%C080324 li

\item \texttt{operator +=}、\texttt{-\allowbreak{}=}、\texttt{\&=}、\texttt{|=}与\texttt{ˆ=}，在实参与原子值之间执行加法、减法或按位与、按位或、按位异或操作，并把新值存入原子对象。这些操作等价于使用\texttt{fetch\_\allowbreak{}add()}、\texttt{fetch\_\allowbreak{}sub()}、\texttt{fetch\_\allowbreak{}and()}、\texttt{fetch\_\allowbreak{}or()}与\texttt{fetch\_\allowbreak{}xor()}。

\end{itemize}


%%C080325 p

假设有一个原子变量，例如\texttt{std::\allowbreak{}atomic<\allowbreak{}int> a}，那么下面的操作并不是原子操作：


%%C080326 code

\begin{codeblockbg}
\begin{Verbatim}[commandchars=\\\{\},breaklines=true,breakanywhere=true,fontsize=\codefont,xleftmargin=17bp]
\PYG{n}{a}\PYG{+w}{ }\PYG{o}{=}\PYG{+w}{ }\PYG{n}{a}\PYG{+w}{ }\PYG{o}{+}\PYG{+w}{ }\PYG{l+m+mi}{42}\PYG{p}{;}
\end{Verbatim}
\end{codeblockbg}


%%C080327 p

这实际上是一系列操作，其中只有一部分是原子的：


%%C080328 li

\begin{itemize}
\item 以原子方式加载原子对象的值。


%%C080329 li

\item 在加载得到的值上加42。


%%C080330 li

\item 以原子方式把结果存入原子对象\texttt{a}。

\end{itemize}


%%C080331 p

另一方面，下面的操作使用了成员\texttt{operator +=}，因此是原子的：


%%C080332 code

\begin{codeblockbg}
\begin{Verbatim}[commandchars=\\\{\},breaklines=true,breakanywhere=true,fontsize=\codefont,xleftmargin=17bp]
\PYG{n}{a}\PYG{+w}{ }\PYG{o}{+}\PYG{o}{=}\PYG{+w}{ }\PYG{l+m+mi}{42}\PYG{p}{;}
\end{Verbatim}
\end{codeblockbg}


%%C080333 p

该操作与下面任意一种写法效果相同：


%%C080334 code

\begin{codeblockbg}
\begin{Verbatim}[commandchars=\\\{\},breaklines=true,breakanywhere=true,fontsize=\codefont,xleftmargin=17bp]
\PYG{n}{a}\PYG{p}{.}\PYG{n}{fetch\PYGZus{}add}\PYG{p}{(}\PYG{l+m+mi}{42}\PYG{p}{)}\PYG{p}{;}\PYG{+w}{               }\PYG{c+c1}{// 使用成员函数}
\PYG{n}{std}\PYG{o}{:}\PYG{o}{:}\PYG{n}{atomic\PYGZus{}fetch\PYGZus{}add}\PYG{p}{(}\PYG{o}{\PYGZam{}}\PYG{n}{a}\PYG{p}{,}\PYG{+w}{ }\PYG{l+m+mi}{42}\PYG{p}{)}\PYG{p}{;}\PYG{+w}{ }\PYG{c+c1}{// 使用非成员函数}
\end{Verbatim}
\end{codeblockbg}


%%C080335 p

虽然\texttt{std::\allowbreak{}atomic}对\texttt{bool}类型有完全特化，名为\texttt{std::\allowbreak{}atomic<\allowbreak{}bool>}，但标准还定义了另一种原子类型\texttt{std::\allowbreak{}atomic\_\allowbreak{}flag}，它保证是无锁的。不过，这种原子类型与\texttt{std::\allowbreak{}atomic\_\allowbreak{}bool}差别很大，它只有两个成员函数：


%%C080336 li

\begin{itemize}
\item \texttt{test\_\allowbreak{}and\_\allowbreak{}set()}，以原子方式把值设为\texttt{true}，并返回此前的值。


%%C080337 li

\item \texttt{clear()}，以原子方式把值设为\texttt{false}。

\end{itemize}


%%C080338 p

前面提到的所有成员函数，无论属于\texttt{std::\allowbreak{}atomic}还是\texttt{std::\allowbreak{}atomic\_\allowbreak{}flag}，都有对应的非成员版本，其名称分别以\texttt{atomic\_}或\texttt{atomic\_\allowbreak{}flag\_}为前缀，具体取决于所引用的类型。例如，\myttbr{std::\allowbreak{}atomic::\allowbreak{}fetch\_\allowbreak{}add()}的对应物是\myttbr{std::\allowbreak{}atomic\_\allowbreak{}fetch\_\allowbreak{}add()}，这些非成员函数的第一个实参始终是指向\texttt{std::\allowbreak{}atomic}对象的指针。在内部，非成员函数会在传入的\texttt{std::\allowbreak{}atomic}实参上调用对应的成员函数。类似地，\myttbr{std::\allowbreak{}atomic\_\allowbreak{}flag::\allowbreak{}test\_\allowbreak{}and\_\allowbreak{}set()}的对应物是\myttbr{std::\allowbreak{}atomic\_\allowbreak{}flag\_\allowbreak{}test\_\allowbreak{}and\_\allowbreak{}set()}，它的第一个参数是指向\texttt{std::\allowbreak{}atomic\_\allowbreak{}flag}对象的指针。


%%C080339 p

所有这些\texttt{std::\allowbreak{}atomic}与\texttt{std::\allowbreak{}atomic\_\allowbreak{}flag}的成员函数都有两组重载；其中一组多出一个表示内存序的实参。类似地，所有非成员函数——例如\myttbr{std::\allowbreak{}atomic\_\allowbreak{}load()}、\myttbr{std::\allowbreak{}atomic\_\allowbreak{}fetch\_\allowbreak{}add()}与\myttbr{std::\allowbreak{}atomic\_\allowbreak{}flag\_\allowbreak{}test\_\allowbreak{}and\_\allowbreak{}set()}——都有一个带\texttt{\_\allowbreak{}explicit}后缀的伴随版本，即\myttbr{std::\allowbreak{}atomic\_\allowbreak{}load\_\allowbreak{}explicit()}、\myttbr{std::\allowbreak{}atomic\_\allowbreak{}fetch\_\allowbreak{}add\_\allowbreak{}explicit()}与\myttbr{std::\allowbreak{}atomic\_\allowbreak{}flag\_\allowbreak{}test\_\allowbreak{}and\_\allowbreak{}set\_\allowbreak{}explicit()}；这些函数多出一个表示内存序的实参。


%%C080340 p

内存序规定了非原子内存访问如何在原子操作周围排序。默认情况下，所有原子类型与原子操作的内存序都是顺序一致性（sequential consistency）。其他顺序类型定义在\texttt{std::\allowbreak{}memory\_\allowbreak{}order}枚举中，可以作为实参传给\texttt{std::\allowbreak{}atomic}与\texttt{std::\allowbreak{}atomic\_\allowbreak{}flag}的成员函数，或传给带\texttt{\_\allowbreak{}explicit()}后缀的非成员函数。


%%C080341 note

\begin{myNote}
\textit{顺序一致性}是一种一致性模型，它要求在多处理器系统中，所有指令都按某种顺序执行，并且所有写入都会立即在整个系统中可见。该模型由Leslie Lamport在70年代首次提出，其描述如下：\textit{任何一次执行的结果，都与所有处理器的操作按某种顺序执行时相同，}\textit{并且每个处理器的操作都按各自程序规定的顺序}\textit{出现在该序列中。}
\end{myNote}


%%C080342 p

下表列出了各种内存序函数，它们取自C++参考网站（http://\allowbreak{}en.\allowbreak{}cppreference.\allowbreak{}com/\allowbreak{}w/\allowbreak{}cpp/\allowbreak{}atomic/\allowbreak{}memory\_\allowbreak{}order）。每一种内存序的具体工作机制都超出了本书的范围，可以查阅标准的C++参考（即我们刚刚提到的链接）：


%%C080343 table

\myTableBegin
\begin{xltabular}{\textwidth}{|>{\hsize=0.296\hsize\raggedright\arraybackslash}X|>{\hsize=0.704\hsize\raggedright\arraybackslash}X|}
\hline
\rowcolor{tabhead} \textcolor{white}{\textbf{模型}} & \textcolor{white}{\textbf{说明}} \\
\hline
\endfirsthead
\rowcolor{tabhead} \textcolor{white}{\textbf{模型}} & \textcolor{white}{\textbf{说明}} \\
\hline
\endhead
\myttbr{memory\_\allowbreak{}order\_\allowbreak{}relaxed} & 这是一种松散操作。没有任何同步或顺序约束；该操作只要求原子性。\\
\hline
\myttbr{memory\_\allowbreak{}order\_\allowbreak{}consume} & 使用该内存序的加载操作会在受影响的内存位置上执行consume操作；当前线程中依赖于刚加载值的读或写，都不能被重排到该加载操作之前。其他线程中释放同一原子变量时对数据依赖变量所做的写入，在当前线程中可见。在大多数平台上，这只影响编译器优化。\\
\hline
\myttbr{memory\_\allowbreak{}order\_\allowbreak{}acquire} & 使用该内存序的加载操作会在受影响的内存位置上执行acquire操作；当前线程中的任何读或写都不能被重排到该加载之前。其他线程中释放同一原子变量的所有写入，在当前线程中可见。\\
\hline
\myttbr{memory\_\allowbreak{}order\_\allowbreak{}release} & 使用该内存序的存储操作会执行release操作；当前线程中的任何读或写都不能被重排到该存储之后。当前线程中的所有写入，对其他获取同一原子变量的线程可见；而携带对该原子变量依赖的写入，会对消费同一原子变量的其他线程可见。\\
\hline
\myttbr{memory\_\allowbreak{}order\_\allowbreak{}acq\_\allowbreak{}rel} & 使用该内存序的读-改-写操作既是acquire操作也是release操作。当前线程中的任何内存读或写都不能被重排到该存储之前或之后。其他线程中释放同一原子变量的所有写入在该修改之前可见，而该修改对其他获取同一原子变量的线程可见。\\
\hline
\myttbr{memory\_\allowbreak{}order\_\allowbreak{}seq\_\allowbreak{}cst} & 使用该内存序的任何操作既是acquire操作也是release操作；存在唯一一个全序，所有线程在其中以相同顺序观察到所有修改。\\
\hline
\end{xltabular}
\myTableEnd


%%C080344 table

% 原书此处是上一张表的续表（PDF分页把表切断，抽取时首行被误当表头），
% 已并入上表，本单元留空。


%%C080345 p

\textit{具体做法}一节中的第一个示例演示了多个线程并发地反复递增计数器。这个示例还可以进一步改进：实现一个类来表示原子计数器，其中提供\texttt{increment()}与\texttt{decrement()}这样修改计数值的方法，以及\texttt{get()}这样获取当前值的方法：


%%C080346 code

\begin{codeblockbg}
\begin{Verbatim}[commandchars=\\\{\},breaklines=true,breakanywhere=true,fontsize=\codefont,xleftmargin=17bp]
\PYG{k}{template}
\PYG{o}{\PYGZlt{}}\PYG{k}{typename}\PYG{+w}{ }\PYG{n+nc}{T}\PYG{p}{,}
\PYG{+w}{ }\PYG{k}{typename}\PYG{+w}{ }\PYG{n+nc}{I}\PYG{+w}{ }\PYG{o}{=}\PYG{+w}{ }\PYG{k}{typename}\PYG{+w}{ }\PYG{n+nc}{std}\PYG{o}{:}\PYG{o}{:}\PYG{n}{enable\PYGZus{}if}\PYG{o}{\PYGZlt{}}\PYG{n}{std}\PYG{o}{:}\PYG{o}{:}\PYG{n}{is\PYGZus{}integral}\PYG{o}{\PYGZlt{}}\PYG{n}{T}\PYG{o}{\PYGZgt{}}
\PYG{+w}{    }\PYG{o}{:}\PYG{o}{:}\PYG{n}{value}\PYG{o}{\PYGZgt{}}\PYG{o}{:}\PYG{o}{:}\PYG{n}{type}\PYG{o}{\PYGZgt{}}
\PYG{k}{class}\PYG{+w}{ }\PYG{n+nc}{atomic\PYGZus{}counter}
\PYG{p}{\PYGZob{}}
\PYG{+w}{  }\PYG{n}{std}\PYG{o}{:}\PYG{o}{:}\PYG{n}{atomic}\PYG{o}{\PYGZlt{}}\PYG{n}{T}\PYG{o}{\PYGZgt{}}\PYG{+w}{ }\PYG{n}{counter}\PYG{+w}{ }\PYG{p}{\PYGZob{}}\PYG{l+m+mi}{0}\PYG{p}{\PYGZcb{}}\PYG{p}{;}
\PYG{+w}{  }\PYG{k}{public}\PYG{o}{:}
\PYG{+w}{  }\PYG{n}{T}\PYG{+w}{ }\PYG{n}{increment}\PYG{p}{(}\PYG{p}{)}
\PYG{+w}{  }\PYG{p}{\PYGZob{}}
\PYG{+w}{    }\PYG{k}{return}\PYG{+w}{ }\PYG{n}{counter}\PYG{p}{.}\PYG{n}{fetch\PYGZus{}add}\PYG{p}{(}\PYG{l+m+mi}{1}\PYG{p}{)}\PYG{p}{;}
\PYG{+w}{  }\PYG{p}{\PYGZcb{}}
\PYG{n}{T}\PYG{+w}{ }\PYG{n}{decrement}\PYG{p}{(}\PYG{p}{)}
\PYG{p}{\PYGZob{}}
\PYG{+w}{  }\PYG{k}{return}\PYG{+w}{ }\PYG{n}{counter}\PYG{p}{.}\PYG{n}{fetch\PYGZus{}sub}\PYG{p}{(}\PYG{l+m+mi}{1}\PYG{p}{)}\PYG{p}{;}
\PYG{p}{\PYGZcb{}}
\PYG{+w}{  }\PYG{n}{T}\PYG{+w}{ }\PYG{n}{get}\PYG{p}{(}\PYG{p}{)}
\PYG{+w}{  }\PYG{p}{\PYGZob{}}
\PYG{+w}{    }\PYG{k}{return}\PYG{+w}{ }\PYG{n}{counter}\PYG{p}{.}\PYG{n}{load}\PYG{p}{(}\PYG{p}{)}\PYG{p}{;}
\PYG{+w}{  }\PYG{p}{\PYGZcb{}}
\PYG{p}{\PYGZcb{}}\PYG{p}{;}
\end{Verbatim}
\end{codeblockbg}


%%C080347 p

借助这个类模板，第一个示例可以改写为下面的形式，结果完全相同：


%%C080348 code

\begin{codeblockbg}
\begin{Verbatim}[commandchars=\\\{\},breaklines=true,breakanywhere=true,fontsize=\codefont,xleftmargin=17bp]
\PYG{n}{atomic\PYGZus{}counter}\PYG{o}{\PYGZlt{}}\PYG{k+kt}{int}\PYG{o}{\PYGZgt{}}\PYG{+w}{ }\PYG{n}{counter}\PYG{p}{;}
\PYG{n}{std}\PYG{o}{:}\PYG{o}{:}\PYG{n}{vector}\PYG{o}{\PYGZlt{}}\PYG{n}{std}\PYG{o}{:}\PYG{o}{:}\PYG{k+kr}{thread}\PYG{o}{\PYGZgt{}}\PYG{+w}{ }\PYG{n}{threads}\PYG{p}{;}
\PYG{k}{for}\PYG{p}{(}\PYG{k+kt}{int}\PYG{+w}{ }\PYG{n}{i}\PYG{+w}{ }\PYG{o}{=}\PYG{+w}{ }\PYG{l+m+mi}{0}\PYG{p}{;}\PYG{+w}{ }\PYG{n}{i}\PYG{+w}{ }\PYG{o}{\PYGZlt{}}\PYG{+w}{ }\PYG{l+m+mi}{10}\PYG{p}{;}\PYG{+w}{ }\PYG{o}{+}\PYG{o}{+}\PYG{n}{i}\PYG{p}{)}
\PYG{p}{\PYGZob{}}
\PYG{+w}{  }\PYG{n}{threads}\PYG{p}{.}\PYG{n}{emplace\PYGZus{}back}\PYG{p}{(}\PYG{p}{[}\PYG{o}{\PYGZam{}}\PYG{n}{counter}\PYG{p}{]}\PYG{p}{(}\PYG{p}{)}\PYG{p}{\PYGZob{}}
\PYG{+w}{    }\PYG{k}{for}\PYG{p}{(}\PYG{k+kt}{int}\PYG{+w}{ }\PYG{n}{i}\PYG{+w}{ }\PYG{o}{=}\PYG{+w}{ }\PYG{l+m+mi}{0}\PYG{p}{;}\PYG{+w}{ }\PYG{n}{i}\PYG{+w}{ }\PYG{o}{\PYGZlt{}}\PYG{+w}{ }\PYG{l+m+mi}{10}\PYG{p}{;}\PYG{+w}{ }\PYG{o}{+}\PYG{o}{+}\PYG{n}{i}\PYG{p}{)}
\PYG{+w}{      }\PYG{n}{counter}\PYG{p}{.}\PYG{n}{increment}\PYG{p}{(}\PYG{p}{)}\PYG{p}{;}
\PYG{+w}{  }\PYG{p}{\PYGZcb{}}\PYG{p}{)}\PYG{p}{;}
\PYG{p}{\PYGZcb{}}
\PYG{k}{for}\PYG{p}{(}\PYG{k}{auto}\PYG{+w}{ }\PYG{o}{\PYGZam{}}\PYG{+w}{ }\PYG{n}{t}\PYG{+w}{ }\PYG{o}{:}\PYG{+w}{ }\PYG{n}{threads}\PYG{p}{)}\PYG{+w}{ }\PYG{n}{t}\PYG{p}{.}\PYG{n}{join}\PYG{p}{(}\PYG{p}{)}\PYG{p}{;}
\PYG{n}{std}\PYG{o}{:}\PYG{o}{:}\PYG{n}{cout}\PYG{+w}{ }\PYG{o}{\PYGZlt{}}\PYG{o}{\PYGZlt{}}\PYG{+w}{ }\PYG{n}{counter}\PYG{p}{.}\PYG{n}{get}\PYG{p}{(}\PYG{p}{)}\PYG{+w}{ }\PYG{o}{\PYGZlt{}}\PYG{o}{\PYGZlt{}}\PYG{+w}{ }\PYG{n}{std}\PYG{o}{:}\PYG{o}{:}\PYG{n}{endl}\PYG{p}{;}\PYG{+w}{ }\PYG{c+c1}{// 100}
\end{Verbatim}
\end{codeblockbg}


%%C080349 h3

\subsection{延伸阅读}


%%C080350 li

\begin{itemize}
\item \textit{使用线程}


%%C080351 li

\item \textit{用互斥量与锁同步对共享数据的访问}


%%C080352 li

\item \textit{异步执行函数}

\end{itemize}


%%C080353 h2

\section{用线程实现并行map与fold}


%%C080354 p

在\texttt{Chapter 3}（\textit{函数探究}）中，我们讨论了两种高阶函数：\texttt{map}，它把函数应用到范围中的元素上，或变换该范围，或产生一个新范围；以及\texttt{fold}，它把范围中的元素合并为单个值。我们给出的各种实现都是串行的。不过在并发、线程与异步任务的语境下，我们可以利用硬件运行这些函数的并行版本，以便在大范围上、或者在变换与聚合比较耗时时加快执行速度。本案例将介绍一种用线程实现\texttt{map}与\texttt{fold}的可行方案。


%%C080355 h3

\subsection{准备工作}


%%C080356 p

你需要熟悉\texttt{map}与\texttt{fold}函数的概念。建议先阅读\texttt{Chapter 3}（\textit{函数探究}）中的\textit{实现高阶函数map与fold}案例。本案例将使用\textit{使用线程}案例中介绍的各种线程功能。为了测量这些函数的执行时间并与串行版本比较，我们将使用\texttt{perf\_\allowbreak{}timer}类模板，它是在\texttt{Chapter 6}（\textit{通用工具}）的\textit{用标准时钟测量函数执行时间}案例中介绍的。


%%C080357 note

\begin{myNote}
算法的并行版本有可能缩短执行时间，但这并非在所有情况下都成立。线程的上下文切换以及对共享数据的同步访问都会带来可观的开销。对某些实现和特定数据集而言，这种开销甚至会让并行版本的执行时间超过串行版本。
\end{myNote}


%%C080358 p

要确定拆分工作需要多少个线程，可以使用下面的函数：


%%C080359 code

\begin{codeblockbg}
\begin{Verbatim}[commandchars=\\\{\},breaklines=true,breakanywhere=true,fontsize=\codefont,xleftmargin=17bp]
\PYG{k+kt}{unsigned}\PYG{+w}{ }\PYG{n+nf}{get\PYGZus{}no\PYGZus{}of\PYGZus{}threads}\PYG{p}{(}\PYG{p}{)}
\PYG{p}{\PYGZob{}}
\PYG{+w}{  }\PYG{k}{return}\PYG{+w}{ }\PYG{n}{std}\PYG{o}{:}\PYG{o}{:}\PYG{k+kr}{thread}\PYG{o}{:}\PYG{o}{:}\PYG{n}{hardware\PYGZus{}concurrency}\PYG{p}{(}\PYG{p}{)}\PYG{p}{;}
\PYG{p}{\PYGZcb{}}
\end{Verbatim}
\end{codeblockbg}


%%C080360 h3

\subsection{具体做法}


%%C080361 p

要实现\texttt{map}函数的并行版本，请按以下步骤操作：


%%C080362 p

定义一个函数模板，它接受指向某个范围的\texttt{begin}与\texttt{end}迭代器，以及


%%C080363 p

一个要应用到所有元素上的函数：


%%C080364 code

\begin{codeblockbg}
\begin{Verbatim}[commandchars=\\\{\},breaklines=true,breakanywhere=true,fontsize=\codefont,xleftmargin=17bp]
\PYG{k}{template}\PYG{+w}{ }\PYG{o}{\PYGZlt{}}\PYG{k}{typename}\PYG{+w}{ }\PYG{n+nc}{Iter}\PYG{p}{,}\PYG{+w}{ }\PYG{k}{typename}\PYG{+w}{ }\PYG{n+nc}{F}\PYG{o}{\PYGZgt{}}
\PYG{k+kt}{void}\PYG{+w}{ }\PYG{n}{parallel\PYGZus{}map}\PYG{p}{(}\PYG{n}{Iter}\PYG{+w}{ }\PYG{n}{begin}\PYG{p}{,}\PYG{+w}{ }\PYG{n}{Iter}\PYG{+w}{ }\PYG{n}{end}\PYG{p}{,}\PYG{+w}{ }\PYG{n}{F}\PYG{+w}{ }\PYG{n}{f}\PYG{p}{)}
\PYG{p}{\PYGZob{}}
\PYG{p}{\PYGZcb{}}
\end{Verbatim}
\end{codeblockbg}


%%C080365 p

检查范围的大小。如果元素个数小于某个


%%C080366 p

预先设定的阈值（本实现中该阈值为10000），就以串行方式执行映射：


%%C080367 code

\begin{codeblockbg}
\begin{Verbatim}[commandchars=\\\{\},breaklines=true,breakanywhere=true,fontsize=\codefont,xleftmargin=17bp]
\PYG{k}{auto}\PYG{+w}{ }\PYG{n}{size}\PYG{+w}{ }\PYG{o}{=}\PYG{+w}{ }\PYG{n}{std}\PYG{o}{:}\PYG{o}{:}\PYG{n}{distance}\PYG{p}{(}\PYG{n}{begin}\PYG{p}{,}\PYG{+w}{ }\PYG{n}{end}\PYG{p}{)}\PYG{p}{;}
\PYG{k}{if}\PYG{p}{(}\PYG{n}{size}\PYG{+w}{ }\PYG{o}{\PYGZlt{}}\PYG{o}{=}\PYG{+w}{ }\PYG{l+m+mi}{10000}\PYG{p}{)}
\PYG{+w}{  }\PYG{n}{std}\PYG{o}{:}\PYG{o}{:}\PYG{n}{transform}\PYG{p}{(}\PYG{n}{begin}\PYG{p}{,}\PYG{+w}{ }\PYG{n}{end}\PYG{p}{,}\PYG{+w}{ }\PYG{n}{begin}\PYG{p}{,}\PYG{+w}{ }\PYG{n}{std}\PYG{o}{:}\PYG{o}{:}\PYG{n}{forward}\PYG{o}{\PYGZlt{}}\PYG{n}{F}\PYG{o}{\PYGZgt{}}\PYG{p}{(}\PYG{n}{f}\PYG{p}{)}\PYG{p}{)}\PYG{p}{;}
\end{Verbatim}
\end{codeblockbg}


%%C080368 p

对于更大的范围，把工作拆分到多个线程上，让每个线程映射范围的


%%C080369 p

一部分。这些部分不应重叠，以避免需要同步对共享数据的访问：


%%C080370 code

\begin{codeblockbg}
\begin{Verbatim}[commandchars=\\\{\},breaklines=true,breakanywhere=true,fontsize=\codefont,xleftmargin=17bp]
\PYG{k}{else}
\PYG{p}{\PYGZob{}}
\PYG{+w}{  }\PYG{k}{auto}\PYG{+w}{ }\PYG{n}{no\PYGZus{}of\PYGZus{}threads}\PYG{+w}{ }\PYG{o}{=}\PYG{+w}{ }\PYG{n}{get\PYGZus{}no\PYGZus{}of\PYGZus{}threads}\PYG{p}{(}\PYG{p}{)}\PYG{p}{;}
\PYG{+w}{  }\PYG{k}{auto}\PYG{+w}{ }\PYG{n}{part}\PYG{+w}{ }\PYG{o}{=}\PYG{+w}{ }\PYG{n}{size}\PYG{+w}{ }\PYG{o}{/}\PYG{+w}{ }\PYG{n}{no\PYGZus{}of\PYGZus{}threads}\PYG{p}{;}
\PYG{+w}{  }\PYG{k}{auto}\PYG{+w}{ }\PYG{n}{last}\PYG{+w}{ }\PYG{o}{=}\PYG{+w}{ }\PYG{n}{begin}\PYG{p}{;}
\PYG{+w}{  }\PYG{c+c1}{// 接第4步与第5步}
\PYG{p}{\PYGZcb{}}
\end{Verbatim}
\end{codeblockbg}


%%C080371 p

启动这些线程，并在每个线程上运行映射的串行版本：


%%C080372 code

\begin{codeblockbg}
\begin{Verbatim}[commandchars=\\\{\},breaklines=true,breakanywhere=true,fontsize=\codefont,xleftmargin=17bp]
\PYG{n}{std}\PYG{o}{:}\PYG{o}{:}\PYG{n}{vector}\PYG{o}{\PYGZlt{}}\PYG{n}{std}\PYG{o}{:}\PYG{o}{:}\PYG{k+kr}{thread}\PYG{o}{\PYGZgt{}}\PYG{+w}{ }\PYG{n}{threads}\PYG{p}{;}
\PYG{k}{for}\PYG{p}{(}\PYG{k+kt}{unsigned}\PYG{+w}{ }\PYG{n}{i}\PYG{+w}{ }\PYG{o}{=}\PYG{+w}{ }\PYG{l+m+mi}{0}\PYG{p}{;}\PYG{+w}{ }\PYG{n}{i}\PYG{+w}{ }\PYG{o}{\PYGZlt{}}\PYG{+w}{ }\PYG{n}{no\PYGZus{}of\PYGZus{}threads}\PYG{p}{;}\PYG{+w}{ }\PYG{o}{+}\PYG{o}{+}\PYG{n}{i}\PYG{p}{)}
\PYG{p}{\PYGZob{}}
\PYG{+w}{  }\PYG{k}{if}\PYG{p}{(}\PYG{n}{i}\PYG{+w}{ }\PYG{o}{=}\PYG{o}{=}\PYG{+w}{ }\PYG{n}{no\PYGZus{}of\PYGZus{}threads}\PYG{+w}{ }\PYG{o}{\PYGZhy{}}\PYG{+w}{ }\PYG{l+m+mi}{1}\PYG{p}{)}\PYG{+w}{ }\PYG{n}{last}\PYG{+w}{ }\PYG{o}{=}\PYG{+w}{ }\PYG{n}{end}\PYG{p}{;}
\PYG{+w}{  }\PYG{k}{else}\PYG{+w}{ }\PYG{n}{std}\PYG{o}{:}\PYG{o}{:}\PYG{n}{advance}\PYG{p}{(}\PYG{n}{last}\PYG{p}{,}\PYG{+w}{ }\PYG{n}{part}\PYG{p}{)}\PYG{p}{;}
\PYG{+w}{  }\PYG{n}{threads}\PYG{p}{.}\PYG{n}{emplace\PYGZus{}back}\PYG{p}{(}
\PYG{+w}{    }\PYG{p}{[}\PYG{o}{=}\PYG{p}{,}\PYG{o}{\PYGZam{}}\PYG{n}{f}\PYG{p}{]}\PYG{p}{\PYGZob{}}\PYG{n}{std}\PYG{o}{:}\PYG{o}{:}\PYG{n}{transform}\PYG{p}{(}\PYG{n}{begin}\PYG{p}{,}\PYG{+w}{ }\PYG{n}{last}\PYG{p}{,}
\PYG{+w}{                          }\PYG{n}{begin}\PYG{p}{,}\PYG{+w}{ }\PYG{n}{std}\PYG{o}{:}\PYG{o}{:}\PYG{n}{forward}\PYG{o}{\PYGZlt{}}\PYG{n}{F}\PYG{o}{\PYGZgt{}}\PYG{p}{(}\PYG{n}{f}\PYG{p}{)}\PYG{p}{)}\PYG{p}{;}\PYG{p}{\PYGZcb{}}\PYG{p}{)}\PYG{p}{;}
\PYG{+w}{  }\PYG{n}{begin}\PYG{+w}{ }\PYG{o}{=}\PYG{+w}{ }\PYG{n}{last}\PYG{p}{;}
\PYG{p}{\PYGZcb{}}
\end{Verbatim}
\end{codeblockbg}


%%C080373 p

等待所有线程执行完毕：


%%C080374 code

\begin{codeblockbg}
\begin{Verbatim}[commandchars=\\\{\},breaklines=true,breakanywhere=true,fontsize=\codefont,xleftmargin=17bp]
\PYG{k}{for}\PYG{p}{(}\PYG{k}{auto}\PYG{+w}{ }\PYG{o}{\PYGZam{}}\PYG{+w}{ }\PYG{n}{t}\PYG{+w}{ }\PYG{o}{:}\PYG{+w}{ }\PYG{n}{threads}\PYG{p}{)}\PYG{+w}{ }\PYG{n}{t}\PYG{p}{.}\PYG{n}{join}\PYG{p}{(}\PYG{p}{)}\PYG{p}{;}
\end{Verbatim}
\end{codeblockbg}


%%C080375 p

把前面的步骤组合起来，就得到下面的实现：


%%C080376 code

\begin{codeblockbg}
\begin{Verbatim}[commandchars=\\\{\},breaklines=true,breakanywhere=true,fontsize=\codefont,xleftmargin=17bp]
\PYG{k}{template}\PYG{+w}{ }\PYG{o}{\PYGZlt{}}\PYG{k}{typename}\PYG{+w}{ }\PYG{n+nc}{Iter}\PYG{p}{,}\PYG{+w}{ }\PYG{k}{typename}\PYG{+w}{ }\PYG{n+nc}{F}\PYG{o}{\PYGZgt{}}
\PYG{k+kt}{void}\PYG{+w}{ }\PYG{n}{parallel\PYGZus{}map}\PYG{p}{(}\PYG{n}{Iter}\PYG{+w}{ }\PYG{n}{begin}\PYG{p}{,}\PYG{+w}{ }\PYG{n}{Iter}\PYG{+w}{ }\PYG{n}{end}\PYG{p}{,}\PYG{+w}{ }\PYG{n}{F}\PYG{+w}{ }\PYG{n}{f}\PYG{p}{)}
\PYG{p}{\PYGZob{}}
\PYG{+w}{  }\PYG{k}{auto}\PYG{+w}{ }\PYG{n}{size}\PYG{+w}{ }\PYG{o}{=}\PYG{+w}{ }\PYG{n}{std}\PYG{o}{:}\PYG{o}{:}\PYG{n}{distance}\PYG{p}{(}\PYG{n}{begin}\PYG{p}{,}\PYG{+w}{ }\PYG{n}{end}\PYG{p}{)}\PYG{p}{;}
\PYG{+w}{  }\PYG{k}{if}\PYG{p}{(}\PYG{n}{size}\PYG{+w}{ }\PYG{o}{\PYGZlt{}}\PYG{o}{=}\PYG{+w}{ }\PYG{l+m+mi}{10000}\PYG{p}{)}
\PYG{+w}{    }\PYG{n}{std}\PYG{o}{:}\PYG{o}{:}\PYG{n}{transform}\PYG{p}{(}\PYG{n}{begin}\PYG{p}{,}\PYG{+w}{ }\PYG{n}{end}\PYG{p}{,}\PYG{+w}{ }\PYG{n}{begin}\PYG{p}{,}\PYG{+w}{ }\PYG{n}{std}\PYG{o}{:}\PYG{o}{:}\PYG{n}{forward}\PYG{o}{\PYGZlt{}}\PYG{n}{F}\PYG{o}{\PYGZgt{}}\PYG{p}{(}\PYG{n}{f}\PYG{p}{)}\PYG{p}{)}\PYG{p}{;}
\PYG{+w}{  }\PYG{k}{else}
\PYG{+w}{  }\PYG{p}{\PYGZob{}}
\PYG{+w}{    }\PYG{k}{auto}\PYG{+w}{ }\PYG{n}{no\PYGZus{}of\PYGZus{}threads}\PYG{+w}{ }\PYG{o}{=}\PYG{+w}{ }\PYG{n}{get\PYGZus{}no\PYGZus{}of\PYGZus{}threads}\PYG{p}{(}\PYG{p}{)}\PYG{p}{;}
\PYG{+w}{    }\PYG{k}{auto}\PYG{+w}{ }\PYG{n}{part}\PYG{+w}{ }\PYG{o}{=}\PYG{+w}{ }\PYG{n}{size}\PYG{+w}{ }\PYG{o}{/}\PYG{+w}{ }\PYG{n}{no\PYGZus{}of\PYGZus{}threads}\PYG{p}{;}
\PYG{+w}{    }\PYG{k}{auto}\PYG{+w}{ }\PYG{n}{last}\PYG{+w}{ }\PYG{o}{=}\PYG{+w}{ }\PYG{n}{begin}\PYG{p}{;}
\PYG{n}{std}\PYG{o}{:}\PYG{o}{:}\PYG{n}{vector}\PYG{o}{\PYGZlt{}}\PYG{n}{std}\PYG{o}{:}\PYG{o}{:}\PYG{k+kr}{thread}\PYG{o}{\PYGZgt{}}\PYG{+w}{ }\PYG{n}{threads}\PYG{p}{;}
\PYG{k}{for}\PYG{p}{(}\PYG{k+kt}{unsigned}\PYG{+w}{ }\PYG{n}{i}\PYG{+w}{ }\PYG{o}{=}\PYG{+w}{ }\PYG{l+m+mi}{0}\PYG{p}{;}\PYG{+w}{ }\PYG{n}{i}\PYG{+w}{ }\PYG{o}{\PYGZlt{}}\PYG{+w}{ }\PYG{n}{no\PYGZus{}of\PYGZus{}threads}\PYG{p}{;}\PYG{+w}{ }\PYG{o}{+}\PYG{o}{+}\PYG{n}{i}\PYG{p}{)}
\PYG{p}{\PYGZob{}}
\PYG{+w}{  }\PYG{k}{if}\PYG{p}{(}\PYG{n}{i}\PYG{+w}{ }\PYG{o}{=}\PYG{o}{=}\PYG{+w}{ }\PYG{n}{no\PYGZus{}of\PYGZus{}threads}\PYG{+w}{ }\PYG{o}{\PYGZhy{}}\PYG{+w}{ }\PYG{l+m+mi}{1}\PYG{p}{)}\PYG{+w}{ }\PYG{n}{last}\PYG{+w}{ }\PYG{o}{=}\PYG{+w}{ }\PYG{n}{end}\PYG{p}{;}
\PYG{k}{else}\PYG{+w}{ }\PYG{n}{std}\PYG{o}{:}\PYG{o}{:}\PYG{n}{advance}\PYG{p}{(}\PYG{n}{last}\PYG{p}{,}\PYG{+w}{ }\PYG{n}{part}\PYG{p}{)}\PYG{p}{;}
\PYG{n}{threads}\PYG{p}{.}\PYG{n}{emplace\PYGZus{}back}\PYG{p}{(}
\PYG{+w}{  }\PYG{p}{[}\PYG{o}{=}\PYG{p}{,}\PYG{o}{\PYGZam{}}\PYG{n}{f}\PYG{p}{]}\PYG{p}{\PYGZob{}}\PYG{n}{std}\PYG{o}{:}\PYG{o}{:}\PYG{n}{transform}\PYG{p}{(}\PYG{n}{begin}\PYG{p}{,}\PYG{+w}{ }\PYG{n}{last}\PYG{p}{,}
\PYG{+w}{                        }\PYG{n}{begin}\PYG{p}{,}\PYG{+w}{ }\PYG{n}{std}\PYG{o}{:}\PYG{o}{:}\PYG{n}{forward}\PYG{o}{\PYGZlt{}}\PYG{n}{F}\PYG{o}{\PYGZgt{}}\PYG{p}{(}\PYG{n}{f}\PYG{p}{)}\PYG{p}{)}\PYG{p}{;}\PYG{p}{\PYGZcb{}}\PYG{p}{)}\PYG{p}{;}
\PYG{+w}{  }\PYG{n}{begin}\PYG{+w}{ }\PYG{o}{=}\PYG{+w}{ }\PYG{n}{last}\PYG{p}{;}
\PYG{p}{\PYGZcb{}}
\PYG{+w}{    }\PYG{k}{for}\PYG{p}{(}\PYG{k}{auto}\PYG{+w}{ }\PYG{o}{\PYGZam{}}\PYG{+w}{ }\PYG{n}{t}\PYG{+w}{ }\PYG{o}{:}\PYG{+w}{ }\PYG{n}{threads}\PYG{p}{)}\PYG{+w}{ }\PYG{n}{t}\PYG{p}{.}\PYG{n}{join}\PYG{p}{(}\PYG{p}{)}\PYG{p}{;}
\PYG{+w}{  }\PYG{p}{\PYGZcb{}}
\PYG{p}{\PYGZcb{}}
\end{Verbatim}
\end{codeblockbg}


%%C080377 p

要实现左\texttt{fold}函数的并行版本，请按以下步骤操作：


%%C080378 p

定义一个函数模板，它接受指向某个范围的\texttt{begin}与\texttt{end}迭代器、一个


%%C080379 p

初始值，以及一个应用到该范围各元素上的二元函数：

%%C080380 code

\begin{codeblockbg}
\begin{Verbatim}[commandchars=\\\{\},breaklines=true,breakanywhere=true,fontsize=\codefont,xleftmargin=17bp]
\PYG{k}{template}\PYG{+w}{ }\PYG{o}{\PYGZlt{}}\PYG{k}{typename}\PYG{+w}{ }\PYG{n+nc}{Iter}\PYG{p}{,}\PYG{+w}{ }\PYG{k}{typename}\PYG{+w}{ }\PYG{n+nc}{R}\PYG{p}{,}\PYG{+w}{ }\PYG{k}{typename}\PYG{+w}{ }\PYG{n+nc}{F}\PYG{o}{\PYGZgt{}}
\PYG{k}{auto}\PYG{+w}{ }\PYG{n}{parallel\PYGZus{}reduce}\PYG{p}{(}\PYG{n}{Iter}\PYG{+w}{ }\PYG{n}{begin}\PYG{p}{,}\PYG{+w}{ }\PYG{n}{Iter}\PYG{+w}{ }\PYG{n}{end}\PYG{p}{,}\PYG{+w}{ }\PYG{n}{R}\PYG{+w}{ }\PYG{n}{init}\PYG{p}{,}\PYG{+w}{ }\PYG{n}{F}\PYG{+w}{ }\PYG{n}{op}\PYG{p}{)}
\PYG{p}{\PYGZob{}}
\PYG{p}{\PYGZcb{}}
\end{Verbatim}
\end{codeblockbg}


%%C080381 p

检查范围的大小。如果元素个数小于一个


%%C080382 p

预先定义的阈值（本实现中为10,000），则以串行方式执行折叠：


%%C080383 code

\begin{codeblockbg}
\begin{Verbatim}[commandchars=\\\{\},breaklines=true,breakanywhere=true,fontsize=\codefont,xleftmargin=17bp]
\PYG{k}{auto}\PYG{+w}{ }\PYG{n}{size}\PYG{+w}{ }\PYG{o}{=}\PYG{+w}{ }\PYG{n}{std}\PYG{o}{:}\PYG{o}{:}\PYG{n}{distance}\PYG{p}{(}\PYG{n}{begin}\PYG{p}{,}\PYG{+w}{ }\PYG{n}{end}\PYG{p}{)}\PYG{p}{;}
\PYG{k}{if}\PYG{p}{(}\PYG{n}{size}\PYG{+w}{ }\PYG{o}{\PYGZlt{}}\PYG{o}{=}\PYG{+w}{ }\PYG{l+m+mi}{10000}\PYG{p}{)}
\PYG{+w}{  }\PYG{k}{return}\PYG{+w}{ }\PYG{n}{std}\PYG{o}{:}\PYG{o}{:}\PYG{n}{accumulate}\PYG{p}{(}\PYG{n}{begin}\PYG{p}{,}\PYG{+w}{ }\PYG{n}{end}\PYG{p}{,}
\PYG{+w}{                         }\PYG{n}{init}\PYG{p}{,}\PYG{+w}{ }\PYG{n}{std}\PYG{o}{:}\PYG{o}{:}\PYG{n}{forward}\PYG{o}{\PYGZlt{}}\PYG{n}{F}\PYG{o}{\PYGZgt{}}\PYG{p}{(}\PYG{n}{op}\PYG{p}{)}\PYG{p}{)}\PYG{p}{;}
\end{Verbatim}
\end{codeblockbg}


%%C080384 p

对于更大的范围，把工作拆分给多个线程，让每个线程折叠范围中的


%%C080385 p

一部分。这些部分不应重叠，以避免对共享数据做线程同步。结果可以通过传给线程函数的引用返回，从而避免数据同步：


%%C080386 code

\begin{codeblockbg}
\begin{Verbatim}[commandchars=\\\{\},breaklines=true,breakanywhere=true,fontsize=\codefont,xleftmargin=17bp]
\PYG{k}{else}
\PYG{p}{\PYGZob{}}
\PYG{+w}{  }\PYG{k}{auto}\PYG{+w}{ }\PYG{n}{no\PYGZus{}of\PYGZus{}threads}\PYG{+w}{ }\PYG{o}{=}\PYG{+w}{ }\PYG{n}{get\PYGZus{}no\PYGZus{}of\PYGZus{}threads}\PYG{p}{(}\PYG{p}{)}\PYG{p}{;}
\PYG{+w}{  }\PYG{k}{auto}\PYG{+w}{ }\PYG{n}{part}\PYG{+w}{ }\PYG{o}{=}\PYG{+w}{ }\PYG{n}{size}\PYG{+w}{ }\PYG{o}{/}\PYG{+w}{ }\PYG{n}{no\PYGZus{}of\PYGZus{}threads}\PYG{p}{;}
\PYG{+w}{  }\PYG{k}{auto}\PYG{+w}{ }\PYG{n}{last}\PYG{+w}{ }\PYG{o}{=}\PYG{+w}{ }\PYG{n}{begin}\PYG{p}{;}
\PYG{+w}{  }\PYG{c+c1}{// 后接第 4、5 步}
\PYG{p}{\PYGZcb{}}
\end{Verbatim}
\end{codeblockbg}


%%C080387 p

启动这些线程，并在每个线程上执行串行版本的折叠：


%%C080388 code

\begin{codeblockbg}
\begin{Verbatim}[commandchars=\\\{\},breaklines=true,breakanywhere=true,fontsize=\codefont,xleftmargin=17bp]
\PYG{n}{std}\PYG{o}{:}\PYG{o}{:}\PYG{n}{vector}\PYG{o}{\PYGZlt{}}\PYG{n}{std}\PYG{o}{:}\PYG{o}{:}\PYG{k+kr}{thread}\PYG{o}{\PYGZgt{}}\PYG{+w}{ }\PYG{n}{threads}\PYG{p}{;}
\PYG{n}{std}\PYG{o}{:}\PYG{o}{:}\PYG{n}{vector}\PYG{o}{\PYGZlt{}}\PYG{n}{R}\PYG{o}{\PYGZgt{}}\PYG{+w}{ }\PYG{n}{values}\PYG{p}{(}\PYG{n}{no\PYGZus{}of\PYGZus{}threads}\PYG{p}{)}\PYG{p}{;}
\PYG{k}{for}\PYG{p}{(}\PYG{k+kt}{unsigned}\PYG{+w}{ }\PYG{n}{i}\PYG{+w}{ }\PYG{o}{=}\PYG{+w}{ }\PYG{l+m+mi}{0}\PYG{p}{;}\PYG{+w}{ }\PYG{n}{i}\PYG{+w}{ }\PYG{o}{\PYGZlt{}}\PYG{+w}{ }\PYG{n}{no\PYGZus{}of\PYGZus{}threads}\PYG{p}{;}\PYG{+w}{ }\PYG{o}{+}\PYG{o}{+}\PYG{n}{i}\PYG{p}{)}
\PYG{p}{\PYGZob{}}
\PYG{+w}{  }\PYG{k}{if}\PYG{p}{(}\PYG{n}{i}\PYG{+w}{ }\PYG{o}{=}\PYG{o}{=}\PYG{+w}{ }\PYG{n}{no\PYGZus{}of\PYGZus{}threads}\PYG{+w}{ }\PYG{o}{\PYGZhy{}}\PYG{+w}{ }\PYG{l+m+mi}{1}\PYG{p}{)}\PYG{+w}{ }\PYG{n}{last}\PYG{+w}{ }\PYG{o}{=}\PYG{+w}{ }\PYG{n}{end}\PYG{p}{;}
\PYG{+w}{  }\PYG{k}{else}\PYG{+w}{ }\PYG{n}{std}\PYG{o}{:}\PYG{o}{:}\PYG{n}{advance}\PYG{p}{(}\PYG{n}{last}\PYG{p}{,}\PYG{+w}{ }\PYG{n}{part}\PYG{p}{)}\PYG{p}{;}
\PYG{+w}{  }\PYG{n}{threads}\PYG{p}{.}\PYG{n}{emplace\PYGZus{}back}\PYG{p}{(}
\PYG{+w}{    }\PYG{p}{[}\PYG{o}{=}\PYG{p}{,}\PYG{o}{\PYGZam{}}\PYG{n}{op}\PYG{p}{]}\PYG{p}{(}\PYG{n}{R}\PYG{o}{\PYGZam{}}\PYG{+w}{ }\PYG{n}{result}\PYG{p}{)}\PYG{p}{\PYGZob{}}
\PYG{+w}{      }\PYG{n}{result}\PYG{+w}{ }\PYG{o}{=}\PYG{+w}{ }\PYG{n}{std}\PYG{o}{:}\PYG{o}{:}\PYG{n}{accumulate}\PYG{p}{(}\PYG{n}{begin}\PYG{p}{,}\PYG{+w}{ }\PYG{n}{last}\PYG{p}{,}\PYG{+w}{ }\PYG{n}{R}\PYG{p}{\PYGZob{}}\PYG{p}{\PYGZcb{}}\PYG{p}{,}
\PYG{+w}{                               }\PYG{n}{std}\PYG{o}{:}\PYG{o}{:}\PYG{n}{forward}\PYG{o}{\PYGZlt{}}\PYG{n}{F}\PYG{o}{\PYGZgt{}}\PYG{p}{(}\PYG{n}{op}\PYG{p}{)}\PYG{p}{)}\PYG{p}{;}\PYG{p}{\PYGZcb{}}\PYG{p}{,}
\PYG{+w}{    }\PYG{n}{std}\PYG{o}{:}\PYG{o}{:}\PYG{n}{ref}\PYG{p}{(}\PYG{n}{values}\PYG{p}{[}\PYG{n}{i}\PYG{p}{]}\PYG{p}{)}\PYG{p}{)}\PYG{p}{;}
\PYG{+w}{  }\PYG{n}{begin}\PYG{+w}{ }\PYG{o}{=}\PYG{+w}{ }\PYG{n}{last}\PYG{p}{;}
\PYG{p}{\PYGZcb{}}
\end{Verbatim}
\end{codeblockbg}


%%C080389 p

等待所有线程执行完毕，并把各部分结果折叠成


%%C080390 p

最终结果：


%%C080391 code

\begin{codeblockbg}
\begin{Verbatim}[commandchars=\\\{\},breaklines=true,breakanywhere=true,fontsize=\codefont,xleftmargin=17bp]
\PYG{k}{for}\PYG{p}{(}\PYG{k}{auto}\PYG{+w}{ }\PYG{o}{\PYGZam{}}\PYG{+w}{ }\PYG{n}{t}\PYG{+w}{ }\PYG{o}{:}\PYG{+w}{ }\PYG{n}{threads}\PYG{p}{)}\PYG{+w}{ }\PYG{n}{t}\PYG{p}{.}\PYG{n}{join}\PYG{p}{(}\PYG{p}{)}\PYG{p}{;}
\PYG{k}{return}\PYG{+w}{ }\PYG{n}{std}\PYG{o}{:}\PYG{o}{:}\PYG{n}{accumulate}\PYG{p}{(}\PYG{n}{std}\PYG{o}{:}\PYG{o}{:}\PYG{n}{begin}\PYG{p}{(}\PYG{n}{values}\PYG{p}{)}\PYG{p}{,}\PYG{+w}{ }\PYG{n}{std}\PYG{o}{:}\PYG{o}{:}\PYG{n}{end}\PYG{p}{(}\PYG{n}{values}\PYG{p}{)}\PYG{p}{,}
\PYG{+w}{                       }\PYG{n}{init}\PYG{p}{,}\PYG{+w}{ }\PYG{n}{std}\PYG{o}{:}\PYG{o}{:}\PYG{n}{forward}\PYG{o}{\PYGZlt{}}\PYG{n}{F}\PYG{o}{\PYGZgt{}}\PYG{p}{(}\PYG{n}{op}\PYG{p}{)}\PYG{p}{)}\PYG{p}{;}
\end{Verbatim}
\end{codeblockbg}


%%C080392 p

把上面的步骤组合起来，就得到下面的实现：


%%C080393 code

\begin{codeblockbg}
\begin{Verbatim}[commandchars=\\\{\},breaklines=true,breakanywhere=true,fontsize=\codefont,xleftmargin=17bp]
\PYG{k}{template}\PYG{+w}{ }\PYG{o}{\PYGZlt{}}\PYG{k}{typename}\PYG{+w}{ }\PYG{n+nc}{Iter}\PYG{p}{,}\PYG{+w}{ }\PYG{k}{typename}\PYG{+w}{ }\PYG{n+nc}{R}\PYG{p}{,}\PYG{+w}{ }\PYG{k}{typename}\PYG{+w}{ }\PYG{n+nc}{F}\PYG{o}{\PYGZgt{}}
\PYG{k}{auto}\PYG{+w}{ }\PYG{n}{parallel\PYGZus{}reduce}\PYG{p}{(}\PYG{n}{Iter}\PYG{+w}{ }\PYG{n}{begin}\PYG{p}{,}\PYG{+w}{ }\PYG{n}{Iter}\PYG{+w}{ }\PYG{n}{end}\PYG{p}{,}\PYG{+w}{ }\PYG{n}{R}\PYG{+w}{ }\PYG{n}{init}\PYG{p}{,}\PYG{+w}{ }\PYG{n}{F}\PYG{+w}{ }\PYG{n}{op}\PYG{p}{)}
\PYG{p}{\PYGZob{}}
\PYG{+w}{  }\PYG{k}{auto}\PYG{+w}{ }\PYG{n}{size}\PYG{+w}{ }\PYG{o}{=}\PYG{+w}{ }\PYG{n}{std}\PYG{o}{:}\PYG{o}{:}\PYG{n}{distance}\PYG{p}{(}\PYG{n}{begin}\PYG{p}{,}\PYG{+w}{ }\PYG{n}{end}\PYG{p}{)}\PYG{p}{;}
\PYG{k}{if}\PYG{p}{(}\PYG{n}{size}\PYG{+w}{ }\PYG{o}{\PYGZlt{}}\PYG{o}{=}\PYG{+w}{ }\PYG{l+m+mi}{10000}\PYG{p}{)}
\PYG{+w}{  }\PYG{k}{return}\PYG{+w}{ }\PYG{n}{std}\PYG{o}{:}\PYG{o}{:}\PYG{n}{accumulate}\PYG{p}{(}\PYG{n}{begin}\PYG{p}{,}\PYG{+w}{ }\PYG{n}{end}\PYG{p}{,}\PYG{+w}{ }\PYG{n}{init}\PYG{p}{,}\PYG{+w}{ }\PYG{n}{std}\PYG{o}{:}\PYG{o}{:}\PYG{n}{forward}\PYG{o}{\PYGZlt{}}\PYG{n}{F}\PYG{o}{\PYGZgt{}}\PYG{p}{(}\PYG{n}{op}\PYG{p}{)}\PYG{p}{)}\PYG{p}{;}
\PYG{k}{else}
\PYG{p}{\PYGZob{}}
\PYG{+w}{  }\PYG{k}{auto}\PYG{+w}{ }\PYG{n}{no\PYGZus{}of\PYGZus{}threads}\PYG{+w}{ }\PYG{o}{=}\PYG{+w}{ }\PYG{n}{get\PYGZus{}no\PYGZus{}of\PYGZus{}threads}\PYG{p}{(}\PYG{p}{)}\PYG{p}{;}
\PYG{+w}{  }\PYG{k}{auto}\PYG{+w}{ }\PYG{n}{part}\PYG{+w}{ }\PYG{o}{=}\PYG{+w}{ }\PYG{n}{size}\PYG{+w}{ }\PYG{o}{/}\PYG{+w}{ }\PYG{n}{no\PYGZus{}of\PYGZus{}threads}\PYG{p}{;}
\PYG{+w}{  }\PYG{k}{auto}\PYG{+w}{ }\PYG{n}{last}\PYG{+w}{ }\PYG{o}{=}\PYG{+w}{ }\PYG{n}{begin}\PYG{p}{;}
\PYG{n}{std}\PYG{o}{:}\PYG{o}{:}\PYG{n}{vector}\PYG{o}{\PYGZlt{}}\PYG{n}{std}\PYG{o}{:}\PYG{o}{:}\PYG{k+kr}{thread}\PYG{o}{\PYGZgt{}}\PYG{+w}{ }\PYG{n}{threads}\PYG{p}{;}
\PYG{n}{std}\PYG{o}{:}\PYG{o}{:}\PYG{n}{vector}\PYG{o}{\PYGZlt{}}\PYG{n}{R}\PYG{o}{\PYGZgt{}}\PYG{+w}{ }\PYG{n}{values}\PYG{p}{(}\PYG{n}{no\PYGZus{}of\PYGZus{}threads}\PYG{p}{)}\PYG{p}{;}
\PYG{k}{for}\PYG{p}{(}\PYG{k+kt}{unsigned}\PYG{+w}{ }\PYG{n}{i}\PYG{+w}{ }\PYG{o}{=}\PYG{+w}{ }\PYG{l+m+mi}{0}\PYG{p}{;}\PYG{+w}{ }\PYG{n}{i}\PYG{+w}{ }\PYG{o}{\PYGZlt{}}\PYG{+w}{ }\PYG{n}{no\PYGZus{}of\PYGZus{}threads}\PYG{p}{;}\PYG{+w}{ }\PYG{o}{+}\PYG{o}{+}\PYG{n}{i}\PYG{p}{)}
\PYG{p}{\PYGZob{}}
\PYG{+w}{  }\PYG{k}{if}\PYG{p}{(}\PYG{n}{i}\PYG{+w}{ }\PYG{o}{=}\PYG{o}{=}\PYG{+w}{ }\PYG{n}{no\PYGZus{}of\PYGZus{}threads}\PYG{+w}{ }\PYG{o}{\PYGZhy{}}\PYG{+w}{ }\PYG{l+m+mi}{1}\PYG{p}{)}\PYG{+w}{ }\PYG{n}{last}\PYG{+w}{ }\PYG{o}{=}\PYG{+w}{ }\PYG{n}{end}\PYG{p}{;}
\PYG{+w}{  }\PYG{k}{else}\PYG{+w}{ }\PYG{n}{std}\PYG{o}{:}\PYG{o}{:}\PYG{n}{advance}\PYG{p}{(}\PYG{n}{last}\PYG{p}{,}\PYG{+w}{ }\PYG{n}{part}\PYG{p}{)}\PYG{p}{;}
\PYG{n}{threads}\PYG{p}{.}\PYG{n}{emplace\PYGZus{}back}\PYG{p}{(}
\PYG{p}{[}\PYG{o}{=}\PYG{p}{,}\PYG{o}{\PYGZam{}}\PYG{n}{op}\PYG{p}{]}\PYG{p}{(}\PYG{n}{R}\PYG{o}{\PYGZam{}}\PYG{+w}{ }\PYG{n}{result}\PYG{p}{)}\PYG{p}{\PYGZob{}}
\PYG{+w}{  }\PYG{n}{result}\PYG{+w}{ }\PYG{o}{=}\PYG{+w}{ }\PYG{n}{std}\PYG{o}{:}\PYG{o}{:}\PYG{n}{accumulate}\PYG{p}{(}\PYG{n}{begin}\PYG{p}{,}\PYG{+w}{ }\PYG{n}{last}\PYG{p}{,}\PYG{+w}{ }\PYG{n}{R}\PYG{p}{\PYGZob{}}\PYG{p}{\PYGZcb{}}\PYG{p}{,}
\PYG{+w}{                           }\PYG{n}{std}\PYG{o}{:}\PYG{o}{:}\PYG{n}{forward}\PYG{o}{\PYGZlt{}}\PYG{n}{F}\PYG{o}{\PYGZgt{}}\PYG{p}{(}\PYG{n}{op}\PYG{p}{)}\PYG{p}{)}\PYG{p}{;}\PYG{p}{\PYGZcb{}}\PYG{p}{,}
\PYG{n}{std}\PYG{o}{:}\PYG{o}{:}\PYG{n}{ref}\PYG{p}{(}\PYG{n}{values}\PYG{p}{[}\PYG{n}{i}\PYG{p}{]}\PYG{p}{)}\PYG{p}{)}\PYG{p}{;}
\PYG{+w}{  }\PYG{n}{begin}\PYG{+w}{ }\PYG{o}{=}\PYG{+w}{ }\PYG{n}{last}\PYG{p}{;}
\PYG{p}{\PYGZcb{}}
\PYG{k}{for}\PYG{p}{(}\PYG{k}{auto}\PYG{+w}{ }\PYG{o}{\PYGZam{}}\PYG{+w}{ }\PYG{n}{t}\PYG{+w}{ }\PYG{o}{:}\PYG{+w}{ }\PYG{n}{threads}\PYG{p}{)}\PYG{+w}{ }\PYG{n}{t}\PYG{p}{.}\PYG{n}{join}\PYG{p}{(}\PYG{p}{)}\PYG{p}{;}
\PYG{+w}{    }\PYG{k}{return}\PYG{+w}{ }\PYG{n}{std}\PYG{o}{:}\PYG{o}{:}\PYG{n}{accumulate}\PYG{p}{(}\PYG{n}{std}\PYG{o}{:}\PYG{o}{:}\PYG{n}{begin}\PYG{p}{(}\PYG{n}{values}\PYG{p}{)}\PYG{p}{,}\PYG{+w}{ }\PYG{n}{std}\PYG{o}{:}\PYG{o}{:}\PYG{n}{end}\PYG{p}{(}\PYG{n}{values}\PYG{p}{)}\PYG{p}{,}
\PYG{+w}{                           }\PYG{n}{init}\PYG{p}{,}\PYG{+w}{ }\PYG{n}{std}\PYG{o}{:}\PYG{o}{:}\PYG{n}{forward}\PYG{o}{\PYGZlt{}}\PYG{n}{F}\PYG{o}{\PYGZgt{}}\PYG{p}{(}\PYG{n}{op}\PYG{p}{)}\PYG{p}{)}\PYG{p}{;}
\PYG{+w}{  }\PYG{p}{\PYGZcb{}}
\PYG{p}{\PYGZcb{}}
\end{Verbatim}
\end{codeblockbg}


%%C080394 h3

\subsection{工作原理}


%%C080395 p

\texttt{map}与\texttt{fold}的这两个并行实现在几个方面是相似的：


%%C080396 li

\begin{itemize}
\item 当范围中的元素个数小于10,000时，两者都会退回到串行版本。


%%C080397 li

\item 两者启动的线程数相同。线程数由静态函数\myttbr{std::\allowbreak{}thread::\allowbreak{}hardware\_\allowbreak{}concurrency()}确定，它返回实现所支持的并发线程数。不过这个值只是一个提示，并不精确，使用时应考虑到这一点。


%%C080398 li

\item 两者都没有使用共享数据，以避免对访问做同步。尽管所有线程都处理同一范围中的元素，但它们处理的是范围中互不重叠的部分。


%%C080399 li

\item 这两个函数都实现为函数模板，接受一对起止迭代器来定义要处理的范围。为了把范围拆分成多个部分、由不同线程独立处理，需要在范围的中间位置使用额外的迭代器。为此，我们用\texttt{std::\allowbreak{}advance()}把迭代器前进指定的位置数。这对vector或数组很有效，但对list之类的容器效率很低。因此，这种实现只适用于具有随机访问迭代器的范围。

\end{itemize}


%%C080400 p

\texttt{map}与\texttt{fold}的串行版本在C++中可以直接用\texttt{std::\allowbreak{}transform()}和\texttt{std::\allowbreak{}accumulate()}实现。事实上，为了验证并行算法的正确性并检查它们是否带来了执行速度的提升，我们可以把它们与这些通用算法的执行情况做比较。


%%C080401 p

为了检验这一点，我们会在一个大小从10,000个元素变化到5000万个元素的vector上使用\texttt{map}与\texttt{fold}。首先对范围做映射（即变换），把每个元素的值翻倍，然后把结果折叠成单个值，也就是把范围中的所有元素相加。为简单起见，范围中每个元素都等于它从1开始的索引（第一个元素是1，第二个元素是2，依此类推）。下面的示例在不同大小的vector上分别运行map与fold的串行版本和并行版本，并以表格形式打印执行时间。


%%C080402 note

\begin{myNote}
作为练习，你可以改变元素个数和线程数，看看并行版本与串行版本相比表现如何。
\end{myNote}


%%C080403 code

\begin{codeblockbg}
\begin{Verbatim}[commandchars=\\\{\},breaklines=true,breakanywhere=true,fontsize=\codefont,xleftmargin=17bp]
\PYG{n}{std}\PYG{o}{:}\PYG{o}{:}\PYG{n}{vector}\PYG{o}{\PYGZlt{}}\PYG{k+kt}{int}\PYG{o}{\PYGZgt{}}\PYG{+w}{ }\PYG{n}{sizes}
\PYG{p}{\PYGZob{}}
\PYG{+w}{  }\PYG{l+m+mi}{10000}\PYG{p}{,}\PYG{+w}{ }\PYG{l+m+mi}{100000}\PYG{p}{,}\PYG{+w}{ }\PYG{l+m+mi}{500000}\PYG{p}{,}
\PYG{+w}{  }\PYG{l+m+mi}{1000000}\PYG{p}{,}\PYG{+w}{ }\PYG{l+m+mi}{2000000}\PYG{p}{,}\PYG{+w}{ }\PYG{l+m+mi}{5000000}\PYG{p}{,}
\PYG{+w}{  }\PYG{l+m+mi}{10000000}\PYG{p}{,}\PYG{+w}{ }\PYG{l+m+mi}{25000000}\PYG{p}{,}\PYG{+w}{ }\PYG{l+m+mi}{50000000}
\PYG{p}{\PYGZcb{}}\PYG{p}{;}
\PYG{n}{std}\PYG{o}{:}\PYG{o}{:}\PYG{n}{cout}
\PYG{+w}{  }\PYG{o}{\PYGZlt{}}\PYG{o}{\PYGZlt{}}\PYG{+w}{ }\PYG{n}{std}\PYG{o}{:}\PYG{o}{:}\PYG{n}{right}\PYG{+w}{ }\PYG{o}{\PYGZlt{}}\PYG{o}{\PYGZlt{}}\PYG{+w}{ }\PYG{n}{std}\PYG{o}{:}\PYG{o}{:}\PYG{n}{setw}\PYG{p}{(}\PYG{l+m+mi}{8}\PYG{p}{)}\PYG{+w}{ }\PYG{o}{\PYGZlt{}}\PYG{o}{\PYGZlt{}}\PYG{+w}{ }\PYG{n}{std}\PYG{o}{:}\PYG{o}{:}\PYG{n}{setfill}\PYG{p}{(}\PYG{l+s+sc}{\PYGZsq{}}\PYG{l+s+sc}{ }\PYG{l+s+sc}{\PYGZsq{}}\PYG{p}{)}\PYG{+w}{ }\PYG{o}{\PYGZlt{}}\PYG{o}{\PYGZlt{}}\PYG{+w}{ }\PYG{l+s}{\PYGZdq{}}\PYG{l+s}{size}\PYG{l+s}{\PYGZdq{}}
\PYG{+w}{  }\PYG{o}{\PYGZlt{}}\PYG{o}{\PYGZlt{}}\PYG{+w}{ }\PYG{n}{std}\PYG{o}{:}\PYG{o}{:}\PYG{n}{right}\PYG{+w}{ }\PYG{o}{\PYGZlt{}}\PYG{o}{\PYGZlt{}}\PYG{+w}{ }\PYG{n}{std}\PYG{o}{:}\PYG{o}{:}\PYG{n}{setw}\PYG{p}{(}\PYG{l+m+mi}{8}\PYG{p}{)}\PYG{+w}{ }\PYG{o}{\PYGZlt{}}\PYG{o}{\PYGZlt{}}\PYG{+w}{ }\PYG{l+s}{\PYGZdq{}}\PYG{l+s}{s map}\PYG{l+s}{\PYGZdq{}}
\PYG{+w}{  }\PYG{o}{\PYGZlt{}}\PYG{o}{\PYGZlt{}}\PYG{+w}{ }\PYG{n}{std}\PYG{o}{:}\PYG{o}{:}\PYG{n}{right}\PYG{+w}{ }\PYG{o}{\PYGZlt{}}\PYG{o}{\PYGZlt{}}\PYG{+w}{ }\PYG{n}{std}\PYG{o}{:}\PYG{o}{:}\PYG{n}{setw}\PYG{p}{(}\PYG{l+m+mi}{8}\PYG{p}{)}\PYG{+w}{ }\PYG{o}{\PYGZlt{}}\PYG{o}{\PYGZlt{}}\PYG{+w}{ }\PYG{l+s}{\PYGZdq{}}\PYG{l+s}{p map}\PYG{l+s}{\PYGZdq{}}
\PYG{+w}{  }\PYG{o}{\PYGZlt{}}\PYG{o}{\PYGZlt{}}\PYG{+w}{ }\PYG{n}{std}\PYG{o}{:}\PYG{o}{:}\PYG{n}{right}\PYG{+w}{ }\PYG{o}{\PYGZlt{}}\PYG{o}{\PYGZlt{}}\PYG{+w}{ }\PYG{n}{std}\PYG{o}{:}\PYG{o}{:}\PYG{n}{setw}\PYG{p}{(}\PYG{l+m+mi}{8}\PYG{p}{)}\PYG{+w}{ }\PYG{o}{\PYGZlt{}}\PYG{o}{\PYGZlt{}}\PYG{+w}{ }\PYG{l+s}{\PYGZdq{}}\PYG{l+s}{s fold}\PYG{l+s}{\PYGZdq{}}
\PYG{+w}{  }\PYG{o}{\PYGZlt{}}\PYG{o}{\PYGZlt{}}\PYG{+w}{ }\PYG{n}{std}\PYG{o}{:}\PYG{o}{:}\PYG{n}{right}\PYG{+w}{ }\PYG{o}{\PYGZlt{}}\PYG{o}{\PYGZlt{}}\PYG{+w}{ }\PYG{n}{std}\PYG{o}{:}\PYG{o}{:}\PYG{n}{setw}\PYG{p}{(}\PYG{l+m+mi}{8}\PYG{p}{)}\PYG{+w}{ }\PYG{o}{\PYGZlt{}}\PYG{o}{\PYGZlt{}}\PYG{+w}{ }\PYG{l+s}{\PYGZdq{}}\PYG{l+s}{p fold}\PYG{l+s}{\PYGZdq{}}
\PYG{+w}{  }\PYG{o}{\PYGZlt{}}\PYG{o}{\PYGZlt{}}\PYG{+w}{ }\PYG{n}{std}\PYG{o}{:}\PYG{o}{:}\PYG{n}{endl}\PYG{p}{;}
\PYG{k}{for}\PYG{+w}{ }\PYG{p}{(}\PYG{k}{auto}\PYG{+w}{ }\PYG{k}{const}\PYG{+w}{ }\PYG{n}{size}\PYG{+w}{ }\PYG{o}{:}\PYG{+w}{ }\PYG{n}{sizes}\PYG{p}{)}
\PYG{p}{\PYGZob{}}
\PYG{+w}{  }\PYG{n}{std}\PYG{o}{:}\PYG{o}{:}\PYG{n}{vector}\PYG{o}{\PYGZlt{}}\PYG{k+kt}{int}\PYG{o}{\PYGZgt{}}\PYG{+w}{ }\PYG{n}{v}\PYG{p}{(}\PYG{n}{size}\PYG{p}{)}\PYG{p}{;}
\PYG{+w}{  }\PYG{n}{std}\PYG{o}{:}\PYG{o}{:}\PYG{n}{iota}\PYG{p}{(}\PYG{n}{std}\PYG{o}{:}\PYG{o}{:}\PYG{n}{begin}\PYG{p}{(}\PYG{n}{v}\PYG{p}{)}\PYG{p}{,}\PYG{+w}{ }\PYG{n}{std}\PYG{o}{:}\PYG{o}{:}\PYG{n}{end}\PYG{p}{(}\PYG{n}{v}\PYG{p}{)}\PYG{p}{,}\PYG{+w}{ }\PYG{l+m+mi}{1}\PYG{p}{)}\PYG{p}{;}
\PYG{k}{auto}\PYG{+w}{ }\PYG{n}{v1}\PYG{+w}{ }\PYG{o}{=}\PYG{+w}{ }\PYG{n}{v}\PYG{p}{;}
\PYG{k}{auto}\PYG{+w}{ }\PYG{n}{s1}\PYG{+w}{ }\PYG{o}{=}\PYG{+w}{ }\PYG{l+m+mf}{0L}\PYG{n}{L}\PYG{p}{;}
\PYG{k}{auto}\PYG{+w}{ }\PYG{n}{tsm}\PYG{+w}{ }\PYG{o}{=}\PYG{+w}{ }\PYG{n}{perf\PYGZus{}timer}\PYG{o}{\PYGZlt{}}\PYG{o}{\PYGZgt{}}\PYG{o}{:}\PYG{o}{:}\PYG{n}{duration}\PYG{p}{(}\PYG{p}{[}\PYG{o}{\PYGZam{}}\PYG{p}{]}\PYG{+w}{ }\PYG{p}{\PYGZob{}}
\PYG{+w}{  }\PYG{n}{std}\PYG{o}{:}\PYG{o}{:}\PYG{n}{transform}\PYG{p}{(}\PYG{n}{std}\PYG{o}{:}\PYG{o}{:}\PYG{n}{begin}\PYG{p}{(}\PYG{n}{v1}\PYG{p}{)}\PYG{p}{,}\PYG{+w}{ }\PYG{n}{std}\PYG{o}{:}\PYG{o}{:}\PYG{n}{end}\PYG{p}{(}\PYG{n}{v1}\PYG{p}{)}\PYG{p}{,}\PYG{+w}{ }\PYG{n}{std}\PYG{o}{:}\PYG{o}{:}\PYG{n}{begin}\PYG{p}{(}\PYG{n}{v1}\PYG{p}{)}\PYG{p}{,}
\PYG{+w}{                 }\PYG{p}{[}\PYG{p}{]}\PYG{p}{(}\PYG{k+kt}{int}\PYG{+w}{ }\PYG{k}{const}\PYG{+w}{ }\PYG{n}{i}\PYG{p}{)}\PYG{+w}{ }\PYG{p}{\PYGZob{}}\PYG{k}{return}\PYG{+w}{ }\PYG{n}{i}\PYG{+w}{ }\PYG{o}{+}\PYG{+w}{ }\PYG{n}{i}\PYG{p}{;}\PYG{+w}{ }\PYG{p}{\PYGZcb{}}\PYG{p}{)}\PYG{p}{;}\PYG{+w}{ }\PYG{p}{\PYGZcb{}}\PYG{p}{)}\PYG{p}{;}
\PYG{k}{auto}\PYG{+w}{ }\PYG{n}{tsf}\PYG{+w}{ }\PYG{o}{=}\PYG{+w}{ }\PYG{n}{perf\PYGZus{}timer}\PYG{o}{\PYGZlt{}}\PYG{o}{\PYGZgt{}}\PYG{o}{:}\PYG{o}{:}\PYG{n}{duration}\PYG{p}{(}\PYG{p}{[}\PYG{o}{\PYGZam{}}\PYG{p}{]}\PYG{+w}{ }\PYG{p}{\PYGZob{}}
\PYG{+w}{  }\PYG{n}{s1}\PYG{+w}{ }\PYG{o}{=}\PYG{+w}{ }\PYG{n}{std}\PYG{o}{:}\PYG{o}{:}\PYG{n}{accumulate}\PYG{p}{(}\PYG{n}{std}\PYG{o}{:}\PYG{o}{:}\PYG{n}{begin}\PYG{p}{(}\PYG{n}{v1}\PYG{p}{)}\PYG{p}{,}\PYG{+w}{ }\PYG{n}{std}\PYG{o}{:}\PYG{o}{:}\PYG{n}{end}\PYG{p}{(}\PYG{n}{v1}\PYG{p}{)}\PYG{p}{,}\PYG{+w}{ }\PYG{l+m+mf}{0L}\PYG{n}{L}\PYG{p}{,}
\PYG{+w}{                       }\PYG{n}{std}\PYG{o}{:}\PYG{o}{:}\PYG{n}{plus}\PYG{o}{\PYGZlt{}}\PYG{o}{\PYGZgt{}}\PYG{p}{(}\PYG{p}{)}\PYG{p}{)}\PYG{p}{;}\PYG{+w}{ }\PYG{p}{\PYGZcb{}}\PYG{p}{)}\PYG{p}{;}
\PYG{k}{auto}\PYG{+w}{ }\PYG{n}{v2}\PYG{+w}{ }\PYG{o}{=}\PYG{+w}{ }\PYG{n}{v}\PYG{p}{;}
\PYG{k}{auto}\PYG{+w}{ }\PYG{n}{s2}\PYG{+w}{ }\PYG{o}{=}\PYG{+w}{ }\PYG{l+m+mf}{0L}\PYG{n}{L}\PYG{p}{;}
\PYG{k}{auto}\PYG{+w}{ }\PYG{n}{tpm}\PYG{+w}{ }\PYG{o}{=}\PYG{+w}{ }\PYG{n}{perf\PYGZus{}timer}\PYG{o}{\PYGZlt{}}\PYG{o}{\PYGZgt{}}\PYG{o}{:}\PYG{o}{:}\PYG{n}{duration}\PYG{p}{(}\PYG{p}{[}\PYG{o}{\PYGZam{}}\PYG{p}{]}\PYG{+w}{ }\PYG{p}{\PYGZob{}}
\PYG{+w}{  }\PYG{n}{parallel\PYGZus{}map}\PYG{p}{(}\PYG{n}{std}\PYG{o}{:}\PYG{o}{:}\PYG{n}{begin}\PYG{p}{(}\PYG{n}{v2}\PYG{p}{)}\PYG{p}{,}\PYG{+w}{ }\PYG{n}{std}\PYG{o}{:}\PYG{o}{:}\PYG{n}{end}\PYG{p}{(}\PYG{n}{v2}\PYG{p}{)}\PYG{p}{,}
\PYG{+w}{               }\PYG{p}{[}\PYG{p}{]}\PYG{p}{(}\PYG{k+kt}{int}\PYG{+w}{ }\PYG{k}{const}\PYG{+w}{ }\PYG{n}{i}\PYG{p}{)}\PYG{+w}{ }\PYG{p}{\PYGZob{}}\PYG{k}{return}\PYG{+w}{ }\PYG{n}{i}\PYG{+w}{ }\PYG{o}{+}\PYG{+w}{ }\PYG{n}{i}\PYG{p}{;}\PYG{+w}{ }\PYG{p}{\PYGZcb{}}\PYG{p}{)}\PYG{p}{;}\PYG{+w}{ }\PYG{p}{\PYGZcb{}}\PYG{p}{)}\PYG{p}{;}
\PYG{k}{auto}\PYG{+w}{ }\PYG{n}{tpf}\PYG{+w}{ }\PYG{o}{=}\PYG{+w}{ }\PYG{n}{perf\PYGZus{}timer}\PYG{o}{\PYGZlt{}}\PYG{o}{\PYGZgt{}}\PYG{o}{:}\PYG{o}{:}\PYG{n}{duration}\PYG{p}{(}\PYG{p}{[}\PYG{o}{\PYGZam{}}\PYG{p}{]}\PYG{+w}{ }\PYG{p}{\PYGZob{}}
\PYG{+w}{  }\PYG{n}{s2}\PYG{+w}{ }\PYG{o}{=}\PYG{+w}{ }\PYG{n}{parallel\PYGZus{}reduce}\PYG{p}{(}\PYG{n}{std}\PYG{o}{:}\PYG{o}{:}\PYG{n}{begin}\PYG{p}{(}\PYG{n}{v2}\PYG{p}{)}\PYG{p}{,}\PYG{+w}{ }\PYG{n}{std}\PYG{o}{:}\PYG{o}{:}\PYG{n}{end}\PYG{p}{(}\PYG{n}{v2}\PYG{p}{)}\PYG{p}{,}\PYG{+w}{ }\PYG{l+m+mf}{0L}\PYG{n}{L}\PYG{p}{,}
\PYG{+w}{                       }\PYG{n}{std}\PYG{o}{:}\PYG{o}{:}\PYG{n}{plus}\PYG{o}{\PYGZlt{}}\PYG{o}{\PYGZgt{}}\PYG{p}{(}\PYG{p}{)}\PYG{p}{)}\PYG{p}{;}\PYG{+w}{ }\PYG{p}{\PYGZcb{}}\PYG{p}{)}\PYG{p}{;}
\PYG{n}{assert}\PYG{p}{(}\PYG{n}{v1}\PYG{+w}{ }\PYG{o}{=}\PYG{o}{=}\PYG{+w}{ }\PYG{n}{v2}\PYG{p}{)}\PYG{p}{;}
\PYG{n}{assert}\PYG{p}{(}\PYG{n}{s1}\PYG{+w}{ }\PYG{o}{=}\PYG{o}{=}\PYG{+w}{ }\PYG{n}{s2}\PYG{p}{)}\PYG{p}{;}
\PYG{+w}{  }\PYG{n}{std}\PYG{o}{:}\PYG{o}{:}\PYG{n}{cout}
\PYG{+w}{    }\PYG{o}{\PYGZlt{}}\PYG{o}{\PYGZlt{}}\PYG{+w}{ }\PYG{n}{std}\PYG{o}{:}\PYG{o}{:}\PYG{n}{right}\PYG{+w}{ }\PYG{o}{\PYGZlt{}}\PYG{o}{\PYGZlt{}}\PYG{+w}{ }\PYG{n}{std}\PYG{o}{:}\PYG{o}{:}\PYG{n}{setw}\PYG{p}{(}\PYG{l+m+mi}{8}\PYG{p}{)}\PYG{+w}{ }\PYG{o}{\PYGZlt{}}\PYG{o}{\PYGZlt{}}\PYG{+w}{ }\PYG{n}{std}\PYG{o}{:}\PYG{o}{:}\PYG{n}{setfill}\PYG{p}{(}\PYG{l+s+sc}{\PYGZsq{}}\PYG{l+s+sc}{ }\PYG{l+s+sc}{\PYGZsq{}}\PYG{p}{)}\PYG{+w}{ }\PYG{o}{\PYGZlt{}}\PYG{o}{\PYGZlt{}}\PYG{+w}{ }\PYG{n}{size}
\PYG{+w}{    }\PYG{o}{\PYGZlt{}}\PYG{o}{\PYGZlt{}}\PYG{+w}{ }\PYG{n}{std}\PYG{o}{:}\PYG{o}{:}\PYG{n}{right}\PYG{+w}{ }\PYG{o}{\PYGZlt{}}\PYG{o}{\PYGZlt{}}\PYG{+w}{ }\PYG{n}{std}\PYG{o}{:}\PYG{o}{:}\PYG{n}{setw}\PYG{p}{(}\PYG{l+m+mi}{8}\PYG{p}{)}
\PYG{+w}{    }\PYG{o}{\PYGZlt{}}\PYG{o}{\PYGZlt{}}\PYG{+w}{ }\PYG{n}{std}\PYG{o}{:}\PYG{o}{:}\PYG{n}{chrono}\PYG{o}{:}\PYG{o}{:}\PYG{n}{duration}\PYG{o}{\PYGZlt{}}\PYG{k+kt}{double}\PYG{p}{,}\PYG{+w}{ }\PYG{n}{td}\PYG{o}{:}\PYG{o}{:}\PYG{n}{micro}\PYG{o}{\PYGZgt{}}\PYG{p}{(}\PYG{n}{tsm}\PYG{p}{)}\PYG{p}{.}\PYG{n}{count}\PYG{p}{(}\PYG{p}{)}
\PYG{+w}{    }\PYG{o}{\PYGZlt{}}\PYG{o}{\PYGZlt{}}\PYG{+w}{ }\PYG{n}{std}\PYG{o}{:}\PYG{o}{:}\PYG{n}{right}\PYG{+w}{ }\PYG{o}{\PYGZlt{}}\PYG{o}{\PYGZlt{}}\PYG{+w}{ }\PYG{n}{std}\PYG{o}{:}\PYG{o}{:}\PYG{n}{setw}\PYG{p}{(}\PYG{l+m+mi}{8}\PYG{p}{)}
\PYG{+w}{    }\PYG{o}{\PYGZlt{}}\PYG{o}{\PYGZlt{}}\PYG{+w}{ }\PYG{n}{std}\PYG{o}{:}\PYG{o}{:}\PYG{n}{chrono}\PYG{o}{:}\PYG{o}{:}\PYG{n}{duration}\PYG{o}{\PYGZlt{}}\PYG{k+kt}{double}\PYG{p}{,}\PYG{+w}{ }\PYG{n}{std}\PYG{o}{:}\PYG{o}{:}\PYG{n}{micro}\PYG{o}{\PYGZgt{}}\PYG{p}{(}\PYG{n}{tpm}\PYG{p}{)}\PYG{p}{.}\PYG{n}{count}\PYG{p}{(}\PYG{p}{)}
\PYG{+w}{    }\PYG{o}{\PYGZlt{}}\PYG{o}{\PYGZlt{}}\PYG{+w}{ }\PYG{n}{std}\PYG{o}{:}\PYG{o}{:}\PYG{n}{right}\PYG{+w}{ }\PYG{o}{\PYGZlt{}}\PYG{o}{\PYGZlt{}}\PYG{+w}{ }\PYG{n}{std}\PYG{o}{:}\PYG{o}{:}\PYG{n}{setw}\PYG{p}{(}\PYG{l+m+mi}{8}\PYG{p}{)}
\PYG{+w}{    }\PYG{o}{\PYGZlt{}}\PYG{o}{\PYGZlt{}}\PYG{+w}{ }\PYG{n}{std}\PYG{o}{:}\PYG{o}{:}\PYG{n}{chrono}\PYG{o}{:}\PYG{o}{:}\PYG{n}{duration}\PYG{o}{\PYGZlt{}}\PYG{k+kt}{double}\PYG{p}{,}\PYG{+w}{ }\PYG{n}{std}\PYG{o}{:}\PYG{o}{:}\PYG{n}{micro}\PYG{o}{\PYGZgt{}}\PYG{p}{(}\PYG{n}{tsf}\PYG{p}{)}\PYG{p}{.}\PYG{n}{count}\PYG{p}{(}\PYG{p}{)}
\PYG{+w}{    }\PYG{o}{\PYGZlt{}}\PYG{o}{\PYGZlt{}}\PYG{+w}{ }\PYG{n}{std}\PYG{o}{:}\PYG{o}{:}\PYG{n}{right}\PYG{+w}{ }\PYG{o}{\PYGZlt{}}\PYG{o}{\PYGZlt{}}\PYG{+w}{ }\PYG{n}{std}\PYG{o}{:}\PYG{o}{:}\PYG{n}{setw}\PYG{p}{(}\PYG{l+m+mi}{8}\PYG{p}{)}
\PYG{+w}{    }\PYG{o}{\PYGZlt{}}\PYG{o}{\PYGZlt{}}\PYG{+w}{ }\PYG{n}{std}\PYG{o}{:}\PYG{o}{:}\PYG{n}{chrono}\PYG{o}{:}\PYG{o}{:}\PYG{n}{duration}\PYG{o}{\PYGZlt{}}\PYG{k+kt}{double}\PYG{p}{,}\PYG{+w}{ }\PYG{n}{std}\PYG{o}{:}\PYG{o}{:}\PYG{n}{micro}\PYG{o}{\PYGZgt{}}\PYG{p}{(}\PYG{n}{tpf}\PYG{p}{)}\PYG{p}{.}\PYG{n}{count}\PYG{p}{(}\PYG{p}{)}
\PYG{+w}{    }\PYG{o}{\PYGZlt{}}\PYG{o}{\PYGZlt{}}\PYG{+w}{ }\PYG{n}{std}\PYG{o}{:}\PYG{o}{:}\PYG{n}{endl}\PYG{p}{;}
\PYG{p}{\PYGZcb{}}
\end{Verbatim}
\end{codeblockbg}


%%C080404 p

该程序的一个可能输出如下面的截图所示（在一台运行64位Windows、配备Intel Core i7处理器、4个物理核与8个逻辑核的机器上执行）。并行版本，尤其是\texttt{fold}的实现，比串行版本表现更好。但这只在vector长度超过一定规模时才成立。从下表中可以看到，元素个数直到100万时，串行版本仍然更快；当vector中有200万个或更多元素时，并行版本执行得更快。要注意，实际耗时每次运行都会略有不同，而在不同机器上可能差别很大：


%%C080405 code

\begin{codeblockbg}
\begin{Verbatim}[commandchars=\\\{\},breaklines=true,breakanywhere=true,fontsize=\codefont,xleftmargin=17bp]
\PYG{+w}{    }\PYG{n}{size}\PYG{+w}{   }\PYG{n}{s}\PYG{+w}{ }\PYG{n}{map}\PYG{+w}{   }\PYG{n}{p}\PYG{+w}{ }\PYG{n}{map}\PYG{+w}{  }\PYG{n}{s}\PYG{+w}{ }\PYG{n}{fold}\PYG{+w}{  }\PYG{n}{p}\PYG{+w}{ }\PYG{n}{fold}
\PYG{+w}{   }\PYG{l+m+mi}{10000}\PYG{+w}{      }\PYG{l+m+mi}{11}\PYG{+w}{      }\PYG{l+m+mi}{10}\PYG{+w}{       }\PYG{l+m+mi}{7}\PYG{+w}{      }\PYG{l+m+mi}{10}
\PYG{+w}{  }\PYG{l+m+mi}{100000}\PYG{+w}{     }\PYG{l+m+mi}{108}\PYG{+w}{    }\PYG{l+m+mi}{1573}\PYG{+w}{      }\PYG{l+m+mi}{72}\PYG{+w}{     }\PYG{l+m+mi}{710}
\PYG{+w}{  }\PYG{l+m+mi}{500000}\PYG{+w}{     }\PYG{l+m+mi}{547}\PYG{+w}{    }\PYG{l+m+mi}{2006}\PYG{+w}{     }\PYG{l+m+mi}{361}\PYG{+w}{     }\PYG{l+m+mi}{862}
\PYG{+w}{ }\PYG{l+m+mi}{1000000}\PYG{+w}{    }\PYG{l+m+mi}{1146}\PYG{+w}{    }\PYG{l+m+mi}{1163}\PYG{+w}{     }\PYG{l+m+mi}{749}\PYG{+w}{     }\PYG{l+m+mi}{862}
\PYG{+w}{ }\PYG{l+m+mi}{2000000}\PYG{+w}{    }\PYG{l+m+mi}{2503}\PYG{+w}{    }\PYG{l+m+mi}{1527}\PYG{+w}{    }\PYG{l+m+mi}{1677}\PYG{+w}{    }\PYG{l+m+mi}{1289}
\PYG{+w}{ }\PYG{l+m+mi}{5000000}\PYG{+w}{    }\PYG{l+m+mi}{5937}\PYG{+w}{    }\PYG{l+m+mi}{3000}\PYG{+w}{    }\PYG{l+m+mi}{4203}\PYG{+w}{    }\PYG{l+m+mi}{2314}
\PYG{l+m+mi}{10000000}\PYG{+w}{   }\PYG{l+m+mi}{11959}\PYG{+w}{    }\PYG{l+m+mi}{6269}\PYG{+w}{    }\PYG{l+m+mi}{8269}\PYG{+w}{    }\PYG{l+m+mi}{3868}
\PYG{l+m+mi}{25000000}\PYG{+w}{   }\PYG{l+m+mi}{29872}\PYG{+w}{   }\PYG{l+m+mi}{13823}\PYG{+w}{   }\PYG{l+m+mi}{20961}\PYG{+w}{    }\PYG{l+m+mi}{9156}
\PYG{l+m+mi}{50000000}\PYG{+w}{   }\PYG{l+m+mi}{60049}\PYG{+w}{   }\PYG{l+m+mi}{27457}\PYG{+w}{   }\PYG{l+m+mi}{41374}\PYG{+w}{   }\PYG{l+m+mi}{19075}
\end{Verbatim}
\end{codeblockbg}


%%C080406 p

为了更好地直观呈现这些结果，我们可以用柱状图表示并行版本的加速比。下图中，蓝色柱表示并行\texttt{map}实现的加速比，橙色柱表示并行\texttt{fold}实现的加速比。正值表示并行版本更快；负值表示串行版本更快：


%%C080407 img

\myFigure{images/ch08/fig-8-1.png}{}


%%C080408 h3

\subsection{延伸阅读}


%%C080409 li

\begin{itemize}
\item \textit{实现高阶函数map与fold}案例，见\texttt{Chapter 3}\textit{探索函数}


%%C080410 li

\item \textit{用任务实现并行map与fold}


%%C080411 li

\item \textit{使用线程}

\end{itemize}


%%C080412 h2

\section{用任务实现并行map与fold}


%%C080413 p

任务（task）是执行并发计算时比线程更高层的替代方案。\texttt{std::\allowbreak{}async()}让我们可以异步执行函数，而无需处理底层的线程细节。本案例要完成与上一案例相同的任务——实现\texttt{map}与\texttt{fold}函数的并行版本，不过这次改用任务，并看看它与线程版本相比如何。


%%C080414 h3

\subsection{准备工作}


%%C080415 p

本案例给出的方案在很多方面与上一案例\textit{用线程实现并行map与fold}中使用线程的方案相似。在继续本案例之前，请务必先阅读那个案例。


%%C080416 h3

\subsection{具体做法}


%%C080417 p

要实现\texttt{map}函数的并行版本，请按以下步骤操作：


%%C080418 p

定义一个函数模板，它接受指向某个范围的起始与结束迭代器，以及一个


%%C080419 p

应用到所有元素上的函数：


%%C080420 code

\begin{codeblockbg}
\begin{Verbatim}[commandchars=\\\{\},breaklines=true,breakanywhere=true,fontsize=\codefont,xleftmargin=17bp]
\PYG{k}{template}\PYG{+w}{ }\PYG{o}{\PYGZlt{}}\PYG{k}{typename}\PYG{+w}{ }\PYG{n+nc}{Iter}\PYG{p}{,}\PYG{+w}{ }\PYG{k}{typename}\PYG{+w}{ }\PYG{n+nc}{F}\PYG{o}{\PYGZgt{}}
\PYG{k+kt}{void}\PYG{+w}{ }\PYG{n}{parallel\PYGZus{}map}\PYG{p}{(}\PYG{n}{Iter}\PYG{+w}{ }\PYG{n}{begin}\PYG{p}{,}\PYG{+w}{ }\PYG{n}{Iter}\PYG{+w}{ }\PYG{n}{end}\PYG{p}{,}\PYG{+w}{ }\PYG{n}{F}\PYG{+w}{ }\PYG{n}{f}\PYG{p}{)}
\PYG{p}{\PYGZob{}}
\PYG{p}{\PYGZcb{}}
\end{Verbatim}
\end{codeblockbg}


%%C080421 p

检查范围的大小。如果元素个数小于


%%C080422 p

预先定义的阈值（本实现中该阈值为10,000），则以串行方式执行映射：


%%C080423 code

\begin{codeblockbg}
\begin{Verbatim}[commandchars=\\\{\},breaklines=true,breakanywhere=true,fontsize=\codefont,xleftmargin=17bp]
\PYG{k}{auto}\PYG{+w}{ }\PYG{n}{size}\PYG{+w}{ }\PYG{o}{=}\PYG{+w}{ }\PYG{n}{std}\PYG{o}{:}\PYG{o}{:}\PYG{n}{distance}\PYG{p}{(}\PYG{n}{begin}\PYG{p}{,}\PYG{+w}{ }\PYG{n}{end}\PYG{p}{)}\PYG{p}{;}
\PYG{k}{if}\PYG{p}{(}\PYG{n}{size}\PYG{+w}{ }\PYG{o}{\PYGZlt{}}\PYG{o}{=}\PYG{+w}{ }\PYG{l+m+mi}{10000}\PYG{p}{)}
\PYG{+w}{  }\PYG{n}{std}\PYG{o}{:}\PYG{o}{:}\PYG{n}{transform}\PYG{p}{(}\PYG{n}{begin}\PYG{p}{,}\PYG{+w}{ }\PYG{n}{end}\PYG{p}{,}\PYG{+w}{ }\PYG{n}{begin}\PYG{p}{,}\PYG{+w}{ }\PYG{n}{std}\PYG{o}{:}\PYG{o}{:}\PYG{n}{forward}\PYG{o}{\PYGZlt{}}\PYG{n}{F}\PYG{o}{\PYGZgt{}}\PYG{p}{(}\PYG{n}{f}\PYG{p}{)}\PYG{p}{)}\PYG{p}{;}
\end{Verbatim}
\end{codeblockbg}


%%C080424 p

对于更大的范围，把工作拆分给多个任务，让每个任务映射


%%C080425 p

范围中的一部分。这些部分不应重叠，以避免同步线程对共享数据的访问：


%%C080426 code

\begin{codeblockbg}
\begin{Verbatim}[commandchars=\\\{\},breaklines=true,breakanywhere=true,fontsize=\codefont,xleftmargin=17bp]
\PYG{k}{else}
\PYG{p}{\PYGZob{}}
\PYG{+w}{  }\PYG{k}{auto}\PYG{+w}{ }\PYG{n}{no\PYGZus{}of\PYGZus{}tasks}\PYG{+w}{ }\PYG{o}{=}\PYG{+w}{ }\PYG{n}{get\PYGZus{}no\PYGZus{}of\PYGZus{}threads}\PYG{p}{(}\PYG{p}{)}\PYG{p}{;}
\PYG{+w}{  }\PYG{k}{auto}\PYG{+w}{ }\PYG{n}{part}\PYG{+w}{ }\PYG{o}{=}\PYG{+w}{ }\PYG{n}{size}\PYG{+w}{ }\PYG{o}{/}\PYG{+w}{ }\PYG{n}{no\PYGZus{}of\PYGZus{}tasks}\PYG{p}{;}
\PYG{+w}{  }\PYG{k}{auto}\PYG{+w}{ }\PYG{n}{last}\PYG{+w}{ }\PYG{o}{=}\PYG{+w}{ }\PYG{n}{begin}\PYG{p}{;}
\PYG{+w}{  }\PYG{c+c1}{// 后接第 4、5 步}
\PYG{p}{\PYGZcb{}}
\end{Verbatim}
\end{codeblockbg}


%%C080427 p

启动这些异步函数，并在其中每一个上运行


%%C080428 p

串行版本的映射：


%%C080429 code

\begin{codeblockbg}
\begin{Verbatim}[commandchars=\\\{\},breaklines=true,breakanywhere=true,fontsize=\codefont,xleftmargin=17bp]
\PYG{n}{std}\PYG{o}{:}\PYG{o}{:}\PYG{n}{vector}\PYG{o}{\PYGZlt{}}\PYG{n}{std}\PYG{o}{:}\PYG{o}{:}\PYG{n}{future}\PYG{o}{\PYGZlt{}}\PYG{k+kt}{void}\PYG{o}{\PYGZgt{}}\PYG{o}{\PYGZgt{}}\PYG{+w}{ }\PYG{n}{tasks}\PYG{p}{;}
\PYG{k}{for}\PYG{p}{(}\PYG{k+kt}{unsigned}\PYG{+w}{ }\PYG{n}{i}\PYG{+w}{ }\PYG{o}{=}\PYG{+w}{ }\PYG{l+m+mi}{0}\PYG{p}{;}\PYG{+w}{ }\PYG{n}{i}\PYG{+w}{ }\PYG{o}{\PYGZlt{}}\PYG{+w}{ }\PYG{n}{no\PYGZus{}of\PYGZus{}tasks}\PYG{p}{;}\PYG{+w}{ }\PYG{o}{+}\PYG{o}{+}\PYG{n}{i}\PYG{p}{)}
\PYG{p}{\PYGZob{}}
\PYG{+w}{  }\PYG{k}{if}\PYG{p}{(}\PYG{n}{i}\PYG{+w}{ }\PYG{o}{=}\PYG{o}{=}\PYG{+w}{ }\PYG{n}{no\PYGZus{}of\PYGZus{}tasks}\PYG{+w}{ }\PYG{o}{\PYGZhy{}}\PYG{+w}{ }\PYG{l+m+mi}{1}\PYG{p}{)}\PYG{+w}{ }\PYG{n}{last}\PYG{+w}{ }\PYG{o}{=}\PYG{+w}{ }\PYG{n}{end}\PYG{p}{;}
\PYG{+w}{  }\PYG{k}{else}\PYG{+w}{ }\PYG{n}{std}\PYG{o}{:}\PYG{o}{:}\PYG{n}{advance}\PYG{p}{(}\PYG{n}{last}\PYG{p}{,}\PYG{+w}{ }\PYG{n}{part}\PYG{p}{)}\PYG{p}{;}
\PYG{+w}{  }\PYG{n}{tasks}\PYG{p}{.}\PYG{n}{emplace\PYGZus{}back}\PYG{p}{(}\PYG{n}{std}\PYG{o}{:}\PYG{o}{:}\PYG{n}{async}\PYG{p}{(}
\PYG{+w}{    }\PYG{n}{std}\PYG{o}{:}\PYG{o}{:}\PYG{n}{launch}\PYG{o}{:}\PYG{o}{:}\PYG{n}{async}\PYG{p}{,}
\PYG{+w}{      }\PYG{p}{[}\PYG{o}{=}\PYG{p}{,}\PYG{o}{\PYGZam{}}\PYG{n}{f}\PYG{p}{]}\PYG{p}{\PYGZob{}}\PYG{n}{std}\PYG{o}{:}\PYG{o}{:}\PYG{n}{transform}\PYG{p}{(}\PYG{n}{begin}\PYG{p}{,}\PYG{+w}{ }\PYG{n}{last}\PYG{p}{,}\PYG{+w}{ }\PYG{n}{begin}\PYG{p}{,}
\PYG{+w}{                            }\PYG{n}{std}\PYG{o}{:}\PYG{o}{:}\PYG{n}{forward}\PYG{o}{\PYGZlt{}}\PYG{n}{F}\PYG{o}{\PYGZgt{}}\PYG{p}{(}\PYG{n}{f}\PYG{p}{)}\PYG{p}{)}\PYG{p}{;}\PYG{p}{\PYGZcb{}}\PYG{p}{)}\PYG{p}{)}\PYG{p}{;}
\PYG{+w}{  }\PYG{n}{begin}\PYG{+w}{ }\PYG{o}{=}\PYG{+w}{ }\PYG{n}{last}\PYG{p}{;}
\PYG{p}{\PYGZcb{}}
\end{Verbatim}
\end{codeblockbg}


%%C080430 p

等待所有异步函数执行完毕：


%%C080431 code

\begin{codeblockbg}
\begin{Verbatim}[commandchars=\\\{\},breaklines=true,breakanywhere=true,fontsize=\codefont,xleftmargin=17bp]
\PYG{k}{for}\PYG{p}{(}\PYG{k}{auto}\PYG{+w}{ }\PYG{o}{\PYGZam{}}\PYG{+w}{ }\PYG{n}{t}\PYG{+w}{ }\PYG{o}{:}\PYG{+w}{ }\PYG{n}{tasks}\PYG{p}{)}\PYG{+w}{ }\PYG{n}{t}\PYG{p}{.}\PYG{n}{wait}\PYG{p}{(}\PYG{p}{)}\PYG{p}{;}
\end{Verbatim}
\end{codeblockbg}


%%C080432 p

把这些步骤组合起来，就得到下面的实现：


%%C080433 code

\begin{codeblockbg}
\begin{Verbatim}[commandchars=\\\{\},breaklines=true,breakanywhere=true,fontsize=\codefont,xleftmargin=17bp]
\PYG{k}{template}\PYG{+w}{ }\PYG{o}{\PYGZlt{}}\PYG{k}{typename}\PYG{+w}{ }\PYG{n+nc}{Iter}\PYG{p}{,}\PYG{+w}{ }\PYG{k}{typename}\PYG{+w}{ }\PYG{n+nc}{F}\PYG{o}{\PYGZgt{}}
\PYG{k+kt}{void}\PYG{+w}{ }\PYG{n}{parallel\PYGZus{}map}\PYG{p}{(}\PYG{n}{Iter}\PYG{+w}{ }\PYG{n}{begin}\PYG{p}{,}\PYG{+w}{ }\PYG{n}{Iter}\PYG{+w}{ }\PYG{n}{end}\PYG{p}{,}\PYG{+w}{ }\PYG{n}{F}\PYG{+w}{ }\PYG{n}{f}\PYG{p}{)}
\PYG{p}{\PYGZob{}}
\PYG{+w}{  }\PYG{k}{auto}\PYG{+w}{ }\PYG{n}{size}\PYG{+w}{ }\PYG{o}{=}\PYG{+w}{ }\PYG{n}{std}\PYG{o}{:}\PYG{o}{:}\PYG{n}{distance}\PYG{p}{(}\PYG{n}{begin}\PYG{p}{,}\PYG{+w}{ }\PYG{n}{end}\PYG{p}{)}\PYG{p}{;}
\PYG{+w}{  }\PYG{k}{if}\PYG{p}{(}\PYG{n}{size}\PYG{+w}{ }\PYG{o}{\PYGZlt{}}\PYG{o}{=}\PYG{+w}{ }\PYG{l+m+mi}{10000}\PYG{p}{)}
\PYG{+w}{    }\PYG{n}{std}\PYG{o}{:}\PYG{o}{:}\PYG{n}{transform}\PYG{p}{(}\PYG{n}{begin}\PYG{p}{,}\PYG{+w}{ }\PYG{n}{end}\PYG{p}{,}\PYG{+w}{ }\PYG{n}{begin}\PYG{p}{,}\PYG{+w}{ }\PYG{n}{std}\PYG{o}{:}\PYG{o}{:}\PYG{n}{forward}\PYG{o}{\PYGZlt{}}\PYG{n}{F}\PYG{o}{\PYGZgt{}}\PYG{p}{(}\PYG{n}{f}\PYG{p}{)}\PYG{p}{)}\PYG{p}{;}
\PYG{+w}{  }\PYG{k}{else}
\PYG{+w}{  }\PYG{p}{\PYGZob{}}
\PYG{+w}{    }\PYG{k}{auto}\PYG{+w}{ }\PYG{n}{no\PYGZus{}of\PYGZus{}tasks}\PYG{+w}{ }\PYG{o}{=}\PYG{+w}{ }\PYG{n}{get\PYGZus{}no\PYGZus{}of\PYGZus{}threads}\PYG{p}{(}\PYG{p}{)}\PYG{p}{;}
\PYG{+w}{    }\PYG{k}{auto}\PYG{+w}{ }\PYG{n}{part}\PYG{+w}{ }\PYG{o}{=}\PYG{+w}{ }\PYG{n}{size}\PYG{+w}{ }\PYG{o}{/}\PYG{+w}{ }\PYG{n}{no\PYGZus{}of\PYGZus{}tasks}\PYG{p}{;}
\PYG{+w}{    }\PYG{k}{auto}\PYG{+w}{ }\PYG{n}{last}\PYG{+w}{ }\PYG{o}{=}\PYG{+w}{ }\PYG{n}{begin}\PYG{p}{;}
\PYG{n}{std}\PYG{o}{:}\PYG{o}{:}\PYG{n}{vector}\PYG{o}{\PYGZlt{}}\PYG{n}{std}\PYG{o}{:}\PYG{o}{:}\PYG{n}{future}\PYG{o}{\PYGZlt{}}\PYG{k+kt}{void}\PYG{o}{\PYGZgt{}}\PYG{o}{\PYGZgt{}}\PYG{+w}{ }\PYG{n}{tasks}\PYG{p}{;}
\PYG{k}{for}\PYG{p}{(}\PYG{k+kt}{unsigned}\PYG{+w}{ }\PYG{n}{i}\PYG{+w}{ }\PYG{o}{=}\PYG{+w}{ }\PYG{l+m+mi}{0}\PYG{p}{;}\PYG{+w}{ }\PYG{n}{i}\PYG{+w}{ }\PYG{o}{\PYGZlt{}}\PYG{+w}{ }\PYG{n}{no\PYGZus{}of\PYGZus{}tasks}\PYG{p}{;}\PYG{+w}{ }\PYG{o}{+}\PYG{o}{+}\PYG{n}{i}\PYG{p}{)}
\PYG{p}{\PYGZob{}}
\PYG{+w}{  }\PYG{k}{if}\PYG{p}{(}\PYG{n}{i}\PYG{+w}{ }\PYG{o}{=}\PYG{o}{=}\PYG{+w}{ }\PYG{n}{no\PYGZus{}of\PYGZus{}tasks}\PYG{+w}{ }\PYG{o}{\PYGZhy{}}\PYG{+w}{ }\PYG{l+m+mi}{1}\PYG{p}{)}\PYG{+w}{ }\PYG{n}{last}\PYG{+w}{ }\PYG{o}{=}\PYG{+w}{ }\PYG{n}{end}\PYG{p}{;}
\PYG{+w}{  }\PYG{k}{else}\PYG{+w}{ }\PYG{n}{std}\PYG{o}{:}\PYG{o}{:}\PYG{n}{advance}\PYG{p}{(}\PYG{n}{last}\PYG{p}{,}\PYG{+w}{ }\PYG{n}{part}\PYG{p}{)}\PYG{p}{;}
\PYG{n}{tasks}\PYG{p}{.}\PYG{n}{emplace\PYGZus{}back}\PYG{p}{(}\PYG{n}{std}\PYG{o}{:}\PYG{o}{:}\PYG{n}{async}\PYG{p}{(}
\PYG{+w}{  }\PYG{n}{std}\PYG{o}{:}\PYG{o}{:}\PYG{n}{launch}\PYG{o}{:}\PYG{o}{:}\PYG{n}{async}\PYG{p}{,}
\PYG{p}{[}\PYG{o}{=}\PYG{p}{,}\PYG{o}{\PYGZam{}}\PYG{n}{f}\PYG{p}{]}\PYG{p}{\PYGZob{}}\PYG{n}{std}\PYG{o}{:}\PYG{o}{:}\PYG{n}{transform}\PYG{p}{(}\PYG{n}{begin}\PYG{p}{,}\PYG{+w}{ }\PYG{n}{last}\PYG{p}{,}\PYG{+w}{ }\PYG{n}{begin}\PYG{p}{,}
\PYG{+w}{                      }\PYG{n}{std}\PYG{o}{:}\PYG{o}{:}\PYG{n}{forward}\PYG{o}{\PYGZlt{}}\PYG{n}{F}\PYG{o}{\PYGZgt{}}\PYG{p}{(}\PYG{n}{f}\PYG{p}{)}\PYG{p}{)}\PYG{p}{;}\PYG{p}{\PYGZcb{}}\PYG{p}{)}\PYG{p}{)}\PYG{p}{;}
\PYG{+w}{  }\PYG{n}{begin}\PYG{+w}{ }\PYG{o}{=}\PYG{+w}{ }\PYG{n}{last}\PYG{p}{;}
\PYG{p}{\PYGZcb{}}
\PYG{+w}{    }\PYG{k}{for}\PYG{p}{(}\PYG{k}{auto}\PYG{+w}{ }\PYG{o}{\PYGZam{}}\PYG{+w}{ }\PYG{n}{t}\PYG{+w}{ }\PYG{o}{:}\PYG{+w}{ }\PYG{n}{tasks}\PYG{p}{)}\PYG{+w}{ }\PYG{n}{t}\PYG{p}{.}\PYG{n}{wait}\PYG{p}{(}\PYG{p}{)}\PYG{p}{;}
\PYG{+w}{  }\PYG{p}{\PYGZcb{}}
\PYG{p}{\PYGZcb{}}
\end{Verbatim}
\end{codeblockbg}


%%C080434 p

要实现左\texttt{fold}函数的并行版本，请按以下步骤操作：


%%C080435 p

定义一个函数模板，它接受指向某个范围的起始与结束迭代器、一个初始


%%C080436 p

值，以及一个应用到该范围各元素上的二元函数：


%%C080437 code

\begin{codeblockbg}
\begin{Verbatim}[commandchars=\\\{\},breaklines=true,breakanywhere=true,fontsize=\codefont,xleftmargin=17bp]
\PYG{k}{template}\PYG{+w}{ }\PYG{o}{\PYGZlt{}}\PYG{k}{typename}\PYG{+w}{ }\PYG{n+nc}{Iter}\PYG{p}{,}\PYG{+w}{ }\PYG{k}{typename}\PYG{+w}{ }\PYG{n+nc}{R}\PYG{p}{,}\PYG{+w}{ }\PYG{k}{typename}\PYG{+w}{ }\PYG{n+nc}{F}\PYG{o}{\PYGZgt{}}
\PYG{k}{auto}\PYG{+w}{ }\PYG{n}{parallel\PYGZus{}reduce}\PYG{p}{(}\PYG{n}{Iter}\PYG{+w}{ }\PYG{n}{begin}\PYG{p}{,}\PYG{+w}{ }\PYG{n}{Iter}\PYG{+w}{ }\PYG{n}{end}\PYG{p}{,}\PYG{+w}{ }\PYG{n}{R}\PYG{+w}{ }\PYG{n}{init}\PYG{p}{,}\PYG{+w}{ }\PYG{n}{F}\PYG{+w}{ }\PYG{n}{op}\PYG{p}{)}
\PYG{p}{\PYGZob{}}
\PYG{p}{\PYGZcb{}}
\end{Verbatim}
\end{codeblockbg}


%%C080438 p

检查范围的大小。如果元素个数小于


%%C080439 p

预先定义的阈值（本实现中该阈值为10,000），则以串行方式执行折叠：


%%C080440 code

\begin{codeblockbg}
\begin{Verbatim}[commandchars=\\\{\},breaklines=true,breakanywhere=true,fontsize=\codefont,xleftmargin=17bp]
\PYG{k}{auto}\PYG{+w}{ }\PYG{n}{size}\PYG{+w}{ }\PYG{o}{=}\PYG{+w}{ }\PYG{n}{std}\PYG{o}{:}\PYG{o}{:}\PYG{n}{distance}\PYG{p}{(}\PYG{n}{begin}\PYG{p}{,}\PYG{+w}{ }\PYG{n}{end}\PYG{p}{)}\PYG{p}{;}
\PYG{k}{if}\PYG{p}{(}\PYG{n}{size}\PYG{+w}{ }\PYG{o}{\PYGZlt{}}\PYG{o}{=}\PYG{+w}{ }\PYG{l+m+mi}{10000}\PYG{p}{)}
\PYG{+w}{  }\PYG{k}{return}\PYG{+w}{ }\PYG{n}{std}\PYG{o}{:}\PYG{o}{:}\PYG{n}{accumulate}\PYG{p}{(}\PYG{n}{begin}\PYG{p}{,}\PYG{+w}{ }\PYG{n}{end}\PYG{p}{,}\PYG{+w}{ }\PYG{n}{init}\PYG{p}{,}
\PYG{+w}{                         }\PYG{n}{std}\PYG{o}{:}\PYG{o}{:}\PYG{n}{forward}\PYG{o}{\PYGZlt{}}\PYG{n}{F}\PYG{o}{\PYGZgt{}}\PYG{p}{(}\PYG{n}{op}\PYG{p}{)}\PYG{p}{)}\PYG{p}{;}
\end{Verbatim}
\end{codeblockbg}


%%C080441 p

对于更大的范围，把工作拆分给多个任务，让每个任务折叠


%%C080442 p

范围中的一部分。这些部分不应重叠，以避免同步线程对共享数据的访问。结果可以通过传给异步函数的引用返回，从而避免同步：


%%C080443 code

\begin{codeblockbg}
\begin{Verbatim}[commandchars=\\\{\},breaklines=true,breakanywhere=true,fontsize=\codefont,xleftmargin=17bp]
\PYG{k}{else}
\PYG{p}{\PYGZob{}}
\PYG{+w}{  }\PYG{k}{auto}\PYG{+w}{ }\PYG{n}{no\PYGZus{}of\PYGZus{}tasks}\PYG{+w}{ }\PYG{o}{=}\PYG{+w}{ }\PYG{n}{get\PYGZus{}no\PYGZus{}of\PYGZus{}threads}\PYG{p}{(}\PYG{p}{)}\PYG{p}{;}
\PYG{+w}{  }\PYG{k}{auto}\PYG{+w}{ }\PYG{n}{part}\PYG{+w}{ }\PYG{o}{=}\PYG{+w}{ }\PYG{n}{size}\PYG{+w}{ }\PYG{o}{/}\PYG{+w}{ }\PYG{n}{no\PYGZus{}of\PYGZus{}tasks}\PYG{p}{;}
\PYG{+w}{  }\PYG{k}{auto}\PYG{+w}{ }\PYG{n}{last}\PYG{+w}{ }\PYG{o}{=}\PYG{+w}{ }\PYG{n}{begin}\PYG{p}{;}
\PYG{+w}{  }\PYG{c+c1}{// 后接第 4、5 步}
\PYG{p}{\PYGZcb{}}
\end{Verbatim}
\end{codeblockbg}


%%C080444 p

启动这些异步函数，并在其中每一个上执行


%%C080445 p

串行版本的折叠：


%%C080446 code

\begin{codeblockbg}
\begin{Verbatim}[commandchars=\\\{\},breaklines=true,breakanywhere=true,fontsize=\codefont,xleftmargin=17bp]
\PYG{n}{std}\PYG{o}{:}\PYG{o}{:}\PYG{n}{vector}\PYG{o}{\PYGZlt{}}\PYG{n}{std}\PYG{o}{:}\PYG{o}{:}\PYG{n}{future}\PYG{o}{\PYGZlt{}}\PYG{n}{R}\PYG{o}{\PYGZgt{}}\PYG{o}{\PYGZgt{}}\PYG{+w}{ }\PYG{n}{tasks}\PYG{p}{;}
\PYG{k}{for}\PYG{p}{(}\PYG{k+kt}{unsigned}\PYG{+w}{ }\PYG{n}{i}\PYG{+w}{ }\PYG{o}{=}\PYG{+w}{ }\PYG{l+m+mi}{0}\PYG{p}{;}\PYG{+w}{ }\PYG{n}{i}\PYG{+w}{ }\PYG{o}{\PYGZlt{}}\PYG{+w}{ }\PYG{n}{no\PYGZus{}of\PYGZus{}tasks}\PYG{p}{;}\PYG{+w}{ }\PYG{o}{+}\PYG{o}{+}\PYG{n}{i}\PYG{p}{)}
\PYG{p}{\PYGZob{}}
\PYG{+w}{  }\PYG{k}{if}\PYG{p}{(}\PYG{n}{i}\PYG{+w}{ }\PYG{o}{=}\PYG{o}{=}\PYG{+w}{ }\PYG{n}{no\PYGZus{}of\PYGZus{}tasks}\PYG{+w}{ }\PYG{o}{\PYGZhy{}}\PYG{+w}{ }\PYG{l+m+mi}{1}\PYG{p}{)}\PYG{+w}{ }\PYG{n}{last}\PYG{+w}{ }\PYG{o}{=}\PYG{+w}{ }\PYG{n}{end}\PYG{p}{;}
\PYG{+w}{  }\PYG{k}{else}\PYG{+w}{ }\PYG{n}{std}\PYG{o}{:}\PYG{o}{:}\PYG{n}{advance}\PYG{p}{(}\PYG{n}{last}\PYG{p}{,}\PYG{+w}{ }\PYG{n}{part}\PYG{p}{)}\PYG{p}{;}
\PYG{+w}{  }\PYG{n}{tasks}\PYG{p}{.}\PYG{n}{emplace\PYGZus{}back}\PYG{p}{(}
\PYG{+w}{    }\PYG{n}{std}\PYG{o}{:}\PYG{o}{:}\PYG{n}{async}\PYG{p}{(}
\PYG{+w}{      }\PYG{n}{std}\PYG{o}{:}\PYG{o}{:}\PYG{n}{launch}\PYG{o}{:}\PYG{o}{:}\PYG{n}{async}\PYG{p}{,}
\PYG{+w}{      }\PYG{p}{[}\PYG{o}{=}\PYG{p}{,}\PYG{o}{\PYGZam{}}\PYG{n}{op}\PYG{p}{]}\PYG{p}{\PYGZob{}}\PYG{k}{return}\PYG{+w}{ }\PYG{n}{std}\PYG{o}{:}\PYG{o}{:}\PYG{n}{accumulate}\PYG{p}{(}
\PYG{+w}{                        }\PYG{n}{begin}\PYG{p}{,}\PYG{+w}{ }\PYG{n}{last}\PYG{p}{,}\PYG{+w}{ }\PYG{n}{R}\PYG{p}{\PYGZob{}}\PYG{p}{\PYGZcb{}}\PYG{p}{,}
\PYG{+w}{                        }\PYG{n}{std}\PYG{o}{:}\PYG{o}{:}\PYG{n}{forward}\PYG{o}{\PYGZlt{}}\PYG{n}{F}\PYG{o}{\PYGZgt{}}\PYG{p}{(}\PYG{n}{op}\PYG{p}{)}\PYG{p}{)}\PYG{p}{;}\PYG{p}{\PYGZcb{}}\PYG{p}{)}\PYG{p}{)}\PYG{p}{;}
\PYG{+w}{  }\PYG{n}{begin}\PYG{+w}{ }\PYG{o}{=}\PYG{+w}{ }\PYG{n}{last}\PYG{p}{;}
\PYG{p}{\PYGZcb{}}
\end{Verbatim}
\end{codeblockbg}


%%C080447 p

等待所有异步函数执行完毕，并把各部分结果


%%C080448 p

折叠成最终结果：


%%C080449 code

\begin{codeblockbg}
\begin{Verbatim}[commandchars=\\\{\},breaklines=true,breakanywhere=true,fontsize=\codefont,xleftmargin=17bp]
\PYG{n}{std}\PYG{o}{:}\PYG{o}{:}\PYG{n}{vector}\PYG{o}{\PYGZlt{}}\PYG{n}{R}\PYG{o}{\PYGZgt{}}\PYG{+w}{ }\PYG{n}{values}\PYG{p}{;}
\PYG{k}{for}\PYG{p}{(}\PYG{k}{auto}\PYG{+w}{ }\PYG{o}{\PYGZam{}}\PYG{+w}{ }\PYG{n}{t}\PYG{+w}{ }\PYG{o}{:}\PYG{+w}{ }\PYG{n}{tasks}\PYG{p}{)}
\PYG{+w}{  }\PYG{n}{values}\PYG{p}{.}\PYG{n}{push\PYGZus{}back}\PYG{p}{(}\PYG{n}{t}\PYG{p}{.}\PYG{n}{get}\PYG{p}{(}\PYG{p}{)}\PYG{p}{)}\PYG{p}{;}
\PYG{k}{return}\PYG{+w}{ }\PYG{n}{std}\PYG{o}{:}\PYG{o}{:}\PYG{n}{accumulate}\PYG{p}{(}\PYG{n}{std}\PYG{o}{:}\PYG{o}{:}\PYG{n}{begin}\PYG{p}{(}\PYG{n}{values}\PYG{p}{)}\PYG{p}{,}\PYG{+w}{ }\PYG{n}{std}\PYG{o}{:}\PYG{o}{:}\PYG{n}{end}\PYG{p}{(}\PYG{n}{values}\PYG{p}{)}\PYG{p}{,}
\PYG{+w}{                       }\PYG{n}{init}\PYG{p}{,}\PYG{+w}{ }\PYG{n}{std}\PYG{o}{:}\PYG{o}{:}\PYG{n}{forward}\PYG{o}{\PYGZlt{}}\PYG{n}{F}\PYG{o}{\PYGZgt{}}\PYG{p}{(}\PYG{n}{op}\PYG{p}{)}\PYG{p}{)}\PYG{p}{;}
\end{Verbatim}
\end{codeblockbg}


%%C080450 p

把这些步骤组合起来，就得到下面的实现：


%%C080451 code

\begin{codeblockbg}
\begin{Verbatim}[commandchars=\\\{\},breaklines=true,breakanywhere=true,fontsize=\codefont,xleftmargin=17bp]
\PYG{k}{template}\PYG{+w}{ }\PYG{o}{\PYGZlt{}}\PYG{k}{typename}\PYG{+w}{ }\PYG{n+nc}{Iter}\PYG{p}{,}\PYG{+w}{ }\PYG{k}{typename}\PYG{+w}{ }\PYG{n+nc}{R}\PYG{p}{,}\PYG{+w}{ }\PYG{k}{typename}\PYG{+w}{ }\PYG{n+nc}{F}\PYG{o}{\PYGZgt{}}
\PYG{k}{auto}\PYG{+w}{ }\PYG{n}{parallel\PYGZus{}reduce}\PYG{p}{(}\PYG{n}{Iter}\PYG{+w}{ }\PYG{n}{begin}\PYG{p}{,}\PYG{+w}{ }\PYG{n}{Iter}\PYG{+w}{ }\PYG{n}{end}\PYG{p}{,}\PYG{+w}{ }\PYG{n}{R}\PYG{+w}{ }\PYG{n}{init}\PYG{p}{,}\PYG{+w}{ }\PYG{n}{F}\PYG{+w}{ }\PYG{n}{op}\PYG{p}{)}
\PYG{p}{\PYGZob{}}
\PYG{+w}{  }\PYG{k}{auto}\PYG{+w}{ }\PYG{n}{size}\PYG{+w}{ }\PYG{o}{=}\PYG{+w}{ }\PYG{n}{std}\PYG{o}{:}\PYG{o}{:}\PYG{n}{distance}\PYG{p}{(}\PYG{n}{begin}\PYG{p}{,}\PYG{+w}{ }\PYG{n}{end}\PYG{p}{)}\PYG{p}{;}
\PYG{k}{if}\PYG{p}{(}\PYG{n}{size}\PYG{+w}{ }\PYG{o}{\PYGZlt{}}\PYG{o}{=}\PYG{+w}{ }\PYG{l+m+mi}{10000}\PYG{p}{)}
\PYG{+w}{  }\PYG{k}{return}\PYG{+w}{ }\PYG{n}{std}\PYG{o}{:}\PYG{o}{:}\PYG{n}{accumulate}\PYG{p}{(}\PYG{n}{begin}\PYG{p}{,}\PYG{+w}{ }\PYG{n}{end}\PYG{p}{,}\PYG{+w}{ }\PYG{n}{init}\PYG{p}{,}\PYG{+w}{ }\PYG{n}{std}\PYG{o}{:}\PYG{o}{:}\PYG{n}{forward}\PYG{o}{\PYGZlt{}}\PYG{n}{F}\PYG{o}{\PYGZgt{}}\PYG{p}{(}\PYG{n}{op}\PYG{p}{)}\PYG{p}{)}\PYG{p}{;}
\PYG{k}{else}
\PYG{p}{\PYGZob{}}
\PYG{+w}{  }\PYG{k}{auto}\PYG{+w}{ }\PYG{n}{no\PYGZus{}of\PYGZus{}tasks}\PYG{+w}{ }\PYG{o}{=}\PYG{+w}{ }\PYG{n}{get\PYGZus{}no\PYGZus{}of\PYGZus{}threads}\PYG{p}{(}\PYG{p}{)}\PYG{p}{;}
\PYG{+w}{  }\PYG{k}{auto}\PYG{+w}{ }\PYG{n}{part}\PYG{+w}{ }\PYG{o}{=}\PYG{+w}{ }\PYG{n}{size}\PYG{+w}{ }\PYG{o}{/}\PYG{+w}{ }\PYG{n}{no\PYGZus{}of\PYGZus{}tasks}\PYG{p}{;}
\PYG{+w}{  }\PYG{k}{auto}\PYG{+w}{ }\PYG{n}{last}\PYG{+w}{ }\PYG{o}{=}\PYG{+w}{ }\PYG{n}{begin}\PYG{p}{;}
\PYG{n}{std}\PYG{o}{:}\PYG{o}{:}\PYG{n}{vector}\PYG{o}{\PYGZlt{}}\PYG{n}{std}\PYG{o}{:}\PYG{o}{:}\PYG{n}{future}\PYG{o}{\PYGZlt{}}\PYG{n}{R}\PYG{o}{\PYGZgt{}}\PYG{o}{\PYGZgt{}}\PYG{+w}{ }\PYG{n}{tasks}\PYG{p}{;}
\PYG{k}{for}\PYG{p}{(}\PYG{k+kt}{unsigned}\PYG{+w}{ }\PYG{n}{i}\PYG{+w}{ }\PYG{o}{=}\PYG{+w}{ }\PYG{l+m+mi}{0}\PYG{p}{;}\PYG{+w}{ }\PYG{n}{i}\PYG{+w}{ }\PYG{o}{\PYGZlt{}}\PYG{+w}{ }\PYG{n}{no\PYGZus{}of\PYGZus{}tasks}\PYG{p}{;}\PYG{+w}{ }\PYG{o}{+}\PYG{o}{+}\PYG{n}{i}\PYG{p}{)}
\PYG{p}{\PYGZob{}}
\PYG{+w}{  }\PYG{k}{if}\PYG{p}{(}\PYG{n}{i}\PYG{+w}{ }\PYG{o}{=}\PYG{o}{=}\PYG{+w}{ }\PYG{n}{no\PYGZus{}of\PYGZus{}tasks}\PYG{+w}{ }\PYG{o}{\PYGZhy{}}\PYG{+w}{ }\PYG{l+m+mi}{1}\PYG{p}{)}\PYG{+w}{ }\PYG{n}{last}\PYG{+w}{ }\PYG{o}{=}\PYG{+w}{ }\PYG{n}{end}\PYG{p}{;}
\PYG{+w}{  }\PYG{k}{else}\PYG{+w}{ }\PYG{n}{std}\PYG{o}{:}\PYG{o}{:}\PYG{n}{advance}\PYG{p}{(}\PYG{n}{last}\PYG{p}{,}\PYG{+w}{ }\PYG{n}{part}\PYG{p}{)}\PYG{p}{;}
\PYG{n}{tasks}\PYG{p}{.}\PYG{n}{emplace\PYGZus{}back}\PYG{p}{(}
\PYG{+w}{  }\PYG{n}{std}\PYG{o}{:}\PYG{o}{:}\PYG{n}{async}\PYG{p}{(}
\PYG{+w}{    }\PYG{n}{std}\PYG{o}{:}\PYG{o}{:}\PYG{n}{launch}\PYG{o}{:}\PYG{o}{:}\PYG{n}{async}\PYG{p}{,}
\PYG{+w}{    }\PYG{p}{[}\PYG{o}{=}\PYG{p}{,}\PYG{o}{\PYGZam{}}\PYG{n}{op}\PYG{p}{]}\PYG{p}{\PYGZob{}}\PYG{k}{return}\PYG{+w}{ }\PYG{n}{std}\PYG{o}{:}\PYG{o}{:}\PYG{n}{accumulate}\PYG{p}{(}
\PYG{n}{begin}\PYG{p}{,}\PYG{+w}{ }\PYG{n}{last}\PYG{p}{,}\PYG{+w}{ }\PYG{n}{R}\PYG{p}{\PYGZob{}}\PYG{p}{\PYGZcb{}}\PYG{p}{,}
\PYG{n}{std}\PYG{o}{:}\PYG{o}{:}\PYG{n}{forward}\PYG{o}{\PYGZlt{}}\PYG{n}{F}\PYG{o}{\PYGZgt{}}\PYG{p}{(}\PYG{n}{op}\PYG{p}{)}\PYG{p}{)}\PYG{p}{;}\PYG{p}{\PYGZcb{}}\PYG{p}{)}\PYG{p}{)}\PYG{p}{;}
\PYG{+w}{  }\PYG{n}{begin}\PYG{+w}{ }\PYG{o}{=}\PYG{+w}{ }\PYG{n}{last}\PYG{p}{;}
\PYG{p}{\PYGZcb{}}
\PYG{n}{std}\PYG{o}{:}\PYG{o}{:}\PYG{n}{vector}\PYG{o}{\PYGZlt{}}\PYG{n}{R}\PYG{o}{\PYGZgt{}}\PYG{+w}{ }\PYG{n}{values}\PYG{p}{;}
\PYG{k}{for}\PYG{p}{(}\PYG{k}{auto}\PYG{+w}{ }\PYG{o}{\PYGZam{}}\PYG{+w}{ }\PYG{n}{t}\PYG{+w}{ }\PYG{o}{:}\PYG{+w}{ }\PYG{n}{tasks}\PYG{p}{)}
\PYG{+w}{  }\PYG{n}{values}\PYG{p}{.}\PYG{n}{push\PYGZus{}back}\PYG{p}{(}\PYG{n}{t}\PYG{p}{.}\PYG{n}{get}\PYG{p}{(}\PYG{p}{)}\PYG{p}{)}\PYG{p}{;}
\PYG{+w}{    }\PYG{k}{return}\PYG{+w}{ }\PYG{n}{std}\PYG{o}{:}\PYG{o}{:}\PYG{n}{accumulate}\PYG{p}{(}\PYG{n}{std}\PYG{o}{:}\PYG{o}{:}\PYG{n}{begin}\PYG{p}{(}\PYG{n}{values}\PYG{p}{)}\PYG{p}{,}\PYG{+w}{ }\PYG{n}{std}\PYG{o}{:}\PYG{o}{:}\PYG{n}{end}\PYG{p}{(}\PYG{n}{values}\PYG{p}{)}\PYG{p}{,}
\PYG{+w}{                           }\PYG{n}{init}\PYG{p}{,}\PYG{+w}{ }\PYG{n}{std}\PYG{o}{:}\PYG{o}{:}\PYG{n}{forward}\PYG{o}{\PYGZlt{}}\PYG{n}{F}\PYG{o}{\PYGZgt{}}\PYG{p}{(}\PYG{n}{op}\PYG{p}{)}\PYG{p}{)}\PYG{p}{;}
\PYG{+w}{  }\PYG{p}{\PYGZcb{}}
\PYG{p}{\PYGZcb{}}
\end{Verbatim}
\end{codeblockbg}


%%C080452 h3

\subsection{工作原理}


%%C080453 p

这里给出的实现与上一案例中的实现只有细微差别。线程被替换为异步函数，用\texttt{std::\allowbreak{}async()}启动，结果通过返回的\texttt{std::\allowbreak{}future}获取。并发启动的异步函数个数等于实现所能支持的线程数，该数目由静态方法\myttbr{std::\allowbreak{}thread::\allowbreak{}hardware\_\allowbreak{}concurrency()}返回，但这个值只是一个提示，不应认为它十分可靠。


%%C080454 p

采用这种做法主要有两个原因：


%%C080455 li

\begin{itemize}
\item 看看一个为并行执行而用线程实现的函数，如何改造成使用异步函数，从而避开线程的底层细节。


%%C080456 li

\item 启动与所支持线程数相同的异步函数，有可能让每个线程运行一个函数；这样可以给并行函数带来最快的执行时间，因为上下文切换与等待时间的开销最小。

\end{itemize}


%%C080457 p

我们可以用与上一案例相同的方法测试新的\texttt{map}与\texttt{fold}实现的性能：


%%C080458 code

\begin{codeblockbg}
\begin{Verbatim}[commandchars=\\\{\},breaklines=true,breakanywhere=true,fontsize=\codefont,xleftmargin=17bp]
\PYG{n}{std}\PYG{o}{:}\PYG{o}{:}\PYG{n}{vector}\PYG{o}{\PYGZlt{}}\PYG{k+kt}{int}\PYG{o}{\PYGZgt{}}\PYG{+w}{ }\PYG{n}{sizes}
\PYG{p}{\PYGZob{}}
\PYG{+w}{  }\PYG{l+m+mi}{10000}\PYG{p}{,}\PYG{+w}{ }\PYG{l+m+mi}{100000}\PYG{p}{,}\PYG{+w}{ }\PYG{l+m+mi}{500000}\PYG{p}{,}
\PYG{+w}{  }\PYG{l+m+mi}{1000000}\PYG{p}{,}\PYG{+w}{ }\PYG{l+m+mi}{2000000}\PYG{p}{,}\PYG{+w}{ }\PYG{l+m+mi}{5000000}\PYG{p}{,}
\PYG{+w}{  }\PYG{l+m+mi}{10000000}\PYG{p}{,}\PYG{+w}{ }\PYG{l+m+mi}{25000000}\PYG{p}{,}\PYG{+w}{ }\PYG{l+m+mi}{50000000}
\PYG{p}{\PYGZcb{}}\PYG{p}{;}
\PYG{n}{std}\PYG{o}{:}\PYG{o}{:}\PYG{n}{cout}
\PYG{+w}{ }\PYG{o}{\PYGZlt{}}\PYG{o}{\PYGZlt{}}\PYG{+w}{ }\PYG{n}{std}\PYG{o}{:}\PYG{o}{:}\PYG{n}{right}\PYG{+w}{ }\PYG{o}{\PYGZlt{}}\PYG{o}{\PYGZlt{}}\PYG{+w}{ }\PYG{n}{std}\PYG{o}{:}\PYG{o}{:}\PYG{n}{setw}\PYG{p}{(}\PYG{l+m+mi}{8}\PYG{p}{)}\PYG{+w}{ }\PYG{o}{\PYGZlt{}}\PYG{o}{\PYGZlt{}}\PYG{+w}{ }\PYG{n}{std}\PYG{o}{:}\PYG{o}{:}\PYG{n}{setfill}\PYG{p}{(}\PYG{l+s+sc}{\PYGZsq{}}\PYG{l+s+sc}{ }\PYG{l+s+sc}{\PYGZsq{}}\PYG{p}{)}\PYG{+w}{ }\PYG{o}{\PYGZlt{}}\PYG{o}{\PYGZlt{}}\PYG{+w}{ }\PYG{l+s}{\PYGZdq{}}\PYG{l+s}{size}\PYG{l+s}{\PYGZdq{}}
\PYG{+w}{ }\PYG{o}{\PYGZlt{}}\PYG{o}{\PYGZlt{}}\PYG{+w}{ }\PYG{n}{std}\PYG{o}{:}\PYG{o}{:}\PYG{n}{right}\PYG{+w}{ }\PYG{o}{\PYGZlt{}}\PYG{o}{\PYGZlt{}}\PYG{+w}{ }\PYG{n}{std}\PYG{o}{:}\PYG{o}{:}\PYG{n}{setw}\PYG{p}{(}\PYG{l+m+mi}{8}\PYG{p}{)}\PYG{+w}{ }\PYG{o}{\PYGZlt{}}\PYG{o}{\PYGZlt{}}\PYG{+w}{ }\PYG{l+s}{\PYGZdq{}}\PYG{l+s}{s map}\PYG{l+s}{\PYGZdq{}}
\PYG{+w}{ }\PYG{o}{\PYGZlt{}}\PYG{o}{\PYGZlt{}}\PYG{+w}{ }\PYG{n}{std}\PYG{o}{:}\PYG{o}{:}\PYG{n}{right}\PYG{+w}{ }\PYG{o}{\PYGZlt{}}\PYG{o}{\PYGZlt{}}\PYG{+w}{ }\PYG{n}{std}\PYG{o}{:}\PYG{o}{:}\PYG{n}{setw}\PYG{p}{(}\PYG{l+m+mi}{8}\PYG{p}{)}\PYG{+w}{ }\PYG{o}{\PYGZlt{}}\PYG{o}{\PYGZlt{}}\PYG{+w}{ }\PYG{l+s}{\PYGZdq{}}\PYG{l+s}{p map}\PYG{l+s}{\PYGZdq{}}
\PYG{+w}{ }\PYG{o}{\PYGZlt{}}\PYG{o}{\PYGZlt{}}\PYG{+w}{ }\PYG{n}{std}\PYG{o}{:}\PYG{o}{:}\PYG{n}{right}\PYG{+w}{ }\PYG{o}{\PYGZlt{}}\PYG{o}{\PYGZlt{}}\PYG{+w}{ }\PYG{n}{std}\PYG{o}{:}\PYG{o}{:}\PYG{n}{setw}\PYG{p}{(}\PYG{l+m+mi}{8}\PYG{p}{)}\PYG{+w}{ }\PYG{o}{\PYGZlt{}}\PYG{o}{\PYGZlt{}}\PYG{+w}{ }\PYG{l+s}{\PYGZdq{}}\PYG{l+s}{s fold}\PYG{l+s}{\PYGZdq{}}
\PYG{+w}{ }\PYG{o}{\PYGZlt{}}\PYG{o}{\PYGZlt{}}\PYG{+w}{ }\PYG{n}{std}\PYG{o}{:}\PYG{o}{:}\PYG{n}{right}\PYG{+w}{ }\PYG{o}{\PYGZlt{}}\PYG{o}{\PYGZlt{}}\PYG{+w}{ }\PYG{n}{std}\PYG{o}{:}\PYG{o}{:}\PYG{n}{setw}\PYG{p}{(}\PYG{l+m+mi}{8}\PYG{p}{)}\PYG{+w}{ }\PYG{o}{\PYGZlt{}}\PYG{o}{\PYGZlt{}}\PYG{+w}{ }\PYG{l+s}{\PYGZdq{}}\PYG{l+s}{p fold}\PYG{l+s}{\PYGZdq{}}
\PYG{+w}{ }\PYG{o}{\PYGZlt{}}\PYG{o}{\PYGZlt{}}\PYG{+w}{ }\PYG{n}{std}\PYG{o}{:}\PYG{o}{:}\PYG{n}{endl}\PYG{p}{;}
\PYG{k}{for}\PYG{p}{(}\PYG{k}{auto}\PYG{+w}{ }\PYG{k}{const}\PYG{+w}{ }\PYG{n}{size}\PYG{+w}{ }\PYG{o}{:}\PYG{+w}{ }\PYG{n}{sizes}\PYG{p}{)}
\PYG{p}{\PYGZob{}}
\PYG{+w}{  }\PYG{n}{std}\PYG{o}{:}\PYG{o}{:}\PYG{n}{vector}\PYG{o}{\PYGZlt{}}\PYG{k+kt}{int}\PYG{o}{\PYGZgt{}}\PYG{+w}{ }\PYG{n}{v}\PYG{p}{(}\PYG{n}{size}\PYG{p}{)}\PYG{p}{;}
\PYG{+w}{  }\PYG{n}{std}\PYG{o}{:}\PYG{o}{:}\PYG{n}{iota}\PYG{p}{(}\PYG{n}{std}\PYG{o}{:}\PYG{o}{:}\PYG{n}{begin}\PYG{p}{(}\PYG{n}{v}\PYG{p}{)}\PYG{p}{,}\PYG{+w}{ }\PYG{n}{std}\PYG{o}{:}\PYG{o}{:}\PYG{n}{end}\PYG{p}{(}\PYG{n}{v}\PYG{p}{)}\PYG{p}{,}\PYG{+w}{ }\PYG{l+m+mi}{1}\PYG{p}{)}\PYG{p}{;}
\PYG{k}{auto}\PYG{+w}{ }\PYG{n}{v1}\PYG{+w}{ }\PYG{o}{=}\PYG{+w}{ }\PYG{n}{v}\PYG{p}{;}
\PYG{k}{auto}\PYG{+w}{ }\PYG{n}{s1}\PYG{+w}{ }\PYG{o}{=}\PYG{+w}{ }\PYG{l+m+mf}{0L}\PYG{n}{L}\PYG{p}{;}
\PYG{k}{auto}\PYG{+w}{ }\PYG{n}{tsm}\PYG{+w}{ }\PYG{o}{=}\PYG{+w}{ }\PYG{n}{perf\PYGZus{}timer}\PYG{o}{\PYGZlt{}}\PYG{o}{\PYGZgt{}}\PYG{o}{:}\PYG{o}{:}\PYG{n}{duration}\PYG{p}{(}\PYG{p}{[}\PYG{o}{\PYGZam{}}\PYG{p}{]}\PYG{+w}{ }\PYG{p}{\PYGZob{}}
\PYG{+w}{  }\PYG{n}{std}\PYG{o}{:}\PYG{o}{:}\PYG{n}{transform}\PYG{p}{(}\PYG{n}{std}\PYG{o}{:}\PYG{o}{:}\PYG{n}{begin}\PYG{p}{(}\PYG{n}{v1}\PYG{p}{)}\PYG{p}{,}\PYG{+w}{ }\PYG{n}{std}\PYG{o}{:}\PYG{o}{:}\PYG{n}{end}\PYG{p}{(}\PYG{n}{v1}\PYG{p}{)}\PYG{p}{,}\PYG{+w}{ }\PYG{n}{std}\PYG{o}{:}\PYG{o}{:}\PYG{n}{begin}\PYG{p}{(}\PYG{n}{v1}\PYG{p}{)}\PYG{p}{,}
\PYG{+w}{                 }\PYG{p}{[}\PYG{p}{]}\PYG{p}{(}\PYG{k+kt}{int}\PYG{+w}{ }\PYG{k}{const}\PYG{+w}{ }\PYG{n}{i}\PYG{p}{)}\PYG{+w}{ }\PYG{p}{\PYGZob{}}\PYG{k}{return}\PYG{+w}{ }\PYG{n}{i}\PYG{+w}{ }\PYG{o}{+}\PYG{+w}{ }\PYG{n}{i}\PYG{p}{;}\PYG{+w}{ }\PYG{p}{\PYGZcb{}}\PYG{p}{)}\PYG{p}{;}\PYG{+w}{ }\PYG{p}{\PYGZcb{}}\PYG{p}{)}\PYG{p}{;}
\PYG{k}{auto}\PYG{+w}{ }\PYG{n}{tsf}\PYG{+w}{ }\PYG{o}{=}\PYG{+w}{ }\PYG{n}{perf\PYGZus{}timer}\PYG{o}{\PYGZlt{}}\PYG{o}{\PYGZgt{}}\PYG{o}{:}\PYG{o}{:}\PYG{n}{duration}\PYG{p}{(}\PYG{p}{[}\PYG{o}{\PYGZam{}}\PYG{p}{]}\PYG{+w}{ }\PYG{p}{\PYGZob{}}
\PYG{+w}{  }\PYG{n}{s1}\PYG{+w}{ }\PYG{o}{=}\PYG{+w}{ }\PYG{n}{std}\PYG{o}{:}\PYG{o}{:}\PYG{n}{accumulate}\PYG{p}{(}\PYG{n}{std}\PYG{o}{:}\PYG{o}{:}\PYG{n}{begin}\PYG{p}{(}\PYG{n}{v1}\PYG{p}{)}\PYG{p}{,}\PYG{+w}{ }\PYG{n}{std}\PYG{o}{:}\PYG{o}{:}\PYG{n}{end}\PYG{p}{(}\PYG{n}{v1}\PYG{p}{)}\PYG{p}{,}\PYG{+w}{ }\PYG{l+m+mf}{0L}\PYG{n}{L}\PYG{p}{,}
\PYG{+w}{                       }\PYG{n}{std}\PYG{o}{:}\PYG{o}{:}\PYG{n}{plus}\PYG{o}{\PYGZlt{}}\PYG{o}{\PYGZgt{}}\PYG{p}{(}\PYG{p}{)}\PYG{p}{)}\PYG{p}{;}\PYG{+w}{ }\PYG{p}{\PYGZcb{}}\PYG{p}{)}\PYG{p}{;}
\PYG{k}{auto}\PYG{+w}{ }\PYG{n}{v2}\PYG{+w}{ }\PYG{o}{=}\PYG{+w}{ }\PYG{n}{v}\PYG{p}{;}
\PYG{k}{auto}\PYG{+w}{ }\PYG{n}{s2}\PYG{+w}{ }\PYG{o}{=}\PYG{+w}{ }\PYG{l+m+mf}{0L}\PYG{n}{L}\PYG{p}{;}
\PYG{k}{auto}\PYG{+w}{ }\PYG{n}{tpm}\PYG{+w}{ }\PYG{o}{=}\PYG{+w}{ }\PYG{n}{perf\PYGZus{}timer}\PYG{o}{\PYGZlt{}}\PYG{o}{\PYGZgt{}}\PYG{o}{:}\PYG{o}{:}\PYG{n}{duration}\PYG{p}{(}\PYG{p}{[}\PYG{o}{\PYGZam{}}\PYG{p}{]}\PYG{+w}{ }\PYG{p}{\PYGZob{}}
\PYG{+w}{  }\PYG{n}{parallel\PYGZus{}map}\PYG{p}{(}\PYG{n}{std}\PYG{o}{:}\PYG{o}{:}\PYG{n}{begin}\PYG{p}{(}\PYG{n}{v2}\PYG{p}{)}\PYG{p}{,}\PYG{+w}{ }\PYG{n}{std}\PYG{o}{:}\PYG{o}{:}\PYG{n}{end}\PYG{p}{(}\PYG{n}{v2}\PYG{p}{)}\PYG{p}{,}
\PYG{+w}{               }\PYG{p}{[}\PYG{p}{]}\PYG{p}{(}\PYG{k+kt}{int}\PYG{+w}{ }\PYG{k}{const}\PYG{+w}{ }\PYG{n}{i}\PYG{p}{)}\PYG{+w}{ }\PYG{p}{\PYGZob{}}\PYG{k}{return}\PYG{+w}{ }\PYG{n}{i}\PYG{+w}{ }\PYG{o}{+}\PYG{+w}{ }\PYG{n}{i}\PYG{p}{;}\PYG{+w}{ }\PYG{p}{\PYGZcb{}}\PYG{p}{)}\PYG{p}{;}\PYG{+w}{ }\PYG{p}{\PYGZcb{}}\PYG{p}{)}\PYG{p}{;}
\PYG{k}{auto}\PYG{+w}{ }\PYG{n}{tpf}\PYG{+w}{ }\PYG{o}{=}\PYG{+w}{ }\PYG{n}{perf\PYGZus{}timer}\PYG{o}{\PYGZlt{}}\PYG{o}{\PYGZgt{}}\PYG{o}{:}\PYG{o}{:}\PYG{n}{duration}\PYG{p}{(}\PYG{p}{[}\PYG{o}{\PYGZam{}}\PYG{p}{]}\PYG{+w}{ }\PYG{p}{\PYGZob{}}
\PYG{+w}{  }\PYG{n}{s2}\PYG{+w}{ }\PYG{o}{=}\PYG{+w}{ }\PYG{n}{parallel\PYGZus{}reduce}\PYG{p}{(}\PYG{n}{std}\PYG{o}{:}\PYG{o}{:}\PYG{n}{begin}\PYG{p}{(}\PYG{n}{v2}\PYG{p}{)}\PYG{p}{,}\PYG{+w}{ }\PYG{n}{std}\PYG{o}{:}\PYG{o}{:}\PYG{n}{end}\PYG{p}{(}\PYG{n}{v2}\PYG{p}{)}\PYG{p}{,}\PYG{+w}{ }\PYG{l+m+mf}{0L}\PYG{n}{L}\PYG{p}{,}
\PYG{+w}{                       }\PYG{n}{std}\PYG{o}{:}\PYG{o}{:}\PYG{n}{plus}\PYG{o}{\PYGZlt{}}\PYG{o}{\PYGZgt{}}\PYG{p}{(}\PYG{p}{)}\PYG{p}{)}\PYG{p}{;}\PYG{+w}{ }\PYG{p}{\PYGZcb{}}\PYG{p}{)}\PYG{p}{;}
\PYG{n}{assert}\PYG{p}{(}\PYG{n}{v1}\PYG{+w}{ }\PYG{o}{=}\PYG{o}{=}\PYG{+w}{ }\PYG{n}{v2}\PYG{p}{)}\PYG{p}{;}
\PYG{n}{assert}\PYG{p}{(}\PYG{n}{s1}\PYG{+w}{ }\PYG{o}{=}\PYG{o}{=}\PYG{+w}{ }\PYG{n}{s2}\PYG{p}{)}\PYG{p}{;}
\PYG{n}{std}\PYG{o}{:}\PYG{o}{:}\PYG{n}{cout}
\PYG{+w}{  }\PYG{o}{\PYGZlt{}}\PYG{o}{\PYGZlt{}}\PYG{+w}{ }\PYG{n}{std}\PYG{o}{:}\PYG{o}{:}\PYG{n}{right}\PYG{+w}{ }\PYG{o}{\PYGZlt{}}\PYG{o}{\PYGZlt{}}\PYG{+w}{ }\PYG{n}{std}\PYG{o}{:}\PYG{o}{:}\PYG{n}{setw}\PYG{p}{(}\PYG{l+m+mi}{8}\PYG{p}{)}\PYG{+w}{ }\PYG{o}{\PYGZlt{}}\PYG{o}{\PYGZlt{}}\PYG{+w}{ }\PYG{n}{std}\PYG{o}{:}\PYG{o}{:}\PYG{n}{setfill}\PYG{p}{(}\PYG{l+s+sc}{\PYGZsq{}}\PYG{l+s+sc}{ }\PYG{l+s+sc}{\PYGZsq{}}\PYG{p}{)}\PYG{+w}{ }\PYG{o}{\PYGZlt{}}\PYG{o}{\PYGZlt{}}\PYG{+w}{ }\PYG{n}{size}
\PYG{+w}{  }\PYG{o}{\PYGZlt{}}\PYG{o}{\PYGZlt{}}\PYG{+w}{ }\PYG{n}{std}\PYG{o}{:}\PYG{o}{:}\PYG{n}{right}\PYG{+w}{ }\PYG{o}{\PYGZlt{}}\PYG{o}{\PYGZlt{}}\PYG{+w}{ }\PYG{n}{std}\PYG{o}{:}\PYG{o}{:}\PYG{n}{setw}\PYG{p}{(}\PYG{l+m+mi}{8}\PYG{p}{)}
\PYG{+w}{  }\PYG{o}{\PYGZlt{}}\PYG{o}{\PYGZlt{}}\PYG{+w}{ }\PYG{n}{std}\PYG{o}{:}\PYG{o}{:}\PYG{n}{chrono}\PYG{o}{:}\PYG{o}{:}\PYG{n}{duration}\PYG{o}{\PYGZlt{}}\PYG{k+kt}{double}\PYG{p}{,}\PYG{+w}{ }\PYG{n}{std}\PYG{o}{:}\PYG{o}{:}\PYG{n}{micro}\PYG{o}{\PYGZgt{}}\PYG{p}{(}\PYG{n}{tsm}\PYG{p}{)}\PYG{p}{.}\PYG{n}{count}\PYG{p}{(}\PYG{p}{)}
\PYG{+w}{  }\PYG{o}{\PYGZlt{}}\PYG{o}{\PYGZlt{}}\PYG{+w}{ }\PYG{n}{std}\PYG{o}{:}\PYG{o}{:}\PYG{n}{right}\PYG{+w}{ }\PYG{o}{\PYGZlt{}}\PYG{o}{\PYGZlt{}}\PYG{+w}{ }\PYG{n}{std}\PYG{o}{:}\PYG{o}{:}\PYG{n}{setw}\PYG{p}{(}\PYG{l+m+mi}{8}\PYG{p}{)}
\PYG{+w}{  }\PYG{o}{\PYGZlt{}}\PYG{o}{\PYGZlt{}}\PYG{+w}{ }\PYG{n}{std}\PYG{o}{:}\PYG{o}{:}\PYG{n}{chrono}\PYG{o}{:}\PYG{o}{:}\PYG{n}{duration}\PYG{o}{\PYGZlt{}}\PYG{k+kt}{double}\PYG{p}{,}\PYG{+w}{ }\PYG{n}{std}\PYG{o}{:}\PYG{o}{:}\PYG{n}{micro}\PYG{o}{\PYGZgt{}}\PYG{p}{(}\PYG{n}{tpm}\PYG{p}{)}\PYG{p}{.}\PYG{n}{count}\PYG{p}{(}\PYG{p}{)}
\PYG{+w}{  }\PYG{o}{\PYGZlt{}}\PYG{o}{\PYGZlt{}}\PYG{+w}{ }\PYG{n}{std}\PYG{o}{:}\PYG{o}{:}\PYG{n}{right}\PYG{+w}{ }\PYG{o}{\PYGZlt{}}\PYG{o}{\PYGZlt{}}\PYG{+w}{ }\PYG{n}{std}\PYG{o}{:}\PYG{o}{:}\PYG{n}{setw}\PYG{p}{(}\PYG{l+m+mi}{8}\PYG{p}{)}
\PYG{+w}{  }\PYG{o}{\PYGZlt{}}\PYG{o}{\PYGZlt{}}\PYG{+w}{ }\PYG{n}{std}\PYG{o}{:}\PYG{o}{:}\PYG{n}{chrono}\PYG{o}{:}\PYG{o}{:}\PYG{n}{duration}\PYG{o}{\PYGZlt{}}\PYG{k+kt}{double}\PYG{p}{,}\PYG{+w}{ }\PYG{n}{std}\PYG{o}{:}\PYG{o}{:}\PYG{n}{micro}\PYG{o}{\PYGZgt{}}\PYG{p}{(}\PYG{n}{tsf}\PYG{p}{)}\PYG{p}{.}\PYG{n}{count}\PYG{p}{(}\PYG{p}{)}
\PYG{+w}{  }\PYG{o}{\PYGZlt{}}\PYG{o}{\PYGZlt{}}\PYG{+w}{ }\PYG{n}{std}\PYG{o}{:}\PYG{o}{:}\PYG{n}{right}\PYG{+w}{ }\PYG{o}{\PYGZlt{}}\PYG{o}{\PYGZlt{}}\PYG{+w}{ }\PYG{n}{std}\PYG{o}{:}\PYG{o}{:}\PYG{n}{setw}\PYG{p}{(}\PYG{l+m+mi}{8}\PYG{p}{)}
\PYG{+w}{  }\PYG{o}{\PYGZlt{}}\PYG{o}{\PYGZlt{}}\PYG{+w}{ }\PYG{n}{std}\PYG{o}{:}\PYG{o}{:}\PYG{n}{chrono}\PYG{o}{:}\PYG{o}{:}\PYG{n}{duration}\PYG{o}{\PYGZlt{}}\PYG{k+kt}{double}\PYG{p}{,}\PYG{+w}{ }\PYG{n}{std}\PYG{o}{:}\PYG{o}{:}\PYG{n}{micro}\PYG{o}{\PYGZgt{}}\PYG{p}{(}\PYG{n}{tpf}\PYG{p}{)}\PYG{p}{.}\PYG{n}{count}\PYG{p}{(}\PYG{p}{)}
\PYG{+w}{  }\PYG{o}{\PYGZlt{}}\PYG{o}{\PYGZlt{}}\PYG{+w}{ }\PYG{n}{std}\PYG{o}{:}\PYG{o}{:}\PYG{n}{endl}\PYG{p}{;}
\PYG{p}{\PYGZcb{}}
\end{Verbatim}
\end{codeblockbg}


%%C080459 p

上面程序的一个可能输出如下，每次执行的结果会略有不同，不同机器之间的差别则可能很大：


%%C080460 code

\begin{codeblockbg}
\begin{Verbatim}[commandchars=\\\{\},breaklines=true,breakanywhere=true,fontsize=\codefont,xleftmargin=17bp]
\PYG{+w}{    }\PYG{n}{size}\PYG{+w}{   }\PYG{n}{s}\PYG{+w}{ }\PYG{n}{map}\PYG{+w}{   }\PYG{n}{p}\PYG{+w}{ }\PYG{n}{map}\PYG{+w}{  }\PYG{n}{s}\PYG{+w}{ }\PYG{n}{fold}\PYG{+w}{  }\PYG{n}{p}\PYG{+w}{ }\PYG{n}{fold}
\PYG{+w}{   }\PYG{l+m+mi}{10000}\PYG{+w}{      }\PYG{l+m+mi}{11}\PYG{+w}{      }\PYG{l+m+mi}{11}\PYG{+w}{      }\PYG{l+m+mi}{11}\PYG{+w}{      }\PYG{l+m+mi}{11}
\PYG{+w}{  }\PYG{l+m+mi}{100000}\PYG{+w}{     }\PYG{l+m+mi}{117}\PYG{+w}{     }\PYG{l+m+mi}{260}\PYG{+w}{     }\PYG{l+m+mi}{113}\PYG{+w}{      }\PYG{l+m+mi}{94}
\PYG{+w}{  }\PYG{l+m+mi}{500000}\PYG{+w}{     }\PYG{l+m+mi}{576}\PYG{+w}{     }\PYG{l+m+mi}{303}\PYG{+w}{     }\PYG{l+m+mi}{571}\PYG{+w}{     }\PYG{l+m+mi}{201}
\PYG{+w}{ }\PYG{l+m+mi}{1000000}\PYG{+w}{    }\PYG{l+m+mi}{1180}\PYG{+w}{     }\PYG{l+m+mi}{573}\PYG{+w}{    }\PYG{l+m+mi}{1165}\PYG{+w}{     }\PYG{l+m+mi}{283}
\PYG{+w}{ }\PYG{l+m+mi}{2000000}\PYG{+w}{    }\PYG{l+m+mi}{2371}\PYG{+w}{     }\PYG{l+m+mi}{911}\PYG{+w}{    }\PYG{l+m+mi}{2330}\PYG{+w}{     }\PYG{l+m+mi}{519}
\PYG{+w}{ }\PYG{l+m+mi}{5000000}\PYG{+w}{    }\PYG{l+m+mi}{5942}\PYG{+w}{    }\PYG{l+m+mi}{2144}\PYG{+w}{    }\PYG{l+m+mi}{5841}\PYG{+w}{    }\PYG{l+m+mi}{1886}
\PYG{l+m+mi}{10000000}\PYG{+w}{   }\PYG{l+m+mi}{11954}\PYG{+w}{    }\PYG{l+m+mi}{4999}\PYG{+w}{   }\PYG{l+m+mi}{11643}\PYG{+w}{    }\PYG{l+m+mi}{2871}
\PYG{l+m+mi}{25000000}\PYG{+w}{   }\PYG{l+m+mi}{30525}\PYG{+w}{   }\PYG{l+m+mi}{11737}\PYG{+w}{   }\PYG{l+m+mi}{29053}\PYG{+w}{    }\PYG{l+m+mi}{9048}
\PYG{l+m+mi}{50000000}\PYG{+w}{   }\PYG{l+m+mi}{59665}\PYG{+w}{   }\PYG{l+m+mi}{22216}\PYG{+w}{   }\PYG{l+m+mi}{58689}\PYG{+w}{   }\PYG{l+m+mi}{12942}
\end{Verbatim}
\end{codeblockbg}


%%C080461 p

与线程方案的图示类似，并行\texttt{map}与\texttt{fold}实现的加速比如下图所示。负值表示串行版本更快：


%%C080462 img

\myFigure{images/ch08/fig-8-2.png}{}

%%C080463 p

把它与使用线程的并行版本的结果相比较，我们会发现任务版本的执行时间更短，加速效果相当明显，尤其是\texttt{fold}函数。下图给出任务实现相对于线程实现的加速比。图中数值小于1表示线程实现更快：


%%C080464 img

\myFigure{images/ch08/fig-8-3.png}{}


%%C080465 h3

\subsection{更多内容}


%%C080466 p

前面给出的实现只是并行化\texttt{map}与\texttt{fold}函数的一种可行做法。另一种可选方案采用以下策略：


%%C080467 li

\begin{itemize}
\item 把要处理的范围划分成两个相等的部分。


%%C080468 li

\item 异步地递归调用并行函数，处理范围的第一部分。


%%C080469 li

\item 同步地递归调用并行函数，处理范围的第二部分。


%%C080470 li

\item 同步递归调用结束后，还要等待异步递归调用也结束，然后才结束执行。

\end{itemize}


%%C080471 p

这种分治算法可能会创建大量任务。视范围大小而定，异步调用的数量可能远远超过线程数量，这时会出现大量等待时间，影响总体执行时间。因此可以采用下面这种替代方案：


%%C080472 code

\begin{codeblockbg}
\begin{Verbatim}[commandchars=\\\{\},breaklines=true,breakanywhere=true,fontsize=\codefont,xleftmargin=17bp]
\PYG{k}{template}\PYG{+w}{ }\PYG{o}{\PYGZlt{}}\PYG{k}{typename}\PYG{+w}{ }\PYG{n+nc}{Iter}\PYG{p}{,}\PYG{+w}{ }\PYG{k}{typename}\PYG{+w}{ }\PYG{n+nc}{F}\PYG{o}{\PYGZgt{}}
\PYG{k+kt}{void}\PYG{+w}{ }\PYG{n}{parallel\PYGZus{}map}\PYG{p}{(}\PYG{n}{Iter}\PYG{+w}{ }\PYG{n}{begin}\PYG{p}{,}\PYG{+w}{ }\PYG{n}{Iter}\PYG{+w}{ }\PYG{n}{end}\PYG{p}{,}\PYG{+w}{ }\PYG{n}{F}\PYG{+w}{ }\PYG{n}{f}\PYG{p}{)}
\PYG{p}{\PYGZob{}}
\PYG{+w}{  }\PYG{k}{auto}\PYG{+w}{ }\PYG{n}{size}\PYG{+w}{ }\PYG{o}{=}\PYG{+w}{ }\PYG{n}{std}\PYG{o}{:}\PYG{o}{:}\PYG{n}{distance}\PYG{p}{(}\PYG{n}{begin}\PYG{p}{,}\PYG{+w}{ }\PYG{n}{end}\PYG{p}{)}\PYG{p}{;}
\PYG{k}{if}\PYG{p}{(}\PYG{n}{size}\PYG{+w}{ }\PYG{o}{\PYGZlt{}}\PYG{o}{=}\PYG{+w}{ }\PYG{l+m+mi}{10000}\PYG{p}{)}
\PYG{p}{\PYGZob{}}
\PYG{+w}{  }\PYG{n}{std}\PYG{o}{:}\PYG{o}{:}\PYG{n}{transform}\PYG{p}{(}\PYG{n}{begin}\PYG{p}{,}\PYG{+w}{ }\PYG{n}{end}\PYG{p}{,}\PYG{+w}{ }\PYG{n}{begin}\PYG{p}{,}\PYG{+w}{ }\PYG{n}{std}\PYG{o}{:}\PYG{o}{:}\PYG{n}{forward}\PYG{o}{\PYGZlt{}}\PYG{n}{F}\PYG{o}{\PYGZgt{}}\PYG{p}{(}\PYG{n}{f}\PYG{p}{)}\PYG{p}{)}\PYG{p}{;}
\PYG{p}{\PYGZcb{}}
\PYG{k}{else}
\PYG{p}{\PYGZob{}}
\PYG{+w}{  }\PYG{k}{auto}\PYG{+w}{ }\PYG{n}{middle}\PYG{+w}{ }\PYG{o}{=}\PYG{+w}{ }\PYG{n}{begin}\PYG{p}{;}
\PYG{+w}{  }\PYG{n}{std}\PYG{o}{:}\PYG{o}{:}\PYG{n}{advance}\PYG{p}{(}\PYG{n}{middle}\PYG{p}{,}\PYG{+w}{ }\PYG{n}{size}\PYG{+w}{ }\PYG{o}{/}\PYG{+w}{ }\PYG{l+m+mi}{2}\PYG{p}{)}\PYG{p}{;}
\PYG{+w}{    }\PYG{k}{auto}\PYG{+w}{ }\PYG{n}{result}\PYG{+w}{ }\PYG{o}{=}\PYG{+w}{ }\PYG{n}{std}\PYG{o}{:}\PYG{o}{:}\PYG{n}{async}\PYG{p}{(}
\PYG{+w}{      }\PYG{n}{std}\PYG{o}{:}\PYG{o}{:}\PYG{n}{launch}\PYG{o}{:}\PYG{o}{:}\PYG{n}{deferred}\PYG{p}{,}
\PYG{+w}{      }\PYG{n}{parallel\PYGZus{}map}\PYG{o}{\PYGZlt{}}\PYG{n}{Iter}\PYG{p}{,}\PYG{+w}{ }\PYG{n}{F}\PYG{o}{\PYGZgt{}}\PYG{p}{,}
\PYG{+w}{      }\PYG{n}{begin}\PYG{p}{,}\PYG{+w}{ }\PYG{n}{middle}\PYG{p}{,}\PYG{+w}{ }\PYG{n}{std}\PYG{o}{:}\PYG{o}{:}\PYG{n}{forward}\PYG{o}{\PYGZlt{}}\PYG{n}{F}\PYG{o}{\PYGZgt{}}\PYG{p}{(}\PYG{n}{f}\PYG{p}{)}\PYG{p}{)}\PYG{p}{;}
\PYG{+w}{    }\PYG{n}{parallel\PYGZus{}map}\PYG{p}{(}\PYG{n}{middle}\PYG{p}{,}\PYG{+w}{ }\PYG{n}{end}\PYG{p}{,}\PYG{+w}{ }\PYG{n}{std}\PYG{o}{:}\PYG{o}{:}\PYG{n}{forward}\PYG{o}{\PYGZlt{}}\PYG{n}{F}\PYG{o}{\PYGZgt{}}\PYG{p}{(}\PYG{n}{f}\PYG{p}{)}\PYG{p}{)}\PYG{p}{;}
\PYG{+w}{    }\PYG{n}{result}\PYG{p}{.}\PYG{n}{wait}\PYG{p}{(}\PYG{p}{)}\PYG{p}{;}
\PYG{+w}{  }\PYG{p}{\PYGZcb{}}
\PYG{p}{\PYGZcb{}}
\PYG{k}{template}\PYG{+w}{ }\PYG{o}{\PYGZlt{}}\PYG{k}{typename}\PYG{+w}{ }\PYG{n+nc}{Iter}\PYG{p}{,}\PYG{+w}{ }\PYG{k}{typename}\PYG{+w}{ }\PYG{n+nc}{R}\PYG{p}{,}\PYG{+w}{ }\PYG{k}{typename}\PYG{+w}{ }\PYG{n+nc}{F}\PYG{o}{\PYGZgt{}}
\PYG{k}{auto}\PYG{+w}{ }\PYG{n}{parallel\PYGZus{}reduce}\PYG{p}{(}\PYG{n}{Iter}\PYG{+w}{ }\PYG{n}{begin}\PYG{p}{,}\PYG{+w}{ }\PYG{n}{Iter}\PYG{+w}{ }\PYG{n}{end}\PYG{p}{,}\PYG{+w}{ }\PYG{n}{R}\PYG{+w}{ }\PYG{n}{init}\PYG{p}{,}\PYG{+w}{ }\PYG{n}{F}\PYG{+w}{ }\PYG{n}{op}\PYG{p}{)}
\PYG{p}{\PYGZob{}}
\PYG{+w}{  }\PYG{k}{auto}\PYG{+w}{ }\PYG{n}{size}\PYG{+w}{ }\PYG{o}{=}\PYG{+w}{ }\PYG{n}{std}\PYG{o}{:}\PYG{o}{:}\PYG{n}{distance}\PYG{p}{(}\PYG{n}{begin}\PYG{p}{,}\PYG{+w}{ }\PYG{n}{end}\PYG{p}{)}\PYG{p}{;}
\PYG{k}{if}\PYG{p}{(}\PYG{n}{size}\PYG{+w}{ }\PYG{o}{\PYGZlt{}}\PYG{o}{=}\PYG{+w}{ }\PYG{l+m+mi}{10000}\PYG{p}{)}
\PYG{+w}{  }\PYG{k}{return}\PYG{+w}{ }\PYG{n}{std}\PYG{o}{:}\PYG{o}{:}\PYG{n}{accumulate}\PYG{p}{(}\PYG{n}{begin}\PYG{p}{,}\PYG{+w}{ }\PYG{n}{end}\PYG{p}{,}\PYG{+w}{ }\PYG{n}{init}\PYG{p}{,}\PYG{+w}{ }\PYG{n}{std}\PYG{o}{:}\PYG{o}{:}\PYG{n}{forward}\PYG{o}{\PYGZlt{}}\PYG{n}{F}\PYG{o}{\PYGZgt{}}\PYG{p}{(}\PYG{n}{op}\PYG{p}{)}\PYG{p}{)}\PYG{p}{;}
\PYG{k}{else}
\PYG{p}{\PYGZob{}}
\PYG{+w}{  }\PYG{k}{auto}\PYG{+w}{ }\PYG{n}{middle}\PYG{+w}{ }\PYG{o}{=}\PYG{+w}{ }\PYG{n}{begin}\PYG{p}{;}
\PYG{+w}{  }\PYG{n}{std}\PYG{o}{:}\PYG{o}{:}\PYG{n}{advance}\PYG{p}{(}\PYG{n}{middle}\PYG{p}{,}\PYG{+w}{ }\PYG{n}{size}\PYG{+w}{ }\PYG{o}{/}\PYG{+w}{ }\PYG{l+m+mi}{2}\PYG{p}{)}\PYG{p}{;}
\PYG{k}{auto}\PYG{+w}{ }\PYG{n}{result1}\PYG{+w}{ }\PYG{o}{=}\PYG{+w}{ }\PYG{n}{std}\PYG{o}{:}\PYG{o}{:}\PYG{n}{async}\PYG{p}{(}
\PYG{+w}{  }\PYG{n}{std}\PYG{o}{:}\PYG{o}{:}\PYG{n}{launch}\PYG{o}{:}\PYG{o}{:}\PYG{n}{async}\PYG{p}{,}
\PYG{+w}{      }\PYG{n}{parallel\PYGZus{}reduce}\PYG{o}{\PYGZlt{}}\PYG{n}{Iter}\PYG{p}{,}\PYG{+w}{ }\PYG{n}{R}\PYG{p}{,}\PYG{+w}{ }\PYG{n}{F}\PYG{o}{\PYGZgt{}}\PYG{p}{,}
\PYG{+w}{      }\PYG{n}{begin}\PYG{p}{,}\PYG{+w}{ }\PYG{n}{middle}\PYG{p}{,}\PYG{+w}{ }\PYG{n}{R}\PYG{p}{\PYGZob{}}\PYG{p}{\PYGZcb{}}\PYG{p}{,}\PYG{+w}{ }\PYG{n}{std}\PYG{o}{:}\PYG{o}{:}\PYG{n}{forward}\PYG{o}{\PYGZlt{}}\PYG{n}{F}\PYG{o}{\PYGZgt{}}\PYG{p}{(}\PYG{n}{op}\PYG{p}{)}\PYG{p}{)}\PYG{p}{;}
\PYG{+w}{    }\PYG{k}{auto}\PYG{+w}{ }\PYG{n}{result2}\PYG{+w}{ }\PYG{o}{=}\PYG{+w}{ }\PYG{n}{parallel\PYGZus{}reduce}\PYG{p}{(}\PYG{n}{middle}\PYG{p}{,}\PYG{+w}{ }\PYG{n}{end}\PYG{p}{,}\PYG{+w}{ }\PYG{n}{init}\PYG{p}{,}
\PYG{+w}{                                   }\PYG{n}{std}\PYG{o}{:}\PYG{o}{:}\PYG{n}{forward}\PYG{o}{\PYGZlt{}}\PYG{n}{F}\PYG{o}{\PYGZgt{}}\PYG{p}{(}\PYG{n}{op}\PYG{p}{)}\PYG{p}{)}\PYG{p}{;}
\PYG{+w}{    }\PYG{k}{return}\PYG{+w}{ }\PYG{n}{result1}\PYG{p}{.}\PYG{n}{get}\PYG{p}{(}\PYG{p}{)}\PYG{+w}{ }\PYG{o}{+}\PYG{+w}{ }\PYG{n}{result2}\PYG{p}{;}
\PYG{+w}{  }\PYG{p}{\PYGZcb{}}
\PYG{p}{\PYGZcb{}}
\end{Verbatim}
\end{codeblockbg}


%%C080473 p

当我们用同样的测试方法，把执行时间与第一个使用异步函数的实现相比较时，会发现这个版本（在下面的输出中用\texttt{p2}表示）无论\texttt{map}还是\texttt{fold}都与串行版本接近，而且比前面给出的第一个并行版本（用\texttt{p1}表示）差得多：


%%C080474 code

\begin{codeblockbg}
\begin{Verbatim}[commandchars=\\\{\},breaklines=true,breakanywhere=true,fontsize=\codefont,xleftmargin=17bp]
\PYG{+w}{    }\PYG{n}{size}\PYG{+w}{   }\PYG{n}{s}\PYG{+w}{ }\PYG{n}{map}\PYG{+w}{ }\PYG{n}{p1}\PYG{+w}{ }\PYG{n}{map}\PYG{+w}{  }\PYG{n}{p2}\PYG{+w}{ }\PYG{n}{map}\PYG{+w}{  }\PYG{n}{s}\PYG{+w}{ }\PYG{n}{fold}\PYG{+w}{ }\PYG{n}{p1}\PYG{+w}{ }\PYG{n}{fold}\PYG{+w}{ }\PYG{n}{p2}\PYG{+w}{ }\PYG{n}{fold}
\PYG{+w}{   }\PYG{l+m+mi}{10000}\PYG{+w}{      }\PYG{l+m+mi}{11}\PYG{+w}{     }\PYG{l+m+mi}{11}\PYG{+w}{      }\PYG{l+m+mi}{10}\PYG{+w}{       }\PYG{l+m+mi}{7}\PYG{+w}{      }\PYG{l+m+mi}{10}\PYG{+w}{      }\PYG{l+m+mi}{10}
\PYG{+w}{  }\PYG{l+m+mi}{100000}\PYG{+w}{     }\PYG{l+m+mi}{111}\PYG{+w}{    }\PYG{l+m+mi}{275}\PYG{+w}{     }\PYG{l+m+mi}{120}\PYG{+w}{      }\PYG{l+m+mi}{72}\PYG{+w}{      }\PYG{l+m+mi}{96}\PYG{+w}{     }\PYG{l+m+mi}{426}
\PYG{+w}{  }\PYG{l+m+mi}{500000}\PYG{+w}{     }\PYG{l+m+mi}{551}\PYG{+w}{    }\PYG{l+m+mi}{230}\PYG{+w}{     }\PYG{l+m+mi}{596}\PYG{+w}{     }\PYG{l+m+mi}{365}\PYG{+w}{     }\PYG{l+m+mi}{210}\PYG{+w}{    }\PYG{l+m+mi}{1802}
\PYG{+w}{ }\PYG{l+m+mi}{1000000}\PYG{+w}{    }\PYG{l+m+mi}{1142}\PYG{+w}{    }\PYG{l+m+mi}{381}\PYG{+w}{    }\PYG{l+m+mi}{1209}\PYG{+w}{     }\PYG{l+m+mi}{753}\PYG{+w}{     }\PYG{l+m+mi}{303}\PYG{+w}{    }\PYG{l+m+mi}{2378}
\PYG{+w}{ }\PYG{l+m+mi}{2000000}\PYG{+w}{    }\PYG{l+m+mi}{2411}\PYG{+w}{    }\PYG{l+m+mi}{981}\PYG{+w}{    }\PYG{l+m+mi}{2488}\PYG{+w}{    }\PYG{l+m+mi}{1679}\PYG{+w}{     }\PYG{l+m+mi}{503}\PYG{+w}{    }\PYG{l+m+mi}{4190}
\PYG{+w}{ }\PYG{l+m+mi}{5000000}\PYG{+w}{    }\PYG{l+m+mi}{5962}\PYG{+w}{   }\PYG{l+m+mi}{2191}\PYG{+w}{    }\PYG{l+m+mi}{6237}\PYG{+w}{    }\PYG{l+m+mi}{4177}\PYG{+w}{    }\PYG{l+m+mi}{1969}\PYG{+w}{    }\PYG{l+m+mi}{7974}
\PYG{l+m+mi}{10000000}\PYG{+w}{   }\PYG{l+m+mi}{11961}\PYG{+w}{   }\PYG{l+m+mi}{4517}\PYG{+w}{   }\PYG{l+m+mi}{12581}\PYG{+w}{    }\PYG{l+m+mi}{8384}\PYG{+w}{    }\PYG{l+m+mi}{2966}\PYG{+w}{   }\PYG{l+m+mi}{15174}
\end{Verbatim}
\end{codeblockbg}


%%C080475 h3

\subsection{延伸阅读}


%%C080476 li

\begin{itemize}
\item \textit{用线程实现并行map与fold}


%%C080477 li

\item \textit{异步执行函数}

\end{itemize}
